% \begin{savequote}
% \includegraphics[width=4cm]{./images/gabriel-christine-aass/2010-06-14_010.jpg}
% \end{savequote}
\def\ChapterCode{EHI}
\chapter
[\CO{[\ChapterCode]} Erste Hilfe]
{Erste Hilfe}
\label{chp:EHI}


\ChapterIntro
{%
Unter Erster Hilfe versteht man Maßnahmen, um das Leben und
die Gesundheit von Patienten zu retten und bis zum
Eintreffen professioneller Hilfe zu erhalten, sowie um
bedrohende Gefahren abzuwenden. Erste-Hilfe-Maßnahmen sind
in der Regel simpel und mit wenig Aufwand durchzuführen. Sie
sind ein essentieller Bestandteil der Rettungskette und
haben, wenn sie frühzeitig und richtig durchgeführt werden,
einen enormen Einfluß auf das spätere Behandlungsergebnis. 
}
{%
\begin{description}
  \item[Maintainer:] Roman Koch
  \item[Co-Maintainer:] Christof Koller
  \item[Reviewer:] Standard-Reviewprozess
  \item[Version:] {\DocVersion} ({\DocVersionInfo})
  \item[SHA1:] {\DocGitCommit}
\end{description}

{\TxTeilkapitelAASS}
}

\section{Einleitung}


\subsection{Allgemeine Einleitung}


\Dictum
[Herr von Lips.\\
In: J. N. Nestroy: Der Zerissene]
{Meiner Seel, 's is a fürchterlich's G'fühl,\\
Wenn man selber nicht weiß, was man will.}

Erste Hilfe ist die Rettung von Menschen aus
lebensbedrohlichen Situationen unter Einsatz von einfachen
Mitteln und Techniken. Sie kann von jedem schnell und
effektiv durchgeführt werden. Auch im professionellem
Bereich stellt sie den Ausgangspunkt und die Basis einer
effizienten Behandlung dar.

\TextSlide{Grundsätze}{
Bei jeder Behandlung hat sich der Ersthelfer an folgende
Grundsätze zu halten:
\label{topic:behandlungspflicht-ehi}%
\begin{enumerate}
  \item Jeder Mensch ist zur
  Durchführung von Erster Hilfe verpflichtet! Sie darf nur bei
  Unzumutbarkeit abgelehnt werden!\ZFootnote{Unzumutbar ist
  die Erste Hilfe für
  medizinisches Fachpersonal
  nur dann, wenn durch Setzen   der Maßnahme das eigene
  Leben bzw. die eigene Gesundheit gefährdet wäre.}
  \item Die zumutbare
  Erste-Hilfe-Maßnahme ist zumindest zu versuchen und darf
  nicht von vornherein als unmöglich abgelehnt
  werden!\ZFootnote{Gemeint sind hierbei Fälle in denen für
  den
  Ersthelfer zwar keine
  Gefahr droht, er sich aber an seine physischen Grenzen
  gesetzt fühlt. Ein denkbares Beispiel wäre die   Rettung
  eines stark übergewichtigen Menschen aus einem
  Gefahrenbereich durch   eine körperlich schwache Person.}
  \item Die Gesamtbehandlung des
  Patienten unterliegt einer Interessensabwägung: Wichtiges
  vor Unwichtigem!\ZFootnote{Es muss immer das nächstgrößere
  Übel beseitigt werden.
  Es kann mitunter auf Grund der Dringlichkeit in der
  Situation nicht immer Rücksicht auf das Eigentum des
  Patienten genommen werden.
  Als Beispiel dient das Aufschneiden der Kleidung des
  Patienten zur besseren Untersuchung und Behandlung.
  }
  \item Bei jeder Einzelbehandlung ist der Patient so
  schnell und so schonend wie möglich zu behandeln!
\end{enumerate}

}{
\begin{enumerate}
  \item Verpflichtung zur Hilfeleistung
  \item Hilfeleistung muss zumindest versucht werden
  \item Interessensabwägung
  \item Hilfeleistung so schnell und schonend wie möglich
\end{enumerate}
}
{}
{}
{}

\VARIANT{VAR-ERNAEHRUNGSPAED}{
\Layout{1}{Hilfeleistungspflicht}{
\subparagraph{Ersthelferische Versorgungspflicht für
Durchschnittsbürger}
Jedermann ist
zur Setzung von Hilfsleistungen verpflichtet, die seinem
Wissens- und Kenntnisstand entsprechen. Daher ist jeder mit
einem Führerschein-Erste-Hilfe-Kurs zur Setzung von
zumutbaren Erste-Hilfe Maßnahmen
verpflichtet.\ZFootnote{{\S}\,95 StGB}

\subparagraph{Unzumutbarkeit der Behandlung}
Für alle Berufe oder Durchschnittsbürger ist die Behandlung dann
nicht zumutbar, wenn man \ldots
\begin{itemize}
  \item das eigene Leben oder die Gesundheit gefährden
  würde, oder
  \item höherwertige Interessen verletzen
  müsste.\ZFootnote{{\S}\,95 Abs. 2 StGB}
\end{itemize}
 


\begin{description}
  \item[{\S}\,95 StGB Unterlassung der Hilfeleistung] ~
  \begin{enumerate}
    \item[(1)] Wer es bei einem Unglücksfall oder einer
    Gemeingefahr ({\S}\,176) unterläßt, die zur Rettung
    eines Menschen aus der Gefahr des Todes oder einer
    beträchtlichen Körperverletzung oder
    Gesundheitsschädigung offensichtlich erforderliche Hilfe
    zu leisten, ist mit Freiheitsstrafe bis zu sechs Monaten
    oder mit Geldstrafe bis zu 360 Tagessätzen, wenn die
    Unterlassung der Hilfeleistung jedoch den Tod eines
    Menschen zur Folge hat, mit Freiheitsstrafe bis zu einem
    Jahr oder mit Geldstrafe bis zu 360 Tagessätzen zu
    bestrafen, es sei denn, daß die Hilfeleistung dem Täter
    nicht zuzumuten ist.
    \item[(2)] Die Hilfeleistung ist insbesondere dann nicht
    zuzumuten, wenn sie nur unter Gefahr für Leib oder Leben
    oder unter Verletzung anderer ins Gewicht fallender
    Interessen möglich
    wäre
  \end{enumerate}
  \item[{\S}\,94 StGB Imstichlassen eines Verletzten] ~
  \begin{enumerate}
    \item[(1)] Wer es unterläßt, einem anderen,
    dessen Verletzung am Körper ({\S}\,83) er, wenn auch nicht
    widerrechtlich, verursacht hat, die erforderliche Hilfe zu
    leisten, ist mit Freiheitsstrafe bis zu einem Jahr oder
    mit Geldstrafe bis zu 360 Tagessätzen zu bestrafen.
    \item[(2)] Hat das Imstichlassen eine schwere
    Körperverletzung ({\S}\,84 Abs. 1) des Verletzten zur Folge,
    so ist der Täter mit Freiheitsstrafe bis zu zwei Jahren,
    hat es seinen Tod zur Folge, mit Freiheitsstrafe bis zu
    drei Jahren zu bestrafen.
    \item[(3)]  Der Täter ist entschuldigt, wenn ihm die
    Hilfeleistung nicht zuzumuten ist. Die Hilfeleistung ist
    insbesondere dann nicht zuzumuten, wenn sie nur unter der
    Gefahr des Todes oder einer beträchtlichen
    Körperverletzung oder Gesundheitsschädigung oder unter
    Verletzung anderer überwiegender Interessen möglich wäre.
    \item[(4)] Der Täter ist nach Abs. 1 und 2 nicht zu
    bestrafen, wenn er schon wegen der Verletzung mit der
    gleichen oder einer strengeren Strafe bedroht ist.
  \end{enumerate}
\end{description}



}{
\begin{itemize}
  \item Ersthelferische Versorgungspflicht für
  Durchschnittsbürger. 
  \item Unzumutbarkeit der Behandlung. 
\end{itemize}
}{}{}{}    

}


\protect\VARIANT{VAR-ERNAEHRUNGSPAED}{
clearpage
\subsection{Unterlassung einer Reanimation}
\label{topic:reanimation-unterlassung}
}


\protect\VARIANT{VAR-ERNAEHRUNGSPAED}{
\MassnahmenNG{Spezielle Maßnahmen}{Unterlassung der Reanimation}
{\label{m:reanimation-unterlassung}}
{
% \fcolorbox{ubuntured}{White}{
% \begin{minipage}[t]{\linewidth - {4\fboxrule} - {4\fboxsep}}
\Ca{Unter den folgenden Voraussetzungen darf (oder muss) eine
Reanimation unterlassen werden:}
\index{Reanimation!Unterlassung der}%
\index{Wiederbelebung|see{Reanimation}}%
\begin{multicols}{2}
\begin{itemize}
  \item \EE{Verwesung}
  \item \EE{Fäulnis}
  \item \EE{Mumifizierung}
  \item \EE{Skelettierung}
  \item \EE{Verletzungen, die nicht mit dem Leben
  vereinbar sind}
  \item \EE{Tierfraß}\footnote{Unter
  \E{\enquote{Tierfraß}} versteht man den Verzehr eines Kadavers
  durch Kadaverfresser. Ein frischer Hundebiss oder
  Ähnliches fällt \underbar{nicht} darunter. Maden können
  sich im Gewebe von Lebenden einnisten und sind daher
  auch kein sicheres Todeszeichen.}
  \item (Verbindliche) \EE{Patientenverfügung}\footnote{\E{Patientenverfügung}: Bei Notfallversorgungen darf \EE{nicht nach einer Verfügung gesucht
  werden, wenn dadurch das Leben oder die Gesundheit des Patienten ernstlich
  gefährdet werden würde}.
  \EE{Aber:} Wenn \E{zweifelsfrei} eine verbindliche
  Patientenverfügung vorliegt (z.\,B. weil eine betreuende
  Pflegekraft diese in der Zwischenzeit gefunden und die
  Verbindlichkeit bestätigt hat), muss dieser Folge
  geleistet werden.
Liegt eine \E{beachtliche Patientenverfügung} vor, so ist im
\E{Einzelfall} zu entscheiden: Das Wohl des Patienten bleibt
oberstes Gebot und die Verfügung soll in die
Entscheidung einfließen.}\index{Patientenverfügung!Unterlassung der
  Reanimation} (Bis zur zweifelsfreien Klärung der Situation muss
  jedenfalls eine Reanimation begonnen bzw. fortgeführt
  werden!)
  \item (Abbruch der Reanimationsmaßnahmen aufgrund
  von \EE{Erschöpfung})\footnote{Erschöpfung darf im
  professionellen Umfeld keine Rolle spielen, für eine
  entsprechende Ablösung ist bei Bedarf zu sorgen.}
  \item Unvereinbarkeit mit
  \EE{Selbstschutz}\footnote{Es müssen jedoch alle Maßnahmen
  getroffen werden, die ohne Gefährdung des Selbstschutzes
  möglich sind. So wird zum Beispiel für Ersthelfer
  empfohlen, wenn
  eine Mund-zu-Mund-Beatmung nicht zumutbar ist, zumindest 
  eine Herzdruckmassage durchzuführen.}
\end{itemize}
\end{multicols}
% \end{minipage}
% }
}

}


\subsection{Rettungskette}
\index{Rettungskette}
\label{topic:rettungskette}


Jede Bewältigung eines Notfalls sollte immer nach dem
selben Schema ablaufen, der Rettungskette.

% TODO EHI: Bild Rettungskette einfügen


\TextSlide{Rettungskette}
{
%%%%%%%%%%%%%%%%%%%%%%%%%%%%%%%%%%%%%%%%%%%%%%%%%%%%%%%%%%%
% TODO Rettungskette: Grafik Neuerstellung
\TODO{GABS}{yyyy-mm-dd}{Rettungskette: Grafik Neuerstellung}{}
%%%%%%%%%%%%%%%%%%%%%%%%%%%%%%%%%%%%%%%%%%%%%%%%%%%%%%%%%%%
Die Rettungskette besteht aus fünf Gliedern. Am Beginn
stehen die \EE{Lebensrettenden Sofortmaßnahmen}, welche noch
in \E{Allgemeine Sofortmaßnahmen} und \E{am Patienten
gesetzten Lebensrettenden Sofortmaßnahmen} unterteilt werden
können. Zu den \E{Allgemeinen Sofortmaßnahmen} zählt das
Absichern der Unfallstelle (Eigensicherung!), das Retten von
Personen aus Gefahrenzonen, sowie das Absetzen des Notrufs
durch eine 2. Person. Zu den am Patienten zu setzenden
\E{Lebensrettenden Sofortmaßnahmen} zählt die stabile
Seitenlage bei Bewusstlosigkeit, die
Herz-Lungen-Wiederbelebung bei Atem-Kreislaufstillstand,
Blutstillung bei starken Blutungen, sowie Schockbekämpfung
beim Schock. Nach den Lebensrettenden Sofortmaßnahmen wird
der \EE{Notruf} abgesetzt, falls dies noch nicht durch eine
2. Person geschehen ist.

Anschließend wird \EE{weitere Erste Hilfe} geleistet, wie
zum Beispiel die Versorgung von einfacheren Verletzungen und
die psychologische Betreuung des Patienten. Der eintreffende
Rettungsdienst übernimmt den Patienten für die
\EE{Sanitätshilfe} und führt den \EE{Transport} durch. Im
Spital erfolgt schließlich die \EE{ärztliche Betreuung} des
Patienten.

}{
\begin{enumerate}
  \item Lebensrettende Sofortmaßnahmen
  \item Notruf
  \item Weitere Erste Hilfe
  \item Sanitätshilfe - Transport
  \item Ärztliche Betreuung
\end{enumerate}
}
{}{}{}




Der Ersthelfer steht am Anfang der Rettungskette. Von seinem
Handeln hängt zunächst alles ab, im Extremfall auch das
Leben des Patienten.


\Fazit{\enquote{Eine Kette ist immer nur so stark 
wie ihr schwächstes Glied.}}

Dieser Satz soll  die Wichtigkeit der Erste-Hilfe-Maßnahmen
verdeutlichen. Schließlich werden von den fünf Gliedern der
Rettungskette die ersten drei vom Ersthelfer durchgeführt.


\section{Lebensrettende Sofortmaßnahmen}

\subsection{Verhalten am Notfallort}
\label{topic:verhalten-notfallort}



\TextSlide{Verhalten am Notfallort
}{
Vor dem Setzen von Maßnahmen sollte der Ersthelfer tief
durchatmen und sich danach einen \EE{Überblick} über die
Situation verschaffen: Was ist passiert? Um welche Art von
Notfall handelt es sich? Welche Gefahren können drohen?


Es gilt \EE{Ruhe zu bewahren} und \EE{Ruhe zu
vermitteln}. Dabei stellt man sich ruhig und gezielt auf die
bevorstehenden Maßnahmen ein und reiht diese nach
\EE{Dringlichkeit}. Erst nach diesen anfänglichen
Überlegungen wird mit einer Behandlung begonnen! 
}{
\begin{itemize}
  \item Überblick verschaffen
  \item Ruhe bewahren, Ruhe vermitteln
  \item Maßnahmen nach Dringlichkeit sortieren
\end{itemize}

}{}{}{}

\subsection{Gefahrenbereich}
\index{Gefahrenbereich}
\label{topic:gefahrenzone}

\TextSlide{Gefahr
}{
Durch das Verschaffen des Überblicks kann der Ersthelfer
\EE{Gefahren erkennen}. Diese können für den Patienten, den
Ersthelfer und dazukommende Dritte bestehen. Der Ersthelfer
muss \EE{Absperrmaßnahmen durchführen} und damit den
Notfallort absichern.

Muss der Patient aus dem Gefahrenbereich gerettet werden,
ist unbedingt auf Selbstschutz zu achten. Anschließend
müssen ggf. \EE{Spezialkräfte} angefordert und ihnen
sämtliche verfügbaren Informationen über die Gefahr bekannt
gegeben werden.
}{
\begin{description}
  \item [G]efahr erkennen
  \item [A]bsperrungsmaßnahmen durchführen - Absichern
  \item [S]pezialkräfte anfordern
\end{description}
}{}{}{}

\subsubsection{Absichern des Gefahrenbereichs}

Die Methoden der Absicherung sind vielfältig und müssen den
jeweiligen Gegebenheiten angepasst werden. Als Grundsatz
jeder Absicherung gilt:

\Fazit{Selbstschutz geht vor Fremdschutz!}

In der Praxis ist es oft möglich, durch
eine einfache Absicherungsmaßnahme das eigene Leben, sowie
das des Patienten, effektiv zu schützen.

\TextSlide{Beispiel: Vorgehen bei einem Verkehrsunfall}
{
Wenn sich der Ersthelfer mit einem Fahrzeug einer
Unfallstelle nähert, hält er nach Möglichkeit sein Fahrzeug
in einem angemessenen \EE{Sicherheitsabstand} zum
Gefahrenbereich an. Besondere \EE{Gefahrenquellen} wie
Kurven, Gefälle, Hangkuppen, sowie die Windrichtung bei
Rauchentwicklung müssen beim Absichern der Unfallstelle
berücksichtigt werden. Zum Absichern selbst werden
Hilfsmittel, wie \EE{Warnleuchte} oder \EE{Warndreieck},
sowie Personen, welche Handzeichen geben, eingesetzt.

\bigskip

\bigskip

}{
\begin{itemize}
  \item Sicherheitsabstand
  \item Gefahrenquellen beachten (Kurven, Gefälle,
  Windrichtung, Gefahrenstoffe,\ldots)
  \item Hilfsmittel
  \begin{itemize}
    \item Warnleuchte, Warndreieck, Handzeichen
  \end{itemize}
  \item Selbstschutz
  \begin{itemize}
    \item Warnweste
    \item Mit geöffnetem Warndreieck
    Verkehr entgegen gehen
  \end{itemize}
\end{itemize}
}{}{}{}



\Fazit{Niemals selbst gefährden!}

Daher hat der Ersthelfer vor dem Aussteigen aus dem Auto
eine \EE{Warnweste} anzuziehen und mit bereits
geöffnetem \EE{Warndreieck} (Pannendreieck) dem 
Verkehr entgegenzugehen. 


\TextSlide{Warndreieck}
{
Das Warndreieck ist
dem jeweiligen Anhalteweg entsprechend zur Unfallstelle
aufzustellen:
\begin{itemize}
  \item \E{Ortsgebiet}: unmittelbar bis zu ca. 50\Meter*
  \item \E{Freilandstraße}: ca. 150\Meter*
  \item \E{Autobahn}: ca. \VonBis{250}{400}\Meter*
\end{itemize}
}{
\begin{itemize}
  \item \E{Ortsgebiet}: unmittelbar bis zu ca. 50\Meter*
  \item \E{Freilandstraße}: ca. 150\Meter*
  \item \E{Autobahn}: ca. \VonBis{250}{400}\Meter*
\end{itemize}
}{}{}{}

\TextSlide{Eisenbahnkreuzung}
{\index{Eisenbahnkreuzung}
Bei einem Unfall an einer Eisenbahnkreuzung ist zu bedenken, dass der
herannahende Zug einen \E{viel längeren Bremsweg} als ein
Auto hat (je nach Geschwindigkeit und
Gewicht des Zuges, ca. \VonBis{600}{2\,600}\Meter*!). Das
\EE{Haltesignal} für den Lokführer ist eine \EE{kreisende Armbewegung}. 
Dadurch wird er gewarnt und leitet 
eine Notbremsung ein.
}{
\begin{itemize}
  \item Viel längerer Bremsweg
  \item Haltesignal: kreisende Armbewegung
\end{itemize}
}{}{}{}



\subsubsection{Gefahrengutransport}

Eine Besonderheit stellen \EE{Gefahrenguttransporte} dar
(\See{topic:gefahrengutunfall}). Sie sind -- wie der Name
schon sagt -- mit gefährlichen Gütern versehen und haben
somit ein erhöhtes Gefährdungspotential. Zu erkennen sind
Gefahrenguttransporter durch eine orange, rechteckige
Warntafel mit schwarzem Rand
(\prettyref{tab:gefahrentafel-ehi}).

\begin{WidePar}
\begin{table}[H]
\Captionof{table}{In der Ausführung gibt es neben der
unbeschrifteten Tafel eine mit zwei Ziffernkombinationen
ausgestattete Form der Warntafel. Die Beschriftung zeigt in
seinem oberen Feld die Gefahrennummer (Kemler-Nummer) und in
seinem unteren Bereich die Stoffnummer (UN-Nummer) an.
\E{Die Nummern sind beim Notruf
anzugeben}.\label{tab:gefahrentafel-ehi}}

\begin{tabular}{m{3.225cm}m{3.225cm}m{3.225cm}m{3.225cm}}
\toprule\addlinespace
\multicolumn{2}{m{6.65cm}}{\centering \TabHeading{Allgemeine Kennzeichnung}
(Sammeltransporte)} &
\multicolumn{2}{m{6.65cm}}{\centering \TabHeading{Spezielle Kennzeichnung}} 
\\\addlinespace\midrule\addlinespace
\includegraphics[width=\linewidth]{./images/khd/warntafel-allgemein.png}
 &
 &
\includegraphics[width=\linewidth]{./images/khd/warntafel-un-60-1710.png}
 &
Gefahrnummer\par (\T{Kemler-Nummer})

Stoffnummer\par (\T{UN-Nummer})
\\\addlinespace\bottomrule
\end{tabular}

\end{table}
\end{WidePar}


\Layout{57}{}{}{}{

}{Kennzeichnung durch Gefahrentafel.}{}



\subsubsection{Retten aus dem Gefahrenbereich}
\label{topic:wegziehen}
\label{topic:rautekgriff}

Durch Anwenden der oben genannten Absicherungsmaßnahmen kann
in vielen Situationen kein vollständiger Schutzbereich
errichtet werden.
Insbesondere im Straßenverkehr drohen, solange sich der
Patient auf der Fahrbahn befindet, weiterhin Gefahren, da
Absicherungsmittel leicht vom nachkommenden Verkehr
übersehen werden können. Daher ist der Patient von der
Fahrbahn in einen vollständig sicheren Bereich neben der
Fahrbahn zu verbringen.

\Fazit{Der Patient ist so rasch und schonend wie nur möglich aus der
Gefahrenzone zu bringen.}

Eine bekannte Technik ist der \T{Rautek-Handgriff}. Er darf
jedoch nur angewandt werden, wenn das \E{Leben des Patienten
akut und unmittelbar in Gefahr} ist, da er den Brustkorb und
die Wirbelsäule stark belastet
\cite{AsboeLehrmeinungRd2011-06}.



\Experimental{
\paragraph{Sitzender Patient:}
%%%%%%%%%%%%%%%%%%%%%%%%%%%%%%%%%%%%%%%%%%%%%%%%%%%%%%%%%%%
% TODO : Bild: Retten aus dem Gefahrenbereich
\TODO{GABS}{yyyy-mm-dd}{Bilder: Retten aus dem
Gefahrenbereich}{}
%%%%%%%%%%%%%%%%%%%%%%%%%%%%%%%%%%%%%%%%%%%%%%%%%%%%%%%%%%%

% \label{topic:rautekgriff}

\TakeNotesLarge
}

% TODO EHI: Bilder Rautek aus Auto einfügen


\Experimental{
\paragraph{Liegender Patient}
%%%%%%%%%%%%%%%%%%%%%%%%%%%%%%%%%%%%%%%%%%%%%%%%%%%%%%%%%%%
% TODO : REtten aus dem GEfahrenbereich, Liegender Patient: Neuerstellung
\TODO{GABS}{yyyy-mm-dd}{Neuerstellung von Text notwendig}{}
%%%%%%%%%%%%%%%%%%%%%%%%%%%%%%%%%%%%%%%%%%%%%%%%%%%%%%%%%%%

\TakeNotesLarge
}

% TODO EHI: Bilder Wegziehen einfügen


\subsection{Helmabnahme}
\label{topic:sturzhelmabnahme}

\TODO{GABS}{2010-08-01}{Mit Aussendung der ASBÖ-Akademie
(\protect\bibentry{ASBOEAkademieUpdate2010Sturzhelm},
\cite{ASBOEAkademieUpdate2010Sturzhelm}) bezüglich Konformität kontrollieren!}{}

\TextSlide{\TxBeschreibung}
{Ein Patient mit aufgesetztem Sturzhelm stellt eine
besondere Herausforderung dar. Es gilt:
% KORR ZWING 2010-03-xx, ok KOCH
\Fazit{Ein Sturzhelm ist grundsätzlich immer abzunehmen!}

Die Annäherung an den Patienten erfolgt in dessen
Blickrichtung, um ihn nicht zu überraschen. Bevor der
Patient angesprochen wird, muss der Helm fixiert
werden. Der Patient wird daraufhin aufgefordert sich nicht
zu bewegen. 

Idealerweise wird der Sturzhelm von zwei Helfern abgenommen
(\I[Zweihelfermethode!Helmabnahme]{Zweihelfermethode}), 
steht jedoch kein zweiter Helfer zur Verfügung, so
ist die Abnahme 
auch mit nur einem Helfer 
(\I[Einhelfermethode!Helmabnahme]{Einhelfermethode}) 
möglich. Der Kopf
muss in Neutralposition stabilisiert werden. Auf
die \EE{Halswirbelsäule} darf \EE{kein Zug} ausgeübt
werden und sie darf \EE{nicht gestaucht}
werden!~\cite{ITLSde6,AsboeLehrmeinungRd2011-07}

Im professionellen Umfeld muss bei Verdacht auf eine
Verletzung der (Hals-)Wirbelsäule
\VARIANT{VAR-STANDARD}{(\prettyref{MK-m:einschaetzung-indikation-ws-immobilisation})}
anschließend eine HWS-Schiene angelegt 
werden\VARIANT{VAR-STANDARD}{ (\prettyref{topic:hws-schiene-anlegen})}. 
}{
\begin{itemize}
  \item Zweihelfermethode (1. Wahl)
  \item Einhelfermethode
  \item Ziel:
  \begin{enumerate}
    \item Kopf in Neutralposition
    stabilisieren
    \item HWS nicht stauchen
    \item Keinen Zug ausüben
  \end{enumerate}
\end{itemize}
}
{}
{}
{}


\begin{PictureSeriesWide}{}{Bilderserie: Helmabnahme zu zweit
(Zweihelfermethode)\index{Bilderserie!Helmabnahme
(Zweihelfermethode)}\index{Helmabnahme!Bilderserie (Zweihelfermethode)}}
\PictureSeriesRowWide{3}
{./images/koch-aass/00800-gamma/HelmZweihelfer1b.JPG}
{Bewegungsverbot: \Speech{\enquote{Bitte bleiben Sie ruhig liegen
und bewegen Sie sich nicht!}} 
Kopf und Helm fixieren, Visier öffnen
\Speech{\enquote{Mein Name ist {\ldots} Wir
werden jetzt Ihren Helm abnehmen.}} }
{./images/koch-aass/00800-gamma/HelmZweihelfer2b.JPG}
{Ansprechen, Kinnriemen öffnen}
{./images/koch-aass/00800-gamma/HelmZweihelfer3b.JPG}
{Kopf stabilisieren, Helm abziehen}
{}
{}
\PictureSeriesRowWide{3}
{./images/koch-aass/00800-gamma/HelmZweihelfer4b.JPG}
{Helm abziehen \ACh{GABS}{2010-02-19}{Helm kippen bis
Nasenspitze sichtbar?}}
{./images/koch-aass/00800-gamma/HelmZweihelfer5b.JPG}
{Stabilisierung übernehmen}
{./images/koch-aass/00800-gamma/HelmZweihelfer6b.JPG}
{Kopf halten: Ohne Zug, nicht stauchen! 
\Z{(Ggfs. HWS-Schiene anlegen\VARIANT{VAR-STANDARD}{,
\prettyref{topic:hws-schiene-anlegen}:
\Speech{\enquote{Bitte bleiben Sie noch ruhig liegen, wir
werden Ihnen jetzt eine Halskrause anlegen, um Ihren Hals zu
stabilisieren.}}})}}
{}
{}
\end{PictureSeriesWide}


\begin{PictureSeriesWide}{}{Bilderserie: Helmabnahme alleine
(Einhelfermethode)\index{Bilderserie!Helmabnahme
(Einhelfermethode)}\index{Helmabnahme!Bilderserie (Einhelfermethode)}}
\PictureSeriesRowWide{3}
{./images/koch-aass/00800-gamma/HelmEinhelfer1b.JPG}
{Bewegungsverbot: \Speech{\enquote{Bitte bleiben Sie ruhig liegen und bewegen Sie
sich nicht!}}
Kopf und Helm fixieren, Visier öffnen.}
{./images/koch-aass/00800-gamma/HelmEinhelfer2b.JPG}
{Ansprechen, Kinnriemen öffnen. \Speech{\enquote{Mein Name ist {\ldots} Ich
werde jetzt Ihren Helm abnehmen.}}}
{./images/koch-aass/00800-gamma/HelmEinhelfer3b.JPG}
{Helm kippen \ldots}
{}
{}
\PictureSeriesRowWide{3}
{./images/koch-aass/00800-gamma/HelmEinhelfer4b.JPG}
{{\ldots} bis
Nasenspitze sichtbar ist.}
{./images/koch-aass/00800-gamma/HelmEinhelfer5b.JPG}
{Kopf stabilisieren, Helm kippen}
{./images/koch-aass/00800-gamma/HelmEinhelfer6b.JPG}
{Kopf ablegen.
\Z{(Ggfs. HWS-Schiene anlegen:
\Speech{\enquote{Bitte bleiben Sie noch ruhig liegen, wir werden
Ihnen jetzt eine Halskrause anlegen, um Ihren Hals zu
stabilisieren.})}}}
{}
{}
\end{PictureSeriesWide}






\subsection{Kontrolle der Vitalfunktionen}
\label{topic:lebensfunktionen-kontrolle}

\TextSlide{Kontrolle der Lebensfunktion
}{
Nachdem der Patient aus dem Gefahrenbereich gerettet wurde,
ist eine \E{Kontrolle seiner Lebensfunktionen} notwendig, um
eine genauere Einschätzung des Gesundheitszustandes zu
erlangen.
Bei der Kontrolle beschränkt man sich auf zwei einfach
festzustellende Vitalwerte -- das \EE{Bewusstsein} und die
\EE{Atmung}.

Eine Kontrolle des Kreislaufs, durch Tasten des Pulses an
der Halsschlagader (oder irgendeiner anderen Stelle), wird
normalerweise von einem Ersthelfer \E{nicht} durchgeführt.
Die Pulskontrolle ist für einen Ersthelfer auch gar nicht
notwendig, denn Kreislauf und Atmung sind derart eng
miteinander verknüpft, sodass ohne Kreislauf keine Atmung 
möglich ist.
}{
\begin{itemize}
  \item Bewusstseinskontrolle
  \item Atemkontrolle
\end{itemize}
}{}{}{}


\TextSlide{Bewusstseinskontrolle
}{
Die Kontrolle baut sich aus folgenden drei Teilschritten
auf. Der erste ist das laute und deutliche \EE{Ansprechen}
des Patienten. Idealerweise mit einfachen Anweisungen wie:

\begin{quote}
   \emph{\enquote{Können Sie mich hören? Machen Sie bitte die Augen
   auf!}}
\end{quote}

\noindent Gleichzeitig kann man den Patienten an seinen
Schultern \EE{berühren} und leicht rütteln.

\Cave{Patienten, die in Folge eines Unfalls verunglückt
sind, dürfen nicht, oder nur sehr schwach, gerüttelt
werden!}
Sollte der Patient darauf nicht reagieren, setzt der 
\I{professionelle} Ersthelfer im Bereich
des Kopfes einen \EE{Schmerzreiz}\footnote{Schmerzreiz: 
Der Laienhelfer setzt keinen Schmerzreiz.}. 

Jede
Reaktion wird als Lebenszeichen gewertet. Erfolgt keine
Reaktion muss der Ersthelfer \EE{um Hilfe rufen}, oder
besser -- wenn möglich -- eine zweite Person mit dem Notruf
beauftragen.

}{
\begin{itemize}
  \item Ansprechen
  \item Berühren
  \item Profi: Schmerzreiz setzen
  \item Ruf um Hilfe
\end{itemize}
}{}{}{}


\TextSlide{Atemkontrolle}{%
Damit die Atmung kontrolliert werden kann, muss der Ersthelfer die
\EE{Atemwege freimachen}. Im Verdachtsfall wird die
Mundhöhle auf Fremdkörper kontrolliert. Anschließend wird der Kopf
nackenwärts \EE{überstreckt}, damit die zurückfallende
Zunge die Atemwege nicht blockiert. Nun muss die
Atemkontrolle durch \EE{Hören}, \EE{Sehen} und
\EE{Fühlen} durchgeführt werden. Dafür wird das Ohr über
den Mund des Patienten gelegt und 10~Sekunden lang
überprüft, ob eine normale Atmung vorhanden ist.



\Fazit{Es
gilt: Wenn eine Atmung vorhanden ist, dann ist auch ein
Kreislauf vorhanden.}
}{
\begin{itemize}
  \item Atemwege freimachen
  \item Atemkontrolle: 10 Sek
  \begin{itemize}
    \item Hören
    \item Sehen
    \item Fühlen
  \end{itemize}
  \item Wenn Atmung vorhanden, ist auch ein Kreislauf vorhanden
\end{itemize}
}{}{}{}


\begin{PictureSeries}{}{Bilderserie: \EE{Atemkontrolle}
\SourcePicture{Hirtler}}
\PictureSeriesRow{2}
{./images/hirtler-aass/Reanimation-Atemkontrolle}
{Atemkontrolle: Kopf überstreckt, Helfer hört, fühlt, und
schaut auf den Brustkorb}
{empty}
{Raum für Notizen}
{empty}
{}
{}
{}
\end{PictureSeries}


\subsection{Notfalldiagnose Bewusstlosigkeit}
\index{Bewusstlosigkeit!Erste-Hilfe}
\label{topic:notfalldiagnose-bewusstlosigkeit}

\TextSlide{\TxBeschreibung
}{%
Nach der Kontrolle der Vitalfunktionen kann der Ersthelfer
zu der Erkenntnis kommen, dass der Patient zwar \EE{kein
Bewusstsein} hat, aber eine \EE{ausreichende Atmung} und
daher einen \EE{vorhandenen Kreislauf}. In einer solchen
Situation liegt die Notfalldiagnose \T{Bewusstlosigkeit}
vor. Der Patient besitzt keine Schutzreflexe mehr und ist
daher schutz- und hilflos. 
}{
\begin{itemize}
    \item Kein Bewusstsein
    \item Ausreichende Atmung
    \item Vorhandener Kreislauf
\end{itemize}
}{}{}{}

\protect\VARIANT{VAR-STANDARD}{
\MassnahmenMK{Erste-Hilfe-Maßnahmen}{Bewusstlosigkeit}{m:ehi-bewusstlosigkeit}{
\EE{Taktik}: \Speech{Stabile Seitenlage}
}{}
}
\protect\VARIANT{VAR-ERNAEHRUNGSPAED}{\SlideCompensationNoParskip{3}

\MassnahmenNG{Erste-Hilfe-Maßnahmen}{Bewusstlosigkeit}{\label{m:ehi-bewusstlosigkeit}}{
\begin{itemize}
    \item Stabile Seitenlage
\end{itemize}
}
}





\begin{PictureSeriesWide}{}{Bilderserie: Stabile
Seitenlage\index{Bilderserie!Stabile Seitenlage}\index{Stabile
Seitenlage!Bilderserie}}
\PictureSeriesRowWide{3}
{./images/koch-aass/00800-gamma/Seitenlage1b.JPG}
{}
{./images/koch-aass/00800-gamma/Seitenlage2b.JPG}
{}
{./images/koch-aass/00800-gamma/Seitenlage3b.JPG}
{}
{}
{}
\PictureSeriesRowWide{3}
{./images/koch-aass/00800-gamma/Seitenlage4b.JPG}
{}
{./images/koch-aass/00800-gamma/Seitenlage5b.JPG}
{}
{empty}
{}
{}
{}
\end{PictureSeriesWide}



\Text{Stabile Seitenlage}{\index{Stabile Seitenlage}
Jeder Ersthelfer ist in der Lage mit den nachfolgend beschriebenen, einfachen
Handgriffen eine \T[Stabile Seitenlage]{stabile Seitenlage} herzustellen.
\begin{enumerate}
  \item  Den zum Ersthelfer näherliegenden Arm
  des Patienten im rechten Winkel auflegen.
  \item  Unter das gegenüberliegende
  Knie des Patienten greifen und anheben. Das
  Handgelenk des gegenüber liegenden Arms des Patienten
  auf das Kniegelenk des Patienten geben.
  \item  Hand- und Kniegelenk des Patienten
  zum Ersthelfer drehen.
  \item  Kopf des Patienten
  überstrecken
\end{enumerate}
Durch die stabile Seitenlage bewegt sich die Zunge auf Grund
der Schwerkraft vom Rachenraum weg und die Atemwege werden
dadurch frei. Bei Erbrechen fließt der Großteil des
Erbrochenen in der stabilen Seitenlage allein aus dem Mund-
und Rachenraum. Die restlichen Bestandteile sind jedoch vom
Ersthelfer mittels eines um zwei Finger gewickelten Tuches
auszuräumen.
}{}{}{}{}


\TextSlide{Weiterführende Maßnahmen
}{
Nachdem der Patient in die stabile Seitenlage gebracht
wurde, muss der Ersthelfer weiterführende Maßnahmen setzen,
um Verschlechterungen des Gesundheitszustandes zu vermeiden.
Sehr wichtig ist der \EE{Wärmeerhalt}. Der Ersthelfer kann 
sich hierfür unterschiedlicher Mittel bedienen, wie z.\,B.
die, in einer Autoapotheke gemäß \Abk{ÖNORM} befindliche
Rettungsdecke. Weiters muss eine \EE{laufende Kontrolle der
Atmung} durchgeführt werden. Der Ersthelfer hat sich alle 30
Sekunden durch Überprüfung der Atemtätigkeit davon zu
überzeugen, dass der Patientenzustand unverändert geblieben
ist. Durch Öffnen eines Fensters in geschlossenen Räumen,
wird dem Patienten die benötigte \EE{Frischluft} zugeführt.
Auch der \EE{psychischen Betreuung} des Patienten durch den
Ersthelfer wird eine große Bedeutung beigemessen.
}{
\begin{itemize}
  \item Wärmeerhalt: zudecken
  \item Laufende Kontrolle der Atmung alle 30\Sec*
  \item Frischluftzufuhr: Fenster öffnen
  \item Psychische Betreuung
\end{itemize}
}{}{}{}




\subsection{Notfalldiagnose Atem- und Kreislaufstillstand}
\label{topic:reanimation-eh}
\label{topic:notfalldiagnose-atem-kreislaufstillstand}

\TextSlide{\TxBeschreibung}
{Nach der Kontrolle der Vitalfunktionen, kann der
Ersthelfer zum Ergebnis kommen, dass der Patient \EE{kein Bewusstsein} und 
\EE{keine normale Atmung} hat, d.\,h. \EE{kein Kreislauf vorhanden} ist.
In einer solchen Situation liegt die Notfalldiagnose
\T{Atem-/Kreislaufstillstand} vor.\bigskip
}{
\begin{itemize}
  \item Kein Bewusstsein
  \item Keine normale Atmung
  \item → Kein Kreislauf
\end{itemize}
}{}{}{}

\TextSlide{Achtung}
{Entgegen der Vorgehensweise der Rettungskette wird bei
einem Atem- und Kreislaufstillstand grundsätzlich bereits
\E{nach der Kontrolle der Vitalfunktionen} der \E{Notruf
abgesetzt} und erst danach mit der \E{lebensrettenden
Sofortmaßnahme} begonnen. \bigskip\bigskip
}{
\begin{itemize}
  \item Nach der Kontrolle der Vitalfunktionen Notruf absetzen
  \item Anschließend lebensrettende Sofortmaßnahmen einleiten
\end{itemize}
}{}{}{}


\protect\VARIANT{VAR-STANDARD}{
\MassnahmenMK{Erste-Hilfe-Maßnahmen}{Atem- und
Kreislaufstillstand}{m:ehi-atem-kreislaufstillstand}{
}{}
}
\protect\VARIANT{VAR-ERNAEHRUNGSPAED}{\SlideCompensationNoParskip{3}

\MassnahmenNG{Erste-Hilfe-Maßnahmen}{Atem- und
Kreislaufstillstand}{\label{m:ehi-atem-kreislaufstillstand}}{
\begin{itemize}
  \item Herz-Lungen-Wiederbelebung (Reanimation)
  \begin{itemize}
    \item Herzdruckmassage
    \item Beatmung
  \end{itemize}
\end{itemize}
}
}



\subsubsection{Herzdruckmassage}
\index{Herzdruckmassage}



\TextSlide{Technik}
{Die Herzdruckmassage bewirkt mittels Kompression des Herzens
zwischen Wirbelsäule und Brustbein ein künstliches Fließgeschehen
(Zirkulation) des Blutes.

\begin{itemize}
  \item \EE{Position und Druckpunkt}: Der Patient muss auf
  einem harten Untergrund liegen. Ggfs. ist der Patient aus einem Bett
  o.\,ä. auf den Boden zu bringen. Der Helfer kniet
  dann seitlich neben dem Körper des Patienten und sucht
  den richtigen \E{Druckpunkt} für die Herzdruckmassage
  auf. Dieser liegt in der \EE{Mitte des Brustkorbes}.
  \item \EE{Arm- und Körperhaltung}: Die richtige
  Arm- und Körperhaltung ist wichtig, um mit wenig Kraftaufwand eine
  hervorragende Leistung zu erzielen. Die \E{Finger
  sollten verschränkt sein}, nur der \E{Handballen liegt am Brustbein auf}.
  Die \E{Arme bleiben gestreckt}, und die \E{Schultern liegen genau über
  dem Druckpunkt}. Nun wird mit dem Gewicht des Oberkörpers
  Druck auf das Brustbein ausgeübt.
  \item \EE{Druckerzeugung}: Die
  Druckerzeugung hat \E{kräftig, aber nicht ruckartig} zu
  erfolgen. Die
  \E{Eindrücktiefe} beträgt \E{\VonBis{5}{6}\CM*}. Beim
  durchschnittlichen Erwachsenen ist das ungefähr \nicefrac{1}{3} des
  Brustkorbes.
  \item \EE{Druckentlastung}: Der Brustkorb muss sich wieder
  ganz ausdehnen können. Es darf daher \E{kein Restdruck}
  bestehen. Die Druck- und Entlastungsphasen dauern beide
  gleich lange.
  \item \EE{Arbeitsgeschwindigkeit}: Die
  Arbeitsgeschwindigkeit der Herzdruckmassage beträgt \E{100
  Massagen pro Minute}.
\end{itemize}


}{
\begin{itemize}
  \item Druckpunkt: \E{Mitte} des Brustkorbes
  \item Arm- und Körperhaltung
  \begin{itemize}
    \item Finger verschränkt,
    \item Handballen am Brustbein,
    \item Arme bleiben gestreckt,
    \item Schultern genau über den Druckpunkt.
  \end{itemize}
  \item Druckerzeugung
  \begin{itemize}
    \item kräftig aber nicht ruckartig
    \item Eindrücktiefe: \VonBis{5}{6}\CM*
  \end{itemize}
  \item Druckentlastung
  \item 100\PerMin*
\end{itemize}
}{}{}{}

\begin{PictureSeriesWide}{}{Bilderserie: \EE{Herzdruckmassage} \SourcePicture{Hirtler}}
\PictureSeriesRowWide{3}
{./images/hirtler-aass/Reanimation-Druckpunkt}
{Druckpunkt: Mitte des Brustkorbs. Die Finger werden
ineinander verschränkt, \ldots}
{./images/hirtler-aass/Reanimation-Position-1}
{{\ldots} die Arme durchgestreckt \ldots}
{./images/hirtler-aass/Reanimation-Position-2}
{{\ldots} und Oberkörper nach vorne gebeugt:
$\rightarrow$~Druck direkt auf Brustbein}
{}
{}
\end{PictureSeriesWide}

\TextSlide{Mögliche Fehlerquellen}
{Das Schlimmste, das der Ersthelfer falsch machen kann, ist
 nichts zu tun. Wenn reanimiert wird, kann es zu dem einen
 oder anderen Fehler kommen. Der Patient muss auf einer
 festen Unterlage gelagert werden, ein Bett ist \EE{keine
 feste Unterlage}! Es darf nicht vergessen werden, den
 Patienten auf den Boden oder eine andere harte Unterlage zu
 legen.

Grundsätzlich kann es auch zu \EE{Verletzungen durch den
Ersthelfer} kommen. Dazu zählen der Bruch des Brustbeines,
Lungen- oder Herzquetschungen und Rippenbrüche durch einen
zu festen Druck, Abweichungen des Druckpunktes, nicht
senkrecht ausgeführten oder zu starken Druck.

Ein zu ruckartiger Druck führt zu einer zu kurzen
Kompression des Herzens, was wiederum eine verkürzte
Auswurfphase bedeutet. Dies alles führt zu einer
\EE{ungenügenden Zirkulation} des Blutkreislaufs.

\Fazit{Wird einmal ein falscher Druck ausgeübt, oder eine
Rippe gebrochen: Nicht aufhören, sondern den richtigen
Druckpunkt erneut aufsuchen und weiter machen!}
}{
\begin{itemize}
  \item keine feste Unterlage
  \item Verletzungen durch den Ersthelfer
  \item ungenügende Zirkulation
\end{itemize}
}{}{}{}

\subsubsection{Beatmung}
\index{Beatmung!Erste-Hilfe}

\TextSlide{Arten}
{
Die bekannteste Methode zur Beatmung durch den Ersthelfer
ist die \T{Mund--zu--Mund--Beatmung}. Dabei wird bei
überstrecktem Kopf mit der stirnseitigen Hand die Nase des
Patienten zugedrückt. Anschließend mit dem eigenen Mund den
des Patienten umschließen und dabei die Lippen luftdicht an
den Patienten drücken. Nun wird die eigene Ausatemluft in
den Patient eingeblasen. Vor der nächsten Atemspende muss
der Mund wieder abgesetzt werden, um selbst frische Luft
einatmen zu können.

Die zweite Möglichkeit ist die \T{Mund--zu--Nase--Beatmung}.
Hier wird bei überstrecktem Kopf mit der kinnseitigen Hand
der Mund des Patienten zugehalten. Anschließend mit dem
eigenem Mund die Nase des Patienten umschließen und die
Lippen luftdicht an den Patienten drücken. Nun wird die
eigene Ausatemluft in den Patienten eingeblasen. Vor der
nächsten Atemspende muss der Mund wieder abgesetzt werden,
um selbst frische Luft einatmen zu können.

Bei Säuglingen und Kleinkindern wäre noch die
\T{Mund--zu--Mund/Nasen-Beatmung} möglich. In diesem
besonderen Fall wird der Mund und die Nase des Patient mit
dem eigenem Mund umschlossen.
}{
\begin{itemize}
  \item Mund--zu--Mund-Beatmung
  \item Mund--zu--Nasen-Beatmung
  \item Mund--zu--Mund/Nasen-Beatmung
\end{itemize}
}{}{}{}

\begin{PictureSeries}{}{Bilderserie:
\EE{Atemspende}
\SourcePicture{Hirtler}}
\PictureSeriesRow{2}
{./images/hirtler-aass/Reanimation-Atemspende}
{Mund-zu-Mund-Beatmung: Der Kopf muss
überstreckt werden!}
{./images/hirtler-aass/Beutelbeatmung1}
{Bei professionellen Helfern kann die Beatmung auch mittels
eines Beatmungsbeutels erfolgen, diese Technik erfordert
jedoch viel Übung.}
{empty}
{}
{}
{}
\end{PictureSeries}

\TextSlide{Wichtige Parameter für die Beatmung}
{
Die \EE{Atemfrequenz} des Erwachsenen beträgt
\VonBis{12}{16}\PerMin*, d.\,h. wir atmen
\VonBis{12}{16}\PerMin* ein und aus. 
Das
\EE{\underline{maximale} (!) Atemzugsvolumen} in ml
entspricht dem 10-fachen des Körpergewichtes, d.\,h. ein 90\Kg*
schwerer Mann kann 900\Ml* Luft ein- und wieder ausatmen.
Ein normaler Atemzug beträgt jedoch nur etwa 500\Ml*.
Die ausreichende \ce{O2}-Konzentration in der Umgebungs-/Einatmungsluft ist die
Voraussetzung für eine funktionierende Atmung. Ist sie
nicht gegeben, ist auch die beste Atemarbeit nutzlos. Der
normale \ce{O2}-Gehalt der Umgebungsluft beträgt
21\PC*. Der \ce{O2}-Gehalt in der \EE{Ausatemluft} beträgt
noch \Range{16}{17}\PC*, d.\,h. der Mensch verbraucht  
nur rund \nicefrac{1}{5} des
eingeatmeten Sauerstoffs. \cite{Silbernagl3-Atmung}
}{
\begin{itemize}
  \item Atemfrequenz \VonBis{12}{16}\PerMin*
  \item Max. Atemzugsvolumen [ml]: 10-fache des Körpergewichtes
  \item Einatemluft 21\PC* \ce{O2}-Gehalt
  \item Ausatemluft 16\PC* \ce{O2}-Gehalt
\end{itemize}
}{}{}{}

\TextSlide{Bei der Beatmung ist zu beachten:}
{
Primär ist auf den \EE{Selbstschutz} zu achten, wenn möglich
ist mit einem Beatmungstuch oder Beatmungsmaske zu beatmen.

Ein \EE{zu hohes Atemzugsvolumen} führt zu einem höheren
Druck im Nasen-Rachenraum, sobald die Lunge vollständig mit
Luft befüllt ist. Die Luft findet nur den Ausweg über die
Speiseröhre in den Magen. Dies kann während einer
Herz-Lungen-Wiederbelebung meist unbemerkt geschehen. Bei
einem \EE{zu niedrigem Atemzugsvolumen} ist die
\ce{O2}-Versorgung nicht sichergestellt, da die Luft nicht
bis in die Lunge kommt. Bei Bartträgern kann es zu einer
\EE{ungenügenden Abdichtung} kommen. Eine \EE{zu hohe
Atemfrequenz} hindert den Patienten am Ausatmen. Wenn es
beim Freimachen der Atemwege zu einer \EE{mangelnden
Überstreckung} des Kopfes kommt, kann es zu einer
unzureichenden Belüftung der Lunge kommen. Bei jeder
Atemspende ist der \EE{Brustkorb} des Patienten zu
\EE{beobachten}. Dieser muss sich deutlich heben und senken.

}{
\begin{itemize}
  \item Selbstschutz (Beatmungstuch)
  \item zu hohes Atemzugsvolumen ${\rightarrow}$ Luft in den Magen
  \item zu niedriges Atemzugsvolumen 
  \item ungenügende Abdichtung
  \item mangelnde Überstreckung
  \item Brustkorb beobachten 
\end{itemize}
}{}{}{}

\subsubsection{Algorithmus Herz-Lungen-Wiederbelebung}

\TextSlide{Vorgehensweise}{
Nach der Kontrolle des Bewusstseins und der Feststellung,
dass der Patient nicht ansprechbar ist, wird um Hilfe
gerufen. Anschließend durch Überstrecken des Kopfes die
Atemwege öffnen und die Atmung kontrollieren. Wenn keine
normale Atmung vorhanden ist, Notruf absetzen. Danach mit
der Herz-Lungen Wiederbelebung beginnen. Diese setzt sich
aus 30 Herzdruckmassagen und 2 Beatmungen zusammen.
}{
\begin{itemize}
  \item Bilderserie: \prettyref{fig:bilderserie-hlw}
  \item Algorithmus
\enquote{Basic Life Support mit AED}:
\FullPrettyRef{tab:algorithmus-bls-aed}
\end{itemize}
}
{}{}{}



\begin{PictureSeriesWide}{}{Bilderserie: Herz-Lungen-Wiederbelebung
(Reanimation)\label{fig:bilderserie-hlw}\index{Bilderserie!Reanimation}\index{Reanimation!Bilderserie}}
\PictureSeriesRowWide{3}
{./images/koch-aass/00800-gamma/Rea1b.JPG}
{Ansprechen, berühren, Schmerzreiz}
{./images/koch-aass/00800-gamma/Rea2b.JPG}
{Hilferuf}
{./images/koch-aass/00800-gamma/Rea3b.JPG}
{Atemwege freimachen}
{}
{}
\PictureSeriesRowWide{3}
{./images/koch-aass/00800-gamma/Rea4b.JPG}
{Atemkontrolle}
{empty}
{Notruf: 144}
{./images/koch-aass/00800-gamma/Rea5b.JPG}
{30 Herzdruckmassagen}
{}
{}
\PictureSeriesRowWide{3}
{./images/koch-aass/00800-gamma/Rea6b.JPG}
{2 Beatmungen}
{./images/koch-aass/00800-gamma/Rea7b.JPG}
{Wenn vorhanden: Defibrillation
(FullPrettyRef{chp:BLS}.\ACh{Instruktor 1}{0.8.0}{Weg
lassen. }\ACh{GABS}{2010-02-19}{Wieso? Hinweis auf AED wäre doch sinnvoll.} }
{empty}
{}
{}
{}
\end{PictureSeriesWide}





\TextSlide{Ende der HLW}
{\index{Reanimation!Beendigung}
Beendet wird eine Herz-Lungenwiederbelebung durch den Ersthelfer nur
dann, wenn einer der folgenden drei Umstände eintritt. Ohne diese
zusätzlichen Umstände darf der Ersthelfer die Reanimation nicht
unterbrechen: 

\begin{enumerate}
  \item \EE{Eintreffen von \E{kompetenter Hilfe}}: Der
  Ersthelfer darf aber beim Eintreffen der Rettungskräfte
  nicht sofort mit der Wiederbelebung  aufhören, sondern
  muss diese weiterführen bis er ausdrücklich vom
  Einsatzpersonal abgelöst  wird.
  \item \EE{Erschöpfung}: Sollte der Ersthelfer all seine
  Kräfte aufgebraucht haben, wird gezwungener Maßen die
  Wiederbelebung ab-, bzw. unterbrochen.
  \item \EE{Der Patient zeigt \EE{Lebenszeichen}}: Entweder
  erwacht der Patient und ist ansprechbar \footnote{eine
  leider äußerst seltene Form der Reanimationsbeendigung},
  oder der Patient gibt ein Husten oder Räuspern von sich
  bzw. atmet beim Beatmungsvorgang dem Ersthelfer entgegen.
  \EE{Erbrechen zählt \E{nicht} als Lebenszeichen.}
\end{enumerate}

}{
\begin{itemize}
  \item Eintreffen von kompetenter Hilfe
  \item Erschöpfung
  \item Patient zeigt Lebenszeichen
\end{itemize}
}{}{}{}


 
% \clearpage % TEMP temp clearpage
\TextSlide{Anschließende Behandlung}
{%
Die anschließenden Behandlungen orientieren sich an den
verschiedenen Beendigungsmöglichkeiten der Reanimation,
wodurch auch sie in ihrer Ausgestaltung verschieden sind.
Zeigt der Patient \E{Lebenszeichen}, ist eine Atemkontrolle
durchzuführen und wenn eine normale Atmung vorhanden ist,
wird der Patient in die stabile Seitenlage umgelagert.
Anschließend weiterführende Maßnahmen gemäß Notfalldiagnose
Bewusstlosigkeit.
Bei Eintreffen von \E{kompetenter Hilfe} muss der Ersthelfer
dem Rettungspersonal Auskunft über relevante Umstände und
die gesetzten Maßnahmen geben.
Wenn der \E{Ersthelfer erschöpft} ist, sollte er (falls
überhaupt konditionell möglich) erneut versuchen andere 
Ersthelfer zu finden. Scheitert auch dieser Versuch, kann
der Ersthelfer keine weitere Maßnahme setzen als den
Patienten zuzudecken und auf das Eintreffen des 
Rettungspersonals zu warten. Der Ersthelfer hat alles in
seiner Macht stehende getan und es ist ihm kein Vorwurf zu
machen.
}{
\begin{itemize}
  \item Lebenszeichen kontrollieren
  \item Atemkontrolle alle 30 Sek
  \item Stabile Seitenlage.
  \item Wärmeerhalt (zudecken)
  \item Frischluftzufuhr (Fensteröffnen, \ldots)
  \item Psychische Betreuung
  \item Kompetente Hilfe kommt ${\rightarrow}$ Auskunft über gesetzte
  Maßnahmen geben
  \item Erschöpfung ${\rightarrow}$ versuchen weitere Ersthelfer zu
  finden
\end{itemize}

}{}{}{}




\subsection{Starke Blutung}
\label{topic:starke-blutung}


\protect\VARIANT{VAR-STANDARD}{
\MassnahmenMK{Erste-Hilfe-Maßnahmen}{Starke
Blutung}{m:ehi-starke-blutung}{
}{}
}
\protect\VARIANT{VAR-ERNAEHRUNGSPAED}{
\MassnahmenNG{Erste-Hilfe-Maßnahmen}{Starke
Blutung}{\label{m:ehi-starke-blutung}}{
\begin{itemize}
  \item Blutstillung
  \begin{itemize}
    \item Hochhalten der betroffenen Gliedmaße
    \item Zudrücken,
    \item Abdrücken einer blutzuführenden Arterie,
    \item Druckverband,
    \item Abbindung.
  \end{itemize}
\end{itemize}

}
}

\SlideCompensationNoParskip{2}

\Layout{22}{}{\index{Blutstillung}
Für die Blutstillung werden folgende Materialien
benötigt: Handschuhe, sterile Wundauflagen (z.\,B.
Zetuvit\texttrademark{} Saugkompresse) in verschiedenen
Größen, Dreiecktücher und Mullbinden. Oft sind auch
Verbandspäckchen vorhanden. Hier ist eine Wundauflage
bereits in die Mullbinde eingearbeitet.
}
{}
{./images/koch-aass/00800-gamma/MaterialDruckverbandb.JPG}
{Material für einen Druckverband: Handschuhe, Wundauflage,
Druckkörper, Dreiecktuch.}
{}

\Layout{60}{Blutstillung durch Zudrücken}
{Mit einem \EE{direkten Druck in die stark blutende Wunde}
mittels Finger, Handballen oder Faust -- je nach Größe der
Verletzung -- mit einer Wundauflage kann diese gestoppt
werden. Ein \EE{Hochhalten der verletzten Extremität}
verstärkt die Wirkung. Dies ist auch die \EE{einzige
Möglichkeit Halsschlagaderblutungen zu stoppen}. Nachteil
dieser Methode ist, dass der Ersthelfer ab nun nicht mehr
loslassen darf und somit gehandicapt ist.
}{
\begin{itemize}
  \item Direkter Druck in die Wunde
  \item Hochhalten der verletzten Extremität
  \item Einzige Möglichkeit Halsschlagaderblutungen zu stoppen
\end{itemize}
}
{./images/hauer-david-aass/blutstillung_druck-00800.jpg}
{Simpel, aber Wirkungsvoll: Draufdrücken!}
{Hauer}

\TextSlide{Blutstillung durch Abdrücken
}{
Der Ersthelfer drückt in diesem Fall die blutzuführende
Arterie herzwärts der stark blutenden Wunde ab. Die
Durchführung ist an verschiedenen Stellen möglich. Bei
Blutungen am Schädeldach kann jeweils an der betroffenen
Seite an der Schläfe abgedrückt werden. Bei Blutungen im
Bereich des Gesichts wird je nach Seite am Kinn abgedrückt.
Bei Blutungen an Schulter oder Oberarm liegt die
Abdruckstelle in der jeweiligen Schlüsselbeingrube.
Unterarmblutungen werden an der betroffenen
Oberarminnenseite abgedrückt. Blutungen im Bereich des
Oberschenkels sind in der Leistenbeuge und
Unterschenkelblutungen an der Oberschenkelinnenseite
abzudrücken.

Nachteil hierbei ist nicht nur dasselbe Handicap wie bei
direkten Druck in die Wunde, sondern auch die Tatsache, dass
 die gesamte betroffene Extremität ohne Blutversorgung
bleibt. Außerdem ist das Auffinden der Druckstellen für
ungeübte Ersthelfer nicht immer möglich.
%%%%%%%%%%%%%%%%%%%%%%%%%%%%%%%%%%%%%%%%%%%%%%%%%%%%%%%%%%%
% TODO : Bild Abdrückpunkte
\TODO{GABS}{yyyy-mm-dd}{Abdruckpunkte: Neuerstellung von Grafik/Text 
                          notwendig}{}
%%%%%%%%%%%%%%%%%%%%%%%%%%%%%%%%%%%%%%%%%%%%%%%%%%%%%%%%%%%
}{
\begin{itemize}
  \item Abdrücken der blutzuführenden Arterie
\end{itemize}
}{}{}{}

\TextSlide{Druckverband
}{\index{Druckverband}
Der Druckverband ist die Standardmethode zur Stillung
einer Blutung an den Extremitäten.
Beim Druckverband wird nur Druck auf das
verletzte Gefäß ausgeübt, somit ist die Durchblutung der
Extremitäten unterhalb der Wunde gewährleistet. 

\Cave{Am Hals darf kein Druckverband angelegt
ewerden!} 

Ein Druckverband besteht aus drei Teilen:

\begin{itemize}
  \item \EE{Sterile Wundauflage}
  \item \EE{Druckkörper}: Es eignen sich nur kompakte, weiche
  und saugfähige Objekte. Ein solches Objekt, das sich auch in
  der Verbandsbox befindet, ist die Mullbinde.
  \item \EE{Befestigungsmittel}: Z.\,B. Dreiecktuchkrawatte
  oder eine nicht-elastische Mullbinde.
\end{itemize}

Das \EE{Hochlagern der
verletzten Extremität} verstärkt die
blutungsstillende Wirkung. 
}{
\begin{itemize}
  \item 3 Teile:
  \begin{itemize}
    \item Sterile Wundauflage
    \item Druckkörper
    \item Befestigungsmittel
  \end{itemize}
  \item Hochhalten der verletzten Extremität
\end{itemize}
}{}{}{}



\begin{PictureSeriesWide}{}{Bilderserie: Druckverband}
\PictureSeriesRowWide{3}
{./images/koch-aass/00800-gamma/Abdrueckenb.JPG}
{Als Erstmaßnahme Exremität hochlagern und mit einer
sterilen Wundauflage Blutung stillen}
{./images/koch-aass/00800-gamma/Druckverband1b.JPG}
{Druckkörper auflegen und mittels Dreiecktuchkrawatte
fixieren.}
{./images/koch-aass/00800-gamma/Druckverband2b.JPG} {}
{Extremität nach Möglichkeit weiterhin hochlagern}
{}
\end{PictureSeriesWide}



\TextSlide{Abbindung
}{\index{Abbindung}
Die Abbindung ist weitestgehend aus der Ersten Hilfe
eliminiert worden. Grund dafür ist, dass sich das Verhältnis
zwischen Nutzen und Folgeschäden massiv zu den Folgeschäden
verschoben hat. Außerdem hat sich die rettungsdienstliche
Vollversorgung inklusive rascher Eintreffzeiten und die
Frühversorgung in unfallchirugischen Fachabteilungen
deutlich verbessert.

\Cave{Abbindung nur bei lebensbedrohlichen Blutungen, die
nicht anders stillbar sind!\footnote{Z.B. bei totaler
Amputation, teilweiser Abtrennung/Zerfetzung oder
Einklemmung.}

\Ca{Die Abbindung ist die letzte Wahl!}}
 
Das Abbinden \EE{funktioniert} aus anatomischen Gründen
\EE{nur am Oberarm und Oberschenkel}, denn dort befindet
sich nur ein Knochen, gegen den sämtliche Blutgefäße
gedrückt werden können. Dazu darf nur ein \EE{breites und
schonendes Material} verwendet werden (z.\,B.
Dreiecktuchkrawatte, Blutdruckmanschette, etc.). Nachteil
ist die nicht mehr gegebenen Blutversorgung der Extremität.
Für die weitere Versorgung des Patienten ist es wichtig, den
\EE{Zeitpunkt der Abbindung} zu notieren, am besten gleich
am Oberarm bzw.
Oberschenkel (oder, wenn wegen Blutung und Verschmutzung 
nicht möglich, auf der Stirn)
 des Patienten.\footnote{Manche Autoren
empfehlen, dass die maximale Abbindezeit 30 Minuten
betragen sollte. Nach heutigem Stand der Wissenschaft wird
empfohlen, einmal angebrachte Abbindungen bis zur
definitiven Versorgung \E{nicht} (auch nicht probeweise) zu
öffnen.
\cite{PHTLS6,BuchfelderErsteHilfe4}}
}{ 
\begin{itemize}
  \item Nur am Oberarm und Oberschenkel
  \item Nur breites und schonendes Material
  \item Zeitpunkt der Abbindung notieren
\end{itemize}
}{}{}{}

\begin{PictureSeries}{}{Bilderserie: Abbinden am Oberarm}
\PictureSeriesRow{2}
{./images/koch-aass/00800-gamma/AbbindungArm1b.JPG}
{Abbinden am Oberarm}
{./images/koch-aass/00800-gamma/AbbindungArm2b.JPG}
{Abbindezeit notieren}
{}
{}
{}
{}
\end{PictureSeries}

\begin{PictureSeries}{}{Bilderserie: Abbinden am Oberschenkel}
\PictureSeriesRow{2}
{./images/koch-aass/00800-gamma/AbbindungBein1b.JPG}
{}
{./images/koch-aass/00800-gamma/AbbindungBein2b.JPG}
{}
{}
{}
{}
{}
\PictureSeriesRow{2}
{./images/koch-aass/00800-gamma/AbbindungBein3b.JPG}
{}
{./images/koch-aass/00800-gamma/AbbindungBein4b.JPG}
{}
{}
{}
{}
{}
\end{PictureSeries}


\clearpage
\TextSlide{Versorgung abgetrennter Gliedmaßen}
{
Die abgetrennten Gliedmaßen sind \EE{immer in das Spital mitzunehmen}. Nach
Versorgung des Patienten ist das \EE{Amputat steril} zu \EE{verbinden}. Für
den Transport gilt es \EE{kühl und trocken} zu \EE{lagern}. Dazu eignen
sich idealerweise Kühlakkus für Kühltaschen falls vorhanden. Wichtig ist aber,
dass das Amputat  nicht direkt mit Eis oder dem Kühlakku in Berührung
kommt, da es sonst zum Gefrieren und Absterben von Gewebe kommt.
}{
\begin{itemize}
  \item immer ins Spital mitnehmen
  \item Amputat ist steril zu verbinden
  \item für den Transport kühl und trocken lagern
\end{itemize}
}{}{}{}




% \clearpage % TEMP temp clearpage
\subsection{Schock}
\index{Schock!Erste-Hilfe}
\label{topic:schock-ehi}

\Definition{{\quotedblbase}Der Schock ist eine lebensbedrohliche
Kreislaufstörung, bei der es zu einer Unterversorgung der
Organe mit Sauerstoff kommt.{\textquotedblleft}}

\VARIANT{VAR-STANDARD}{\noindent\emph{Vgl. \FullPrettyRef{topic:schock}.}}

Das Krankheitsbild, welches umgangssprachlich als Schock
bezeichnet wird, ist eine Schreckreaktion mit Muskelzittern,
weichen Knien, Schwächegefühl und Blässe. Es hat aber mit
dem echten Schock nichts zu tun, sondern ist nur eine
Reaktion des Nervensystems und normalerweise nicht
bedrohlich. Beruhigen und hinlegen lassen, führt bei den
Patienten zu einer schnellen Erholung.

\TextSlide{\TxBeschreibung{} und Verlauf}
{
Der Schock ist eine \EE{lebensbedrohliche Kreislaufstörung},
bei der es zu einer \EE{Unterversorgung der lebenswichtigen
Organe}, wie Hirn, Herz, Lunge, Leber und Nieren mit
Sauerstoff kommt.

\begin{itemize}
  \item \EE{Zentralisation}: Der Körper hilft sich selbst
  indem er das Blut zu den lebenswichtigen Organen bringt
  und äußere Körpergebiete wie Arme, Beine und Haut weniger
  bis gar nicht durchblutet.
  \item \EE{Dezentralisation}: Hält der Zustand länger an,
  bzw. ist nach wie vor zu wenig Blut vorhanden, beginnt der
  Körper wieder zu dezentralisieren und lässt das Blut in
  weniger wichtige Körperregionen fließen. Der dadurch
  entstehende Sauerstoffmangel bei den lebenswichtigen
  Organen führt zu einer \EE{Schädigung der Organe}, die in
  weiterer Folge zum \EE{Versagen der Organe} führt. Der
  Schock ist daher unbehandelt in seiner finalen Wirkung
  tödlich!
\end{itemize}

}{
\begin{itemize}
  \item Lebensbedrohliche Kreislaufstörung
  \item Unterversorgung lebenswichtiger Organe:
  
  Hirn, Herz,
  Lunge, Leber, Nieren
  \item Zentralisation
  \item Dezentralisation
  \item Schädigung der Organe, Versagen der Organe
  \item Unbehandelt führt der Schock zum Tod!
\end{itemize}
}{}{}{}



\TextSlide{Ursachen}
{
Es gibt drei mögliche Ursachen, welche zum Schock führen
können:

\begin{itemize}
  \item \EE{Blut- und Flüssigkeitsverlust}: Der Blutverlust
  kommt von Verletzungen mit starken Blutungen.
  Flüssigkeitsverlust kann bei großflächigen Verbrennungen
  und
  bei extremen Brechdurchfall vorkommen.
  \item Bei \EE{Herzversagen} ist zwar
  genügend Blut vorhanden, aber die Pumpe -- das Herz -- ist
  zu schwach um es durch den Körper zu pumpen.
  \item Bei \EE{Vergiftungen und allergischen Reaktionen}
  bleibt zwar auch die Blutmenge konstant, aber die Gefäße
  erweitern sich. Daher entsteht ein unzureichender Druck in
  den Gefäßen, und das Blut kann nicht weiter transportiert
  werden.
\end{itemize}
}{
\begin{itemize}
  \item Blut- und Flüssigkeitsverlust
  \item Herzversagen
  \item Vergiftungen und allergische Reaktionen
\end{itemize}
}{}{}{}

\TextSlide{\TxSymptome{} und Kennzeichen
}{
Das Problem bei den Schockkennzeichen ist, dass nicht immer
alle sofort erkannt werden. Die häufigsten bzw. typischsten
sind hier angeführt.
Der Patient kann eine \EE{blasse Haut} haben. Weiters
besteht die Möglichkeit, dass er \EE{zittert und friert},
aber paradoxerweise dazu einen \EE{kalten, klebrigen
Schweiß} auf der Stirn hat. Der Patient kann außerdem
\EE{verwirrt, unruhig und ängstlich} sein. Er besitzt eine
\EE{beschleunigte, flache Atmung} und einen
\EE{beschleunigten, kaum tastbaren Puls}.
\bigskip
\bigskip

\bigskip

\bigskip
}{
\begin{itemize}
  \item Blasse, kühle Haut
  \item Zittern und frieren
  \item Kalter, klebriger Schweiß
  \item Verwirrt, unruhig, ängstlich
  \item Beschleunigte, flache Atmung
  \item Beschleunigter, kaum tastbarer Puls
\end{itemize}
}{}{}{}

\TextSlide{Lagerung
}{
Da der Schock aus drei verschiedenen Ursachen entstehen
kann, muss auch in der Behandlung ursachengerecht und damit
differenziert vorgegangen werden.  Das wichtigste ist jedoch
die \EE{situationsgerechte Lagerung} des Patienten. Bei
Bewusstlosigkeit stabile Seitenlage, bei Blut- und
Flüssigkeitsverlust wird der Oberkörper flach und die Beine
hoch gelagert. Bei Herzversagen wird Oberkörper leicht
erhöht und bei Atemnot stark erhöht gelagert.

\bigskip

\bigskip

}{
\begin{itemize}
  \item Situationsgerecht, z.\,B.:
  \begin{itemize}
    \item Oberkörper hoch
    \item Beine hoch
    \item Stabile Seitenlage
    \item \ldots
  \end{itemize}
\end{itemize}
}{}{}{}


\protect\VARIANT{VAR-STANDARD}{
\MassnahmenMK{Erste-Hilfe-Maßnahmen}{Schockbekämpfung}{m:ehi-schockbekaempfung}{
}{}
}
\protect\VARIANT{VAR-ERNAEHRUNGSPAED}{
\MassnahmenNG{Erste-Hilfe-Maßnahmen}{Schockbekämpfung}{\label{m:ehi-schockbekaempfung}}{

\begin{itemize}
  \item Ursache beseitigen (wenn möglich), ggfs. Blutstillung
  \item Beengende Kleidungsstücke öffnen
  \item Situationsgerechte Lagerung
  \item Wärmeerhalt
  \item Beruhigender Zuspruch
  \item Beobachtung
\end{itemize}
}
}


\section{Notruf}
\index{Notruf}


\TextSlide{Absetzen eines Notrufes}{
Der Ersthelfer muss darauf vorbereitet sein beim Notruf
folgende Fragen beantworten zu können.

\begin{itemize}
  \item \EE{Wo} ist der Notfallort? Ist es nicht möglich eine
  Adresse anzugeben, muss der Notfallort so eindeutig wie
  möglich beschrieben werden.
  \item \EE{Was} ist passiert? Die Art des Notfalls bekannt geben. Es
  reicht die Wahrnehmungen weiterzugeben, es müssen keine
  Diagnosen gestellt werden.
  \item \EE{Wie viele} Leute sind betroffen? Es ist
  notwendig, die Anzahl der Patienten bekannt zu geben, um
  eine entsprechende Anzahl an Rettungsmitteln einzuteilen.
  \item \EE{Wer} ruft an? Es ist sinnvoll
  eine Rückrufnummer bekannt zu geben. Sollte der Notfallort z.\,B. nicht gefunden
  werden, besteht die Möglichkeit zurückzurufen und nachzufragen.
\end{itemize}



Wichtig für die Aufnahme vollständiger Daten ist es, dass der
Ersthelfer nicht als Erster auflegt. Er muss warten bis der
Rettungsdienstmitarbeiter das Gespräch beendet und damit zu
verstehen gegeben hat, alle notwendigen Informationen erhalten zu haben.
}{
\begin{itemize}
  \item Wo?
  \item Was?
  \item Wie viele?
  \item Wer?
  \item Niemals als Erster auflegen!
\end{itemize}
}{}{}{}

\TextSlide{Notrufnummern
}{
\index{Notruf!-nummern}
\index{112}
\index{122}
\index{133}
\index{144}
Die drei in Österreich geltenden Notrufnummern
lauten:
\EE{\Code{122} Feuerwehr}, \EE{\Code{133} Polizei} und \EE{\Code{144}
Rettung}. Vor einigen Jahren kam durch die EU eine weitere
Notrufnummer dazu:
\EE{\Code{112} Euro-Notruf}. Die Besonderheit dieser Nummer
liegt darin, dass sie nicht nur in Österreich, sondern in
allen Mitgliedsstaaten der EU anwendbar ist. Bei Verwendung mit 
einem Mobiltelefon ist dies auch ohne SIM-Karte oder PIN-Code-Eingabe 
über das stärkste verfügbare Handynetz möglich.
Über die Nummer 112 wird man zu einer Stelle
weitergeleitet, bei der das Personal mehrere Sprachen
spricht, die Heimatsprache und Englisch sind verpflichtend.
}{
\begin{itemize}
  \item \Code{122} -- Feuerwehr
  \item \Code{133} -- Polizei
  \item \Code{144} -- Rettung
  \item \Code{112} -- Euro-Notruf
\end{itemize}
}{}{}{}

\TextSlide{Besondere Dienste}
{\index{Aerztefunkdienst@Ärztefunkdienst}
\index{AeFD@ÄFD|see{Ärztefunkdienst}}
\index{Vergiftungsinformationszentrale}
Die
\EE{Vergiftungsinformationszentrale} (VIZ) dient nur dazu
Auskunft über die Risiken und Behandlungen von sicher
festgestellten Vergiftungen zu geben (\E{Expertenrat}). Die
Vergiftungsinformationszentrale dient grundsätzlich dem
medizinischen Fachpersonal um Expertenrat einholen zu
können. Der Ersthelfer muss im Vergiftungsnotfall in erster
Linie die Rettung alarmieren.

Der \EE{Ärztefunkdienst} (\Abk{ÄFD})
ist nur von Montag bis Freitag in der Nacht und Samstag, Sonn- und
Feiertag rund um die Uhr besetzt. Er
ist weder personal- noch materialtechnisch dafür
ausgerüstet, Notfälle bzw. Unfälle zu bewältigen. 
Sein Zweck liegt in der
Behandlung von allgemeinmedizinischen
\enquote*{Ordinationsfällen}. Außerdem kann es zu Wartezeiten von
bis zu drei Stunden kommen.\footnote{Ein Umstand,
welcher leider viele Patienten dazu verleitet, lieber bei
der Rettung anzurufen, welche durchschnittlich in 11 Minuten
vor Ort ist.}
}{
\begin{description}
  \item[Vergiftungsinformationszentrale] \Code{01\,/\,406\,43\,43} --
  Expertenrat
  \item[Ärztefunkdienst] \Code{141} -- Allgemeinmedizinische Versorgung in
  der Nacht, an Feiertagen und Wochenenden
\end{description}
}{}{}{}



\section{Weitere Erste Hilfe und wichtige Krankheitsbilder}

Die weitere Erste Hilfe umfasst Maßnahmen zur Versorgung von
Verletzungen, Erkrankungen und Vergiftungen. Manche Autoren erfassen
unter diesem Begriff Umstände, die
angeblich weniger lebensbedrohlich sein sollen, als die zuvor behandelten. Tatsächlich
können sich auch die folgenden Beeinträchtigungen der
Patientenintegrität zu lebensbedrohlichen Situationen entwickeln. 

\VARIANT{VAR-ERNAEHRUNGSPAED}{
\subsection{Standardmaßnahmen bei vital bedrohten Patienten}
\subsubsection{Vitale Bedrohung}

\TextSlide{Vitale Bedrohung} {%
Eine vitale Bedrohung kann sich aus drei wesentlichen
Säulen ergeben: 
\begin{itemize}
  \item \EE{Massive BeABC-Störung}: Massive Störungen von
  Vitalfunktionen 1. Ordnung (Bewußtsein, Atmung,
  Kreislauf)
  \item \T{Alarm-Symptome}: Zu den Alarmsymptomen gehören
  u.\,a.
  Atemnot, die sich nicht bessert, brodelndes Atemgeräusch,
  Thoraxschmerz,
  Schocksymptome,
  entgleiste Vitalwerte,
  schwere Verletzungen,
  Hirndruckzeichen, etc.
  \item \T{Alarm-Diagnosen}: Manche Diagnosen
  sind automatisch mit einer vitalen Bedrohung gleichzusetzen,
  auch wenn eventuell Symptome fehlen oder nur schwach
  ausgeprägt sind. Dazu gehören beispielsweise der
  Herzinfarkt, das kardiale oder toxische Lungenödem,
  der Status Epilepticus, eine Beckenfraktur, u.\,s.\,w.
\end{itemize}
Bei vital bedrohten Patienten werden die \EE{Standardmaßnahmen 
für vital bedrohte Patienten} ergriffen,
sofern das entsprechende Material für den Ersthelfer
verfügbar bzw. eine entsprechende Ausbildung vorhanden ist.
% 
}
{\vspace{1em}

\fcolorbox{DefaultRed}{White}{
\begin{minipage}[t]{\linewidth - 4\Korrektur}
{
\begin{itemize}
    \item Massive Störungen von Vitalfunktionen
    \item Alarm-Symptome
\item Alarm-Diagnosen
\end{itemize}



}
\end{minipage}
}
}{./images/teaser/minen_hinweis2.jpg}
{Vorsicht!}
{GABS.}
\label{topic:standardmassnahmen-vital}%
\MassnahmenNG{Standardmaßnahmen}{Standardmaßnahmen bei vital
bedrohten
Patienten}{\label{m:ehi-standardmassnahmen-vital}}{

\begin{enumerate}
  \item \EE{Situationsgerechte Lagerung}
  \item \EE{Sauerstoffgabe}: je nach Indikation,
  allgemeiner Zielwert: Sp\ce{O2} von \VonBis{94}{98}\PC*
  \item Ggfs. \EE{Notarzt-Nachforderung}:
  mit kurzer Begründung
  
  \Z{Bei manchen zeitkritischen Notfällen kann vom 
  \E{professionellen Helfer} von der Notarztnachforderung 
  zugunsten eines raschen Transportes in eine geeignete 
  Einrichtung abgesehen werden.}
  \item \EE{Engmaschiges, bestmögliches Monitoring}: Je nach
  vorhandenem Material
  
  (Blutdruck, Herzfrequenz, Pulsoxymetrie, EKG,
  Sitzwache/Patientenbeobachtung, etc.\footnote{Zuerst werden die Patientenbeobachtung
  und -- wenn vorhanden -- die
  Pulsoxymetrie eingesetzt. Überwachungsgeräte, deren Anlage
  zeitintensiv ist (EKG, \ldots), sollen erst dann verwendet
  werden, wenn alle dringlicheren Maßnahmen durchgeführt und
  der Einschätzungsblock beendet ist! Die
  Patientenbeobachtung bleibt immer ein wesentlicher Teil
  des Monitorings!}
  \item \EE{Reanimationsbereitschaft} herstellen
  (Platz schaffen, \E{Defibrillator} und ggfs.
  Absauggerät organisieren)
\end{enumerate}



{\scriptsize(Reihenfolge zählt!)}
}

}

\subsection{Wunden und andere spezielle Situationen}
\index{Wunden!Erste-Hilfe}
\label{topic:wunden-ehi}


Jede Wunde ist für unseren Körper eine Gefahr, auch wenn sie noch so klein
ist. Grund dafür ist das mögliche Eindringen von Krankheitserregern wie z.\,B.
Tetanus. 

\TextSlide{Bagatellverletzungen}
{%
Bagatellverletzungen sind \EE{kleine, oberflächliche
Wunden}, wie z.\,B. Schürfwunden und kleine Schnitte, welche
nicht tiefer als \nicefrac{1}{2}\CM* und nicht länger als
2\CM* sind, gemeint. Bei idealer Behandlung durch den
Ersthelfer müssen diese nicht von einem Arzt versorgt
werden.
}{
\begin{itemize}
  \item kleine, oberflächliche Wunden
\end{itemize}
}{}{}{}


\protect\VARIANT{VAR-STANDARD}{

\MassnahmenMK{Erste-Hilfe-Maßnahmen}{Bagatellverletzungen}{m:ehi-bagatellverletzungen}{
}{}
\SlideCompensationNoParskip{1}
}
\protect\VARIANT{VAR-ERNAEHRUNGSPAED}{

\MassnahmenNG{Erste-Hilfe-Maßnahmen}{Bagatellverletzungen}{\label{m:ehi-bagatellverletzungen}}{

\begin{itemize}
  \item Wunde reinigen (Leitungswasser)
  \item Evtl. steriler Wundverband
  \item Auf den Tetanusschutz achten
\end{itemize}
}
}


\VARIANT{VAR-ERNAEHRUNGSPAED}{
\subsubsection{Einteilung von Wunden}
\label{topic:wunden-einteilung}



Wunden lassen sich \E{nach verschiedenen Gesichtspunkten}
einteilen bzw. unterscheiden. In den folgende Abschnitten
sind Einteilungen nach der physikalischen
Verletzungsursache, nach der Tiefe der Schädigung sowie nach
der Art der Schädigung (Wundart) angeführt.




\TextSlide{Physikalische Verletzungsursache}{%
Eine Einteilung bezieht sich auf die \EE{physikalische
Verletzungsursache} \cite[S.1796]{Pschy259}:

\begin{enumerate}
  \item \T{Mechanische Wunden}:
  \index{Wunden!mechanische}Entstehen durch spitze oder
  stumpfe Gewalteinwirkung auf den Körper (Stichwunde,
  Schürfwunde, Schusswunde, Bisswunde, ...)
  \item \T{Thermische Wunden}: \index{Wunden!thermische}
  Entstehen durch Kälte- oder Hitzeeinwirkung auf den Körper
   (Verbrennungen, Stromeinwirkung, Erfrierung, ...)
  \item \T{Chemische Wunden}:
  \index{Wunden!chemische}Entstehen durch Verätzungen mit
  Laugen oder Säuren
  \item \T{Strahlenbedingte Wunden}: \index{Wunden!strahlenbedingte} Entstehen
  durch Ultraviolettstrahlung oder ionisierende Strahlung
  (radioaktive Strahlung, Röntgenstrahlung)
\end{enumerate}
}{
\begin{itemize}
  \item Mechanische Wunden
  \item Thermische Wunden
  \item Chemische Wunden
  \item Strahlenbedingte Wunden
\end{itemize}
}{}{}{}




\TextSlide{Tiefe der
Schädigung}{%
Eine weitere Unterteilung kann nach der \EE{Tiefe der
Schädigung} des Gewebes vorgenommen werden.

\begin{itemize}
  \item \EE{Oberflächliche Wunden:} Leichte Verletzungen,
  welche max. bis zur Lederhaut reichen.
  \item \EE{Tiefe Wunden:} Neben der Verletzung der Haut
  können hier auch Muskeln, Sehnen, größere Blutgefäße
  innere Organe und Knochen betroffen sein \cite{LPN3}.
  \item \EE{Penetrierende Wunden:} Penetrierende Wunden
  sind tiefe Wunden mit Eröffnung mindestens einer
  Körperhöhle (Schädel, Thorax, Bauch) \cite{LPN3}.
\end{itemize}
}{
\begin{itemize}
  \item Oberflächliche Wunden
  \item Tiefe Wunden
  \item Penetrierende Wunden
\end{itemize}
}{}{}{}






Unterschiedliche \E{Entstehungsmechanismen} erzeugen
unterschiedliche \E{Wundarten} bzw. \E{Wundformen}.

\Layout{57}{Wundarten}
{%

}{

}{
\begin{tabulary}{\linewidth}{LL}\toprule\addlinespace
% \addlinespace
\noindent\EE{Schnittwunde}
&
Glatte Wundränder; tiefer liegende Nerven und Blutgefäße können verletzt sein
\\\addlinespace
\noindent\EE{Stichwunde}
&
Glatte Wundränder; meist kleine Wunde dafür aber tiefer, kann mitunter in
Körperhöhlen penetrieren; Das harmlose Aussehen kann über
das tatsächliche Verletzungsausmaß hinweg täuschen!
\\\addlinespace 
\noindent\EE{Risswunde}
&
Unregelmäßige Wundränder; meist nur Schädigung der Haut und
	keine tieferliegenden Nerven und Blutgefäßen
\\\addlinespace
\noindent\EE{Quetschwunde}
&
Unregelmäßige Wundränder; es kann zum Aufplatzen der Haut kommen (Platzwunde)
\\\addlinespace
\noindent\EE{Rissquetschwunde} (RQW)
&
Kombination aus Riss- und Quetschwunde
\\\addlinespace
\noindent\EE{Schürf- und Kratzwunde} 
&
Oberflächliche Wunde
\\\addlinespace
\noindent\EE{Bisswunden} 
&
Kombination aus Riss-, Stich- und Quetschwunde
\\\addlinespace
\noindent\EE{Ablederung} 
&
Großflächiger Gewebedefekt, z.\,B. Skalpierung, große Lappenwunden
\\\addlinespace
\noindent\EE{Amputation} 
&
Abtrennung eines Körperteiles
\\\addlinespace
\noindent\EE{Pfählung~/ Einspießung} 
&
Eindringen (oder Durchdringen) und Verbleiben von spitzen,
u.\,U. unregelmäßigen Fremdkörpern (z.\,B. Messer,
Injektionsnadel) \\\addlinespace 
\noindent\EE{Schusswunde} 
&
Meist klein aussehende Wunde, das eigentliche Schadensausmaß im Körperinneren
ist meist nicht ersichtlich (Druckwelle kann starke innere
Verletzungen verursachen),
\Z{evtl. Ausschuss größer als Einschuss}
\\\addlinespace
\noindent\EE{(Hämatom)} 
&
Bluterguss; Ansammlung von Blut im Gewebe außerhalb der Blutgefäße 
\\\addlinespace\bottomrule
\end{tabulary}
}{Wundarten.}{}

}


\TextSlide{Wunden, die unbedingt ärztlich versorgt werden müssen.}
{
Alle anderen Wunden, welche tiefer, größer etc. sind, müssen
unbedingt von einem Arzt versorgt werden. Für diese Fälle
gibt es \E{zwei Gebote} für den Ersthelfer, welche er
unbedingt einhalten muss:

\begin{enumerate}
    \item \EE{Es kommt nichts in die Wunde
hinein}, wie z.\,B. Wundpuder oder Wundsalben. Dies muss sonst im
Spital wieder mühevoll aus der Wunde entfernt werden, um sie
gründlich zu reinigen.
\item \EE{Es kommt nichts aus der Wunde heraus.}
Fremdkörper, welche ein Blutgefäß durchtrennen, können
dieses evtl. abdichten. Beim Entfernen kann es dann zu
starken Blutungen kommen.

\end{enumerate}
 Die anschließende
Tabelle gibt einen groben Überblick über Wunden, welche unbedingt von einem
Arzt versorgt werden müssen.
}{
\begin{itemize}
  \item Zwei Gebote:
  \begin{itemize}
    \item Es kommt nichts hinein
    \item Es kommt nichts heraus
  \end{itemize}
\end{itemize}
}{}{}{}

\Layout{57}{}
{}{}
{
\begin{tabulary}{\linewidth}{LL}
\addlinespace\toprule\addlinespace
\noindent\EE{Kein ausreichender Tetanus-Impfschutz}
&
Impfung innerhalb von 72 Stunden notwendig
\\\addlinespace
\noindent\EE{Klaffende Wunden}
&
\ldots\ müssen oft genäht werden. 
\\\addlinespace
\noindent\EE{Verbrennungen, Verbrühungen}
&
Einschätzung der Bedrohung anhand von Verbrennungsgrad und Ausdehnung
(\prettyref{topic:verbrennung-ehi}) 
\\\addlinespace
\noindent\EE{Wunden mit Fremdkörper}
&
sofern keine Bagatellverletzung (z.\,B. Schiefer, \ldots )
\\\addlinespace
\noindent\EE{Augenverletzungen aller Art}
&
Bei Augenverletzungen immer beide Augen verbinden
\\\addlinespace
\noindent\EE{Stichwunden}
&
Vom Ersthelfer kann nicht festgestellt werden wie tief diese sind
\\\addlinespace
\noindent\EE{Wunden an Gelenken, Hand- und Fußrücken}
&
\ldots\ mit gleichzeitiger Bewegungseinschränkung
\\\addlinespace
\noindent\EE{Bisswunden}
&
Infektionsgefahr, nicht nur von Tieren, sondern auch von Menschen
\\\addlinespace
\noindent\EE{Insektenstiche}
&
Bei allergischen Reaktionen oder bekannter entsprechender Allergie
\\\addlinespace
\noindent\EE{Schusswunden}
&

\\\addlinespace
\noindent\EE{Wunden im Genitalbereich}
&

\\\addlinespace\bottomrule
\end{tabulary}
}
{Wunden, die unbedingt vom Arzt versorgt werden müssen!}{}



\subsubsection{Verbrennungen, Verbrühungen}
\index{Verbrennung!Erste-Hilfe}
\label{topic:verbrennung-ehi}

\TextSlide{Einteilung
}{%
\index{Verbrennungspanzer}%
Um festzustellen, wie bedrohlich die Verbrennung für den
Patienten ist, muss der Schweregrad der Schädigung der Haut
eingeteilt werden. Dies geschieht mittels drei Grade. Beim
\EE{1.~Grad} kommt es zur Rötung der Haut, einer Schwellung
des Gewebes und zu Schmerzen. Der \EE{2.~Grad} ist
charakterisiert
durch Blasenbildung. Es bestehen weiterhin
Schmerzen. Beim \EE{3.~Grad} kommt es zur \T{Schorfbildung}
(\E{Verbrennungspanzer}). Die Schmerzen sind abhängig von
der Tiefe der Schädigung der Haut. }{
\begin{itemize}
  \item 1\textdegree : Rötung, Schwellung, Schmerzen
  \item 2\textdegree : Blasenbildung, Schmerzen
  \item 3\textdegree : Schorfbildung, Verbrennungspanzer
\end{itemize}
}{}{}{}



\TextSlide{Ausdehnung: Handregel und Neunerregel}
{
Genauso wichtig wie der Schweregrad ist die Ausdehnung der
Verbrennung. Damit ist gemeint, wieviel Prozent der Körperoberfläche betroffen sind.
Dazu gibt es zwei Möglichkeiten. Einerseits die \T{Handregel}, welche
besagt, dass die \EE{Handfläche des Patienten  1\PC*}
seiner Körperoberfläche entspricht. Da dies in der Praxis
bei großflächigen Verbrennungen meist umständlich ist, gibt
es auch noch die \T{Neunerregel}. Hier wird der Körper in
9\PC*-Teile eingeteilt.

}{
\begin{itemize}
  \item Handregel
  \item Neunerregel
\end{itemize}
}
{}
{}{}

\begin{table}[H]
\caption{Neunerregel}
\begin{tabulary}{\linewidth}[H]{LCC}
\toprule
 & \% &
\\\midrule
Kopf:
 & 9  &
\\
Arm, je:  &  9  &
\\
Rumpf vorne: & 18  &
\\
Rumpf hinten: & 18 &
\\
Bein, je: & 18 &
\\
Genitalien:  & 1 &
\\\bottomrule
\end{tabulary}
\end{table}

\Cave{\Ca{Schockgefahr!} 

Bei Kindern ist ab \EE{5\PC*}, bei Erwachsenen ab \EE{10\PC*}
mind. 2\textdegree{} verbrannter Körperoberfläche mit einem Schock zu
rechnen!}


\protect\VARIANT{VAR-STANDARD}{
\MassnahmenMK{Erste-Hilfe-Maßnahmen}{Verbrennungen}{m:verbrennung-eh}{
}{}
}
\protect\VARIANT{VAR-ERNAEHRUNGSPAED}{
\MassnahmenNG{Erste-Hilfe-Maßnahmen}{Verbrennungen}{\label{m:verbrennung-eh}}{
\begin{itemize}
  \item Hitzeeinwirkung sofort stoppen.
  \item Wenn 
  \begin{itemize}
    \item innerhalb von 2 Minuten
  möglich,
    \item Ausdehnung <\,20\PC* (Kinder <\,10\PC, Säuglinge
    <\,10\PC*) 
  \end{itemize}  für \E{maximal 10\Min* mit} mit
    handwarmen Wasser kühlen (Richtwert 20\,\DegC); 
    
    andernfalls ist keine Kühlung durchzuführen
  \cite{Asanger2010,ASBOEAkademieUpdate2010Verbrennung}.
  
  \Ca{Achtung: Gefahr der Unterkühlung!}
  \item Wärmeerhalt ist besonders wichtig!
  \item Keimfreier Wundverband
  \item Schockbekämpfung
  \item \Ca{Keine Salben oder Puder etc. verwenden}
\end{itemize}
}
}


% \clearpage % TODO temporäres clearpage
\subsubsection{Erfrierungen}

\Definition{Durchblutungsstörungen mit Gewebstod an den Körperenden, 
die durch lokale 
Durchblutungsstörungen aufgrund von  Kälteeinwirkung (Eis,
Wind, Nässe,  Kälte)
entstanden sind.}


\protect\VARIANT{VAR-STANDARD}{
\MassnahmenMK{Erste-Hilfe-Maßnahmen}{Erfrierungen}{m:ehi-erfrierungen}{
}{}
}
\protect\VARIANT{VAR-ERNAEHRUNGSPAED}{
\MassnahmenNG{Erste-Hilfe-Maßnahmen}{Erfrierungen}{\label{m:ehi-erfrierungen}}{

\begin{itemize}
  \item Keimfreier Wundverband
  \item Restlichen Körper aufwärmen
  \item Warme, gezuckerte Getränke verabreichen
  \item Keinen Alkohol zu trinken geben!
  \item Nicht frottieren! 
\end{itemize}
}
}


\clearpage
\subsubsection{Unterkühlung}

\Definition{Absinken der Körperkerntemperatur unter 36\,\DegC{}}

\TextSlide{\TxSymptome{} und Kennzeichen}
{
Am Beginn hat der Patient \EE{Schmerzen}, Atem- und
Herzfrequenz sind erhöht.
Je weiter die Temperatur sinkt, desto \EE{teilnahmsloser}
und \EE{müder} wird der Patient. Zu Letzt kommt es zum
\EE{Absinken von Atem- und Herzfrequenz}. Bei einer
Temperatur unter 27\,\DegC{} tritt der Tod ein.

\bigskip

\bigskip
}{
\begin{itemize}
  \item Schmerzen
  \item Teilnahmslosigkeit/Müdigkeit
  \item Absinken von Atem- und Herzfrequenz
\end{itemize}
}{}{}{}


\protect\VARIANT{VAR-STANDARD}{
\MassnahmenMK{Erste-Hilfe-Maßnahmen}{Unterkühlung}{m:ehi-unterkuehlung}{
}{}
}
\protect\VARIANT{VAR-ERNAEHRUNGSPAED}{\SlideCompensationNoParskip{1}

\MassnahmenNG{Erste-Hilfe-Maßnahmen}{Unterkühlung}{\label{m:ehi-unterkuehlung}}{

\begin{itemize}
  \item Aufwecken und wach halten
  \item Vorsichtig in Decken oder Kleidungsstücke einwickeln
  \item Rettungsdienst verständigen
  \item Warme Getränke verabreichen, wenn Patient bewusstseinsklar ist
  \item Keinen Alkohol zu trinken geben!
  \item Keinesfalls bewegen oder frottieren: sog.
  \Ca{Bergungstod!}\footnote{Unter dem Begriff
  \emph{Bergungstod} versteht man das Versterben des
  Patienten während der Rettung in Folge eines
  Blutrückflusses aus der kalten Körperschale in den
  wärmeren Körperkern, welcher durch die Bewegung des
  Patienten während der Rettung ausgelöst wird.}
\end{itemize}
}
}


\subsubsection{Verstauchung}
\label{topic:verstauchung-eh}

\Definition{Überdehnung der Bänder bzw. des Kapselapparates eines Gelenkes.}

\TextSlide{\TxSymptome{} und Kennzeichen}
{%
Der Patient hat \EE{Schmerzen} im Bereich des betroffenen Gelenkes, es
kann auch zu einer \EE{Schwellung} rund um die geschädigte Stelle kommen. Mit
unter nicht gleich vorhanden, sondern erst am nächsten Tag ersichtlich,
kann sich auch ein \EE{Bluterguss} bilden. Auch eine \EE{Einschränkung der
Bewegung} ist -- schmerzbedingt -- möglich. 
\bigskip
\bigskip

\bigskip
}{
\begin{itemize}
  \item Schmerzen
  \item Schwellung
  \item Bluterguss
  \item Bewegungseinschränkung
\end{itemize}
}{}{}{}


\protect\VARIANT{VAR-STANDARD}{
\MassnahmenMK{Erste-Hilfe-Maßnahmen}{Verstauchung}{m:ehi-verstauchung}{
}{}
}
\protect\VARIANT{VAR-ERNAEHRUNGSPAED}{
\MassnahmenNG{Erste-Hilfe-Maßnahmen}{Verstauchung}{\label{m:ehi-verstauchung}}{

\begin{itemize}
  \item Ruhigstellung
  \item Kalte Umschläge
  \item Unfallabteilung aufsuchen
  \item Hochlagern % ZWING 2010-03-xx
\end{itemize}
}
}



\subsubsection{Verrenkung}
\label{topic:verrenkung-eh}

\Definition{Gelenkverletzung, bei der zwei das Gelenk bildende Knochenenden
aus ihrer normalen Stellung zueinander verschoben werden. Dabei kommt es oft zu
Begleitverletzungen an der Gelenkkapsel bzw. den Gelenkbändern. \cite{Gorgass7}}



\TextSlide{\TxSymptome{} und Kennzeichen}
{%
Der Patient hat starke \EE{Schmerzen} im Bereich des verletzten Gelenkes,
die betroffene Extremität kann nicht mehr bewegt werden
(\EE{Bewegungsunfähigkeit}). Außerdem kommt es zu einer \EE{abnormen
Stellung}, welche \EE{federnd fixiert} gelagert ist. 

\bigskip

\bigskip

\bigskip
}{
\begin{itemize}
    \item Schmerzen
    \item Bewegungsunfähigkeit
    \item Abnorme Stellung, \enquote{federnd fixiert}
\end{itemize}
}{}{}{}


\protect\VARIANT{VAR-STANDARD}{
\MassnahmenMK{Erste-Hilfe-Maßnahmen}{Verrenkung}{m:ehi-verrenkung}{
}{}
}
\protect\VARIANT{VAR-ERNAEHRUNGSPAED}{\SlideCompensationNoParskip{1}

\MassnahmenNG{Erste-Hilfe-Maßnahmen}{Verrenkung}{\label{m:ehi-verrenkung}}{

\begin{itemize}
    \item Vorgefundene Stellung beibehalten
    \item Ruhigstellung
    \item Rettungsdienst verständigen
    \item Keine Einrenkversuche!
\end{itemize}
}
}



\subsubsection{Knochenbrüche}
\label{topic:knochenbruch-eh}

\TextSlide{Arten}
{
Grundsätzlich wird zwischen zwei Arten von Knochenbrüchen
unterschieden: Beim \EE{geschlossenen Knochenbruch} ist die
Haut über der Bruchstelle noch intakt, während beim
\EE{offenen Knochenbruch} aufgrund der Gewalteinwirkung von
außen die Haut über der Bruchstelle verletzt ist. Es muss
dabei nicht der Knochen herausragen.


\bigskip
}{
\begin{itemize}
    \item Geschlossener Knochenbruch
    \item Offener Knochenbruch
\end{itemize}
}{}{}{}

\TextSlide{\TxSymptome{} und Kennzeichen}
{
Es gibt \E{sichere} und \E{unsichere} Kennzeichen für einen Knochenbruch. 

Zu
den \T{sicheren Kennzeichen} zählt die 

\begin{enumerate}
  \item \EE{Fehlstellung} (abnorme Stellung,
  Verkürzung, Treppenbildung)
  \item \EE{Abnorme Beweglichkeit}
  \item \EE{Krepitation} (das hörbare und fühlbare Aneinanderreiben
  der Knochen)
  \item \EE{Sichtbare freie Knochenteile} (offener Bruch)
\end{enumerate}



\T{Unsichere
Knochenbruchzeichen}\index{Knochenbruchzeichen!Unsichere}
sind:

\begin{enumerate}
  \item \EE{Schmerzen} 
  \item \EE{Schwellungen}
  \item \EE{Bewegungseinschränkungen} 
  \item \EE{Blutergüsse}
\end{enumerate}


}{
\begin{itemize}
    \item Sichere Kennzeichen
    \begin{compactitem}
        \item Fehlstellung
        \item Abnorme Beweglichkeit
        \item Krepitation
        \item Sichtbare Knochenteile
    \end{compactitem}
    \item Unsichere Kennzeichen
    \begin{compactitem}
        \item Schmerzen
        \item Schwellung
        \item Bewegungseinschränkung
        \item Bluterguss
     \end{compactitem}
\end{itemize}
}{}{}{}

\TextSlide{\TxGefahren{}}
{
Ein Knochen ist lebendes
Gewebe, daher ist er auch gut durchblutet. Eine der Gefahren
des Knochenbruchs ist daher der \EE{Blutverlust}. Je größer
der Blutverlust ist umso größer ist wiederum die Gefahr,
dass der Patient in den \EE{Schock} gerät. Bei einem offenen
Knochenbruch besteht außerdem die Gefahr einer
\EE{Wundinfektion}. Bei Rippenbrüchen kann
es zu \EE{Begleitverletzungen innerer Organe} kommen.
\bigskip

\bigskip
}{
\begin{itemize}
    \item Blutverlust
    \item Schock
    \item Wundinfektion
    \item Begleitverletzungen innerer Organe
\end{itemize}
}{}{}{}


\protect\VARIANT{VAR-STANDARD}{
\MassnahmenMK{Erste-Hilfe-Maßnahmen}{Knochenbrüche}{m:ehi-knochenbrueche}{
}{}
}
\protect\VARIANT{VAR-ERNAEHRUNGSPAED}{\SlideCompensationNoParskip{1}

\MassnahmenNG{Erste-Hilfe-Maßnahmen}{Knochenbrüche}{\label{m:ehi-knochenbrueche}}{

\begin{itemize}
    \item Ruhig liegen lassen
    \item Wundversorgung bei offenen Knochenbrüchen
    \item Beengende Kleidungsstücke lockern
    \item Schuhe öffnen, nicht ausziehen
    \item Ruhigstellung durch unterstützte Lagerung
    \item Schockbekämpfung
\end{itemize}
}
}


\subsubsection{Vergiftungen}
\label{topic:vergiftungen-ehi}

\Text{Allgemeines}
{
Es gibt mehrere Möglichkeiten der Giftaufnahme. Man kann es
beispielsweise verschlucken, einatmen oder über die Haut
aufnehmen. Das Gift selbst kann aus dem Haushalt, der Umwelt
(Pflanzen, Pilze, \ldots), der Industrie, dem Gewerbe oder
der Landwirtschaft kommen.

}{
}{}{}{}

\TextSlide{\TxSymptome{} und Kennzeichen}
{
Je nach Art des Giftes und der Menge, die dem Körper zugeführt wurden, gibt es
unterschiedliche Kennzeichen. Es kann mit \EE{Rauschzuständen} und
\EE{Erbrechen bzw. Durchfall} beginnen, über \EE{Bewusstseinsstörung} zu
\EE{Atem- und Kreislaufstörung} führen.
\bigskip
\bigskip

\bigskip

\bigskip
}{
\begin{itemize}
    \item Rauschzustände
    \item Erbrechen, Durchfall
    \item Bewusstseinsstörung
    \item Atem- und Kreislaufstörung
\end{itemize}
}{}{}{}


\protect\VARIANT{VAR-STANDARD}{
\MassnahmenMK{Erste-Hilfe-Maßnahmen}{Vergiftungen}{m:ehi-vergiftungen}{
}{}
}
\protect\VARIANT{VAR-ERNAEHRUNGSPAED}{
\MassnahmenNG{Erste-Hilfe-Maßnahmen}{Vergiftungen}{\label{m:ehi-vergiftungen}}{
\begin{itemize}
    \item Patient nicht ansprechbar:
    \begin{itemize}
        \item Retten aus der Gefahr (Selbstschutz!)
        \item Lebensrettende Sofortmaßnahmen
    \end{itemize}
    \item Patient ansprechbar:
    \begin{itemize}
        \item Wenn das Gift bekannt ist, Vergiftungsinformationszentrale
        (\Code{01\,/\,406\,43\,43}) anrufen und \E{Expertenrat} einholen.
        \item Wenn Gift unbekannt: Betreuen und → Notruf/Krankenhaus.
        \item Selbstschutz!
        \item Giftstoffe sichern und mitnehmen.
    \end{itemize}
\end{itemize}
}
}


\subsubsection{Verätzungen}


\protect\VARIANT{VAR-STANDARD}{
\MassnahmenMK{Erste-Hilfe-Maßnahmen}{Verätzungen der
Haut}{m:ehi-veraetzungen-haut}{
}{}\SlideCompensationNoParskip{1}
}
\protect\VARIANT{VAR-ERNAEHRUNGSPAED}{
\MassnahmenNG{Erste-Hilfe-Maßnahmen}{Verätzungen der
Haut}{\label{m:ehi-veraetzungen-haut}}{

\begin{itemize}
   \item \Ca{Selbstschutz} beachten!
    \item Mit ätzender Substanz getränkte Kleidung entfernen
    \item Haut mit reinem Wasser abspülen
    \item Keimfreier Wundverband
    \item Evtl. Schockbekämpfung
\end{itemize}
}
}


\protect\VARIANT{VAR-STANDARD}{
\MassnahmenMK{Erste-Hilfe-Maßnahmen}{Verätzungen der
Augen}{m:ehi-veraetzungen-augen}{
}{}\SlideCompensationNoParskip{1}
}
\protect\VARIANT{VAR-ERNAEHRUNGSPAED}{
\MassnahmenNG{Erste-Hilfe-Maßnahmen}{Verätzungen der
Augen}{\label{m:ehi-veraetzungen-augen}}{

\begin{itemize}
    \item Auge mit reinem Wasser von innen nach außen ausspülen. Abwasser
    \E{nicht über das gesunde Auge} abrinnen lassen
    \item Auge mit keimfreiem Wundverband bedecken; zum Ruhigstellen \E{beide}
    Augen bedecken
    \item Selbstschutz beachten!
\end{itemize}
}
}


\protect\VARIANT{VAR-STANDARD}{
\MassnahmenMK{Erste-Hilfe-Maßnahmen}{Verätzungen des
Verdauungstraktes}{m:ehi-veraetzungen-verdauungstrakt}
{
}{}\SlideCompensationNoParskip{1}
}
\protect\VARIANT{VAR-ERNAEHRUNGSPAED}{
\MassnahmenNG{Erste-Hilfe-Maßnahmen}{Verätzungen des
Verdauungstraktes}{\label{m:ehi-veraetzungen-verdauungstrakt}}
{ 
\begin{itemize}
    \item Mund mit Wasser ausspülen
    \item Nicht erbrechen lassen!
    \item Substanz bekannt → Vergiftungsinformationszentrale anrufen
    \item evtl. Schockbekämpfung
    \item Selbstschutz beachten!
\end{itemize}
}
}



\TODO{????}{0.8.0}{Fehlt: EHI Tierbisse, Tollwut, Schlangenbiss, ZEckenbiss,
Insektenstiche Mund/Rachenraum}{}
\subsection{Wundverbände}
\label{topic:wundverbaende}

\TextSlide{\TxBeschreibung}{
Ein Wundverband besteht aus einer keimfreien Wundauflage und einer Befestigung.
}{
\begin{itemize}
  \item keimfreie Wundauflage
  \item Befestigung
\end{itemize}
}{}{}{} 

\subsubsection{Pflasterwundverband}
\Layout{20}{\TxBeschreibung}{
Pflasterverbände können nahezu überall eingesetzt werden,
sie müssen dazu unter Umständen passend zugeschnitten
werden. Wenn ein Pflasterverband über ein bewegliches
Gelenk angelegt werden muss, wird an der beweglichen Stelle
der Klebestreifen von außen keilförmig eingeschnitten.
}{}{./images/koch-aass/00800-gamma/Pflasterb.JPG}{}{}

\subsubsection{Dreiecktuchverbände}

Mittels des Dreiecktuchs können verschiedene Verbände
angelegt werden, typische Beispiele sind:

\begin{itemize}
  \item Kopfverband \dotfill
  \FullPrettyRef{fig:mitella-kopfverband}
  \item Schulterverband \dotfill
  \FullPrettyRef{fig:mitella-schulterverband}
  \item Handverband \dotfill \FullPrettyRef{fig:mitella-handverband}
  \item Knieverband \dotfill \FullPrettyRef{fig:mitella-knieverband}
\end{itemize}




\begin{PictureSeriesWide}{}{Bilderserie:
Kopfverbände\label{fig:mitella-kopfverband}}
\PictureSeriesRow{3}
{./images/koch-aass/00800-gamma/Kopfverband1b.JPG}
{Verband 1}
{./images/koch-aass/00800-gamma/Kopfverband2b.JPG}
{Verband 2}
{empty}
{}
{}
{}
\end{PictureSeriesWide}

\begin{PictureSeriesWide}{}{Bilderserie:
Schulterverband\label{fig:mitella-schulterverband}}
\PictureSeriesRow{3}
{./images/koch-aass/00800-gamma/Schulterverband1b.JPG}
{Schritt 1}
{./images/koch-aass/00800-gamma/Schulterverband2b.JPG}
{Schritt 2}
{./images/koch-aass/00800-gamma/Schulterverband3b.JPG}
{Schritt 3}
{}
{}
\end{PictureSeriesWide}



\begin{PictureSeriesWide}{}{Bilderserie: Zwei verschiedene
Handverbände\label{fig:mitella-handverband}}
\PictureSeriesRow{3}
{./images/koch-aass/00800-gamma/Handverband1b.JPG}
{Schritt 1}
{./images/koch-aass/00800-gamma/Handverband2b.JPG}
{Schritt 2}
{./images/koch-aass/00800-gamma/Handverband3b.JPG}
{Schritt 3}
{}
{}
\PictureSeriesRow{3}
{./images/koch-aass/00800-gamma/Handkrawatte1b.JPG}
{Alternative: Schritt 1}
{./images/koch-aass/00800-gamma/Handkrawatte2b.JPG}
{Schritt 2}
{./images/koch-aass/00800-gamma/Handkrawatte3b.JPG}
{Schritt 3}
{}
{}
\end{PictureSeriesWide}

\begin{PictureSeriesWide}{}{Bilderserie:
Knieverband\label{fig:mitella-knieverband}}
\PictureSeriesRow{3}
{./images/koch-aass/00800-gamma/Knieverband1b.JPG}
{Schritt 1}
{./images/koch-aass/00800-gamma/Knieverband2b.JPG}
{Schritt 2}
{./images/koch-aass/00800-gamma/Knieverband3b.JPG}
{Schritt 3}
{}
{}
\end{PictureSeriesWide}







%%%%%%%%%%%%%%%%%%% ERNÄHRUNGSPÄDAGOGIK
\VARIANT{VAR-ERNAEHRUNGSPAED}{
\clearpage

\subsection{Koronare Herzkrankheit und Akutes Koronarsyndrom}
\index{Koronare Herzkrankheit}
\index{Herzkrankheit!Koronare}
\index{KHK}
\index{Akutes Koronarsyndrom}
\index{Koronarsyndrom!Akutes}
\label{topic:khk}
\label{topic:aks}
\label{topic:akutes-koronarsyndrom}

\Text{Krankheitsbilder}
{Die \T{Koronare Herzkrankheit} (\T{KHK})] ist eine
Erkrankung der \EE{Herzkranzgefäße} (Koronargefäße) bei der
es zu einer Verengung der Gefäße kommt. Das 
\T[Akutes Koronarsyndrom]{Akute Koronarsyndrom}
\index{Koronarsyndrom!akutes} 
ist ein
Sammelbegriff für \E{Herzinfarkt}, \E{stabile} und
\E{instabile Angina pectoris}. Bei diesen Krankheitsbildern
kommt es zu einer akuten \EE{Unterversorgung des
Herzmuskels} mit Blut und Sauerstoff aufgrund der
Durchblutungsstörung. Die Ursache sind in erster Linie
verengte oder verstopfte Herzkranzgefäße
(\I[Ischämie!Akutes Koronarsyndrom]{Ischämie}). 
Das akute Koronarsyndrom
ist häufig das Ergebnis einer vorbestehenden koronaren Herzkrankheit.}
{}
{}
{}
{}

\Commentary[Koronare Herzkrankheit]{Eine koronare
Herzkrankheit ist oft anamnestisch bereits bekannt. Ein
Hinweis kann auch die Einnahme von ASS-Präparaten (z.\,B.
ThromboASS{\texttrademark}, HerzschutzASS{\texttrademark},
\ldots) oder ein Stent-Ausweis sein. 
Je nachdem welche Koronargefäße betroffen sind, spricht man
auch von einer 1-, 2- oder 3-Gefäßerkrankung (\Lang{engl.}
1-, 2-, 3-Vessel-Disease; bzw. \Lang{Abkz.} 1-, 2.-, 3.-VD).}

\begin{PictureSeriesWide}{}{Bilderserie:
\EEE{Koronargefäße}}
\PictureSeriesRowWide{3}
{./images/hirtler-aass/Herz-edited-0800.png}
{Das Herz mit seinen Koronargefäßen \SourcePicture{Hirtler}}
{./images/wikipedia-cc-by/Hk_coronary_bionerd.jpg}
{Darstellung der
Herzkranzgefäße während einer Herzkatheteruntersuchung.
Dabei wird über die Leistenarterie ein Katheter
eingebracht und bis knapp vor das Herz zu den Abgängen der
Koronargefäße vorgeschoben und anschließend ein
Kontrastmittel appliziert.
\SourcePicture{WmCo
\enquote{Bionerd}, CC-BY-3.0}}
{./images/wikipedia-cc-by/Heart_coronary_artery_lesion-lq.jpg}
{Die Koronargefäße versorgen den
Herzmuskel von außen nach innen. \SourcePicture{Patrick J.
Lynch, CC-BY-2.5}} {./images/wikipedia-cc-by/Hk_coronary_bionerd.jpg}
{empty}
{}

\end{PictureSeriesWide}


\subsubsection{Angina pectoris}
\index{Angina pectoris}
\label{topic:angina-pectoris}


\HeadingSub{\Lang{lat.} \enquote{Brustenge}}

\Definition{Plötzliches Engegefühl in der Brust meist mit
Thoraxschmerz infolge einer \underline{vorübergehenden}
Unter\-ver\-sorgung des Herzmuskels mit sauerstoffreichem Blut.}

\TextSlide{\TxBeschreibung}
{%
Es sind \E{zwei Auslösemechanismen} möglich:
Einerseits kann eine Minderdurchblutung des Herzmuskels
(Myokard)
  aufgrund einer oder mehrerer \E{Einengung(en) von
  Herzkranzgefäßen} (Koronararterien) die Ursache sein, oder
es kommt zu einem \E{erhöhtem Sauerstoffbedarf} des Herzens
bei
  Anstrengung.
In beiden Fällen kommt es zu einem \EE{Missverhältnis
zwischen Sauerstoffangebot und Sauerstoffbedarf} des
Herzens(\I[Ischämie!Angina pectoris]{Ischämie}), 
die Folge sind Schmerzen,
Atemnot und ein Leistungsverlust!

Der akute Angina-Pectoris-Anfall gehört zu der Familie des
\EE{Akuten Koronarsyndroms}. Dazu gehört auch der
Herzinfarkt, dem der gleiche Mechanismus wie der Angina Pectoris zu
Grunde liegt, wobei es beim Infarkt im \E{Gegensatz} zur Angina
pectoris zu einem \E{totalen} Herzkranzgefäß\E{verschluss}
mit Absterben von Herzmuskelgewebe kommt.
\A{KOCH: nein; GABS: Doch. }


}{
Missverhältnis zwischen Sauerstoffangebot und  
Sauerstoffbedarf
\begin{itemize}
  \item Einengung(en) von Herzkranzgefäßen
  \item Erhöhter Sauerstoffbedarf bei Anstrengung
\end{itemize}
}{}{}{}


\Layout{20}{\TxSymptome}
{%
Das Leitsymptom ist der \EE{akute, nicht atemabhängige,
Brustschmerz}. Er ist meist hinter dem Brustbein
lokalisiert und \EE{strahlt häufig aus}, häufig in den
Kiefer, (linken) Arm und in den Rücken. Unter Umständen
kann es auch zu einem \EE{Oberbauchschmerz} kommen,
welcher unter Umständen als Verdauungsstörung fehlgedeutet
werden kann. Der Schmerz wird als brennend bzw. stechend
beschrieben. Dazu kommt i.\,d.\,R. ein \EE{Druckgefühl} am
Brustkorb, \emph{\enquote{als würde jemand auf dem Brustkorb
stehen.}} \EE{Todesangst und Vernichtungsgefühl} sind
häufig.
Ein weiteres wichtiges Leitsymptom ist die \EE{Atemnot}
(\E{Dyspnoe}).

Ein akuter Angina pectoris-Anfall ist primär \EE{nicht
sicher von einem Herzinfarkt unterscheidbar}. Die Symptome
treten eher bei Belastung auf und vergehen oft bei
Entspannung, das Vorübergehen der Beschwerden darf aber
natürlich nicht abgewartet werden!
}{
\begin{itemize}
  \item Brustschmerz, brenndend/stechend, Druckgefühl
  \item Oft Ausstrahlung in den Kiefer,
  linken Arm, Rücken
  \item Evtl. Oberbauchschmerzen
  \item Todesangst, Vernichtungsgefühl
  \item Dyspnoe
  \item Primär \EE{nicht sicher von MCI unterscheidbar}
  \item Symptome eher bei Belastung, vergehen oft bei Entspannung
\end{itemize}

}
{./images/hirtler-aass/Herzschmerz.pdf}
{Schmerzausdehnung beim ischämischen
Herzschmerz.}
{Hirtler.}




\subsubsection{Herzinfarkt}
\HeadingSub{\Lang{Term.} \T{Myokardinfarkt}, \Lang{Abkz.} \T{MCI}}
\index{Herzinfarkt}
\index{Infarkt!Herz-}
\index{Myokardinfarkt}
\index{MCI}
\index{Herzkranzgefäße!Verschluss}
\index{Koronargefäße!Verschluss}
\label{topic:herzinfarkt}
\label{topic:myocardinfarkt}
\label{topic:mci}

% \DefinitionExp{Herzinfarkt}

\Definition{Verschluss eines oder mehrerer Herzkranzgefäße mit
  anschließendem Absterben des betroffenen Herzmuskelgewebes
  (\I[Ischämie!Herzinfarkt]{Ischämie}).}

\Layout{20}{\TxSymptome}
{%
Die Symptome entsprechen denen der Angina
pectoris, doch vergehen diese nicht mehr und es bleiben --
selbst bei optimaler Behandlung -- Schäden am Herzmuskel
bestehen. Bei Diabetikern kann es aufgrund der
Nervenschädigungen auch zu einem schmerzlosen Infarkt
kommen, man spricht dabei auch vom \enquote{\T[stummer
Infarkt]{stummen Infarkt}}\index{Infarkt, Stummer}.
Sonst steht normalerweise der \E{stechende/brennende}, nicht
atemabhängige \II[Brustschmerz]{Brust-} oder
\II{Oberbauchschmerz} zusammen mit \II{Atemnot} und
\II{Todesangst} im Vordergrund. Der Brustschmerz kann in
Richtung \E{Magen}, rechter oder linker \E{Hand}, \E{Kiefer}
oder auch \E{Rücken} \EE{ausstrahlen}. Zusätzlich zu dem
Schmerz verspürt der Patient auch oft ein \EE{Engegefühl} im
Brustkorb, \Speech{\enquote{als würde jemand darauf
stehen}}. Die Atemnot kann massiv ausfallen.

Bei Auftreten eines \EE{kardiogenen Schocks} können die
entsprechenden Schockzeichen beobachtet werden: Blässe,
Kaltschweißigkeit, Hypotonie, etc.
}{
\begin{itemize}
  \item Symptome wie bei Angina Pectoris, diese vergehen aber nicht
  \begin{compactitem}
    \item Stechender/brennender \EE{Brustschmerz} mit
    Ausstrahlung oder \EE{Oberbauchschmerz}
    \item \enquote{Stummer Infarkt} bei Diabetikern
    \item Druckgefühl, \Speech{\enquote{als würde jemand darauf
stehen}}
    \item Todesangst
    \item \EE{Dyspnoe}
    \item Kardiogener Schock
  \end{compactitem}
\end{itemize}
}
{./images/hirtler-aass/Herzschmerz.pdf}
{Schmerzausdehnung beim ischämischen
Herzschmerz.}
{Hirtler.}
  
\TextSlide{\TxKomplikationen}
{%
\begin{itemize}
  \item \EE{Herzrhythmusstörungen} (alle Arten), bis hin
  zum Atem-/Kreislaufstillstand und Reanimation.
  Besonders häufig kommt es als Frühkomplikation zum
  plötzlichen Kammerflimmern.
  \item Kardiogener Schock \index{Kardiogener
  Schock!beim Herzinfarkt}\index{Schock!Kardiogener!beim
  Herzinfarkt}
\end{itemize}
}{
\begin{itemize}
  \item \EE{Lebensgefahr!}
  \item Jede Herzrhythmusstörung als Komplikation möglich!
  \item Kammerflimmern als Frühkomplikation häufig!
\end{itemize}

}{}{}{}

\begin{PictureSeriesWide}{}{Bilderserie:
\EEE{ACS und Herzinfarkt}}
\PictureSeriesRowWide{3}
{./images/hirtler-aass/Herzschmerz.pdf}
{Schmerzausdehnung beim ischämischen
Herzschmerz. \SourcePicture{Hirtler.}}
{./images/wikipedia-cc-by/Heart_ant_wall_infarction-00800.jpg}
{Abgestorbenes Muskelgewebe. \SourcePicture{Patrick J. Lynch,
CC-BY-2.5}}
{./images/wikipedia-pd/12_lead_generated_inferior_MI}
{\Z{Typische Veränderungen im EKG bei einem
Hinterwandinfarkt:
Deutliche ST-Hebungen in den Ableitungen \EKG{II}, \EKG{III}
und
\EKG{aVF}} \SourcePicture{WMC/PD}}
{empty}
{}
\end{PictureSeriesWide}


\subsubsection{Akutes Koronarsyndrom}

Da bei Behandlungsbeginn oft nicht klar ist, ob es sich um
einen Herzinfarkt oder einen Angina-pectoris-Anfall handelt,
wird bei Vorliegen der entsprechenden Symptome (Atemnot,
Thoraxschmerz, Todesangst, \ldots) einfach vom \T{Akuten
Koronarsyndrom}\ZFootnote{\Lang{Engl} Acute Coronary
Syndrome, \Abk{ACS}.} gesprochen. 


\TextSlide{Weitere Therapie}
{%
Herzinfarktpatienten werden \enquote*{besonders} behandelt
und normalerweise nicht auf eine beliebige Internistische
Station gebracht\citep{Kalla2006}. Der Notarzt muss
i.\,d.\,R. zwischen dem Einleiten einer Lyse-Therapie und
der Einweisung in ein Herzkatheterlabor entscheiden.

\begin{itemize}
  \item Bei der \II{Lyse-Therapie} versucht man mittels eines
  speziellen Medikaments eine Verstopfung aufzulösen.
  \item \index{PTCA}
  \index{Herzkatheterlabor}
  Mit einem \EE{Herzkatheter}
   versucht man die
  verengten oder verstopften Gefäße aufzudehnen. Dazu
  benötigt man ein Spital, welches über ein entsprechend
  ausgestattetes und
  \E{dienstbereites}
  \T{Herzkatheterlabor} verfügt.\opt{REGVIE}{\footnote{In
  Wien gibt es einen \E{Dienstplan} für die Katheterlabore,
  d.h. jedes Katheterlabor hat an bestimmten Tagen der Woche
  rund um die Uhr Bereitschaftsdienst. Eine Voranmeldung und
  vorangehende Absprache zwischen Notarzt und diensthabendem
  Oberarzt der Abteilung ist erforderlich. Die PTCA ist in
  Wien im AKH, Donauspital, KH Hietzing, Rudolfstiftung,
  Hanusch-Krankenhaus und Wilhelminenspital verfügbar.(Angaben
  ohne Gewähr und können sich
  jederzeit ändern.)}} \hfill\Commentary[Herzkatheter]{%
  Bei
  einer Herzkatheteruntersuchung (\I{Koronarangiographie})
  wird über die Leistenarterie ein Katheter eingebracht und 
  bis knapp vor das Herz zu den Abgängen der Koronargefäße 
  vorgeschoben und anschließend ein
Röntgen-Kontrastmittel appliziert. Dadurch kann mittels
Durchleuchtungsgeräten der Blutfluß in den  Koronargefäßen
beobachtet und beurteilt werden, sowie Verengungen und
Verschlüsse identifiziert werden (diagnostische
Herzkatheteruntersuchung).
Wird eine solche Stelle gefunden, kann diese oft auch gleich mittels des
Herzkatheters behandelt werden
(Intervention; \I{PCI} (Percutane Coronare
Intervention) oder \I{PTCA} (Perkutane Transluminale
Coronare Angioplastie)). Dabei wird ein Ballon
in die Engstelle eingebracht, welcher das Gefäß aufdehnt.
Anschließend wird ein zylindrisches Metallgeflecht, ein
sog.
\I{Stent}, eingebracht, welcher als Stütze für das Gefäß
dient und dieses längerfristig offen halten soll. Vgl.
\cite{Harrisons18DeKap230,Harrisons18DeKap246}
}
\end{itemize}

}
{
\begin{itemize}
  \item Lyse
  \item Herzkatheterlabor (PCI, PTCA)%
  \opt{REGVIE}{%
\begin{itemize}
  \item Dienstbereites Katheterlabor: eigener Dienstplan
\end{itemize}%
}  	
\end{itemize}

}{}{}{}


\bigskip



\clearpage
\subsection{Herzrhythmusstörungen}
\label{topic:herzrhythmusstoerungen}
\index{Herzrhythmusstörungen}
\index{Rhythmusstörungen!Herz-}

Herzrhythmusstörungen sind häufig, allerdings auch oft sehr
unterschiedlich in ihrer Bedeutung: Manche sind chronisch
und eher wenig gefährlich, andere können plötzlich auftreten
und sogar akut lebensbedrohlich sein. Rhythmusstörungen
können sowohl selber das Problem darstellen, als auch \E{als
Komplikationen} im Rahmen anderer Krankheitsbilder (z.\,B.
Herzinfarkt, etc.) auftreten.

\TextSlide{Arten von Herzrhythmusstörungen}
{%
Herzrhythmusstörungen können hervorgerufen werden durch eine
\EE{Störung der Frequenz} (zu langsam, zu schnell) oder
einer \EE{Störung der Regelmäßigkeit}. Wenn man die
elektrischen Ströme mittels eines EKG darstellt, kann man
zusätzlich auch \EE{Störungen der Form der elektrischen
Herzströme} erkennen und beschreiben. Ein halbautomatischer
Defibrillator (AED) erledigt dies automatisch, er zeichnet
die elektrischen Herzströme auf und beurteilt diese.

Im Rahmen von Herzrhythmusstörungen kann es zu einer
\EE{Entkoppelung der elektrischen und der mechanischen
Herzaktion} kommen, z.\,B. wenn nicht jedem elektrischen
Impuls ein Herzschlag folgt.
}{
\begin{itemize}
  \item Störung der
  \begin{compactitem}
    \item Frequenz
    \item Regelmäßigkeit
    \item Form der elektrischen Herzströme
  \end{compactitem}
\end{itemize}

Entkoppelung von elektrischer und der mechanischer
Herzaktion möglich!
}{}{}{}

\SlideCompensationNoParskip{1}

\TextSlide{Hier zählt die Geschwindigkeit: Tachy- und Bradykardie}
{%
%
\begin{itemize}
  \item \T{Bradykardie}:  Zu langsamer Herzschlag
  ({HF \textless} 60\,\nicefrac{}{\Min} {\tiny (beim Erwachsenen)})
  
  Leistungssportler können aufgrund ihres trainierten Körpers auch sehr
  niedrige Ruhepulsfrequenzen haben.
  \item \T{Tachykardie}:  Zu schneller
  \marginpar{{\footnotesize Tachometer im Auto:
  Geschwindigkeitsanzeiger, meistens zu schnell ;)}}
  Herzschlag
  (\Abk{HF} {\textgreater} 100\,\nicefrac{}{\Min} {\tiny
  (beim Erwachsenen)})
\end{itemize}

Denke bei der Beurteilung auch immer an die Umstände und den
Erregungszustand des Patienten!
}{
\begin{itemize}
  \item Bradykardie \dotfill\ zu langsam
  \item Tachykardie \dotfill\ zu schnell
\end{itemize}
}{}{}{}


\begin{table}[H]
\FontTableBase
\begin{WidePar}
\caption{Übersicht: Wichtige und häufige
Herzrhythmusstörungen.\label{tab:herzrhythmusstoerungen}}

% Index-Einträge für Tabelle ---
\index{Bradykardie} 
\index{Tachykardie}
\index{Kammerflimmern}
\index{Flimmern!Kammer-}
\index{Vorhofflimmern}
\index{Flimmern!Vorhof-}
\index{Asystolie}
\index{Pulslose Ventrikuläre Tachykardie} 
\index{Tachylardie!Pulslose Ventrikuläre}
% --- Ende


\begin{minipage}{\linewidth}
\FontTableBase
\begin{tabulary}{\textwidth}{LLL}\addlinespace\toprule\addlinespace
\noindent\TabHeading{Rhythmus}
& 
\noindent\TabHeading{Bedeutung}
& 
\noindent\TabHeading{Bemerkung}
\\\addlinespace\midrule\addlinespace
\TabSub{Bradykardie}

{\scriptsize (HF\,{\textless}\,60\PerMin*)}
& 
Situations\-abhängig 
& 
\emph{\enquote{zu langsam}}
\\\addlinespace
\TabSub{Tachykardie} 

{\scriptsize (HF\,{\textgreater}\,100\PerMin*) }
& 
Situations\-abhängig 
& 
\emph{\enquote{zu schnell}}
\\\addlinespace\midrule\addlinespace
\TabSub{Kammerflimmern}
& Reanimation  & \emph{\enquote{schockbarer}} Rhythmus 
\\\addlinespace
\TabSub{Pulslose Ventrikuläre Tachy\-kardie} 
& Reanimation 
& \emph{\enquote{schockbarer}} Rhythmus 
\\\addlinespace
\TabSub{Asystolie} 

{\scriptsize (HF\,=\,0)} 
& 
Reanimation 
& 
Nulllinie, \emph{\enquote{nicht-schockbarer}} Rhythmus
\\\addlinespace
\TabSub{Vorhofflimmern}
& regellose Erregung im Vorhof

oft symptomlos 
& 
\E{Chronische Erkrankung.} Sehr viele Leute haben
Vorhofflimmern und leben ganz gut damit. 

Risikofaktor für:

\E{Tachykardie-Anfälle}

\E{Thrombenbildung}, dadurch erhöhtes Lungenembolie- und
Schlaganfallrisiko. Deshalb Vorbeugung mit
gerinnungs\-hemmenden (\enquote{blut\-ver\-dünnenden})
Medi\-ka\-menten:
z.\,B. \E{\enquote{Marcoumar}} ${\rightarrow}$ Patient ist
blutungsgefährdet
\\\addlinespace\midrule\addlinespace
\TabSub{Extrasystole}
&
Extraschläge
&
Vereinzelt bis \enquote{Salven} möglich

Von symptom-/bedeutungslos bis lebensbedrohlich
\\\addlinespace\bottomrule
\end{tabulary}
\end{minipage}
\end{WidePar}
\end{table}





\subsubsection{Besondere Rhythmen}





\TextSlide{Reanimationspflichtige Rhythmusstörungen}
{
\begin{description}
  \item[\T{Kammerflimmern}]
%   \index{Kammerflimmern|TLike}
  \label{topic:kammerflimmern}
  %
  NICHT verwechseln mit Vorhofflimmern! Die elektrische
  Erregung im Herzen ist ungerichtet, als würden alle Autos
  eines Parkplatzes gleichzeitig in alle Richtungen losfahren
  ${\rightarrow}$ der Herzmuskel kontrahiert nicht geordnet,
  er \enquote{zittert} nur (wie ein Pudding), es wird \EE{keine
  Auswurfleistung} erreicht!
  \item[\T{Pulslose Ventrikuläre Tachykardie}]
%   \index{Pulslose Ventrikuläre Tachykardie|TLike}
  \label{topic:pulslose-ventrikulaere-tachykardie}
  %
  Vorstufe zum Kammerflimmern. Das Herz kontrahiert
  so schnell, dass es sich nicht mehr füllen kann
  ${\rightarrow}$ \EE{keine Auswurfleistung}.
  \item[\T{Asystolie}] 
%   \index{Asystolie|TLike}
  \label{topic:asystolie}
  
  Im Herz entsteht keine elektrische
  Erregung, der Herzmuskel kontrahiert daher nicht und es erfolgt
  {keine Auswurfleistung}
\end{description}


}
{
\begin{compactitem}
  \item Kammerflimmern
  \begin{compactitem}  \item {Reanimation!}
    \item Wirre Erregung, keine sinnvolle Muskelarbeit ${\rightarrow}$ kein
    Auswurf.
  \end{compactitem}
  \item Pulslose Ventrikuläre Tachykardie
  \begin{compactitem}
    \item Keine Zeit zum Füllen ${\rightarrow}$ kein Auswurf
  \end{compactitem}
  \item Asystolie
  \begin{compactitem}
    \item Keine elektrische Erregung, keine Muskelarbeit, kein
    Auswurf
  \end{compactitem}
\end{compactitem}

}{}{}{}



\SlideCompensationNoParskip{1}
    
\TextSlide{Vorhofflimmern}
{
\index{Vorhofflimmern|TLike}
\label{topic:vorhofflimmern}
Regellose Erregung im \E{Vorhof}, die oft symptomlos ist. 
(Die Herzkammern arbeiten dabei normal!) Vorhofflimmern kann
ständig bestehen oder auch nur manchmal (episodenhaft)
auftreten. Sehr viele Leute haben Vorhofflimmern und leben
ganz gut damit.
Häufige \EE{Komplikationen} sind jedoch:
\begin{compactitem}
  \item Tachykardie-Anfälle
  \item Thrombenbildung: Dies stellt wiederum ein Risko dar für:
  \begin{compactitem}
    \item Arterielle Gefäßverschlüsse
    \item Schlaganfälle (\FullPrettyRef{topic:insult})
    \item Mesenterialinfarkte (\FullPrettyRef{topic:mesenterialinfarkt})
  \end{compactitem}
  Daher Vorbeugung mit \EE{gerinnungshemmenden}
  (\enquote*{blutverdünnenden}) Medikamenten (Marcoumar
  ${\rightarrow}$ Patient ist blutungsgefährdet)
\end{compactitem}
}
{%
\begin{compactitem}
  \item Flimmern im \EE{Vorhof}
  \item Häufige und oft chronische Erkrankung
  \item Thrombenbildung, Tachykardieanfälle
  \item Oft \E{gerinnngshemmende Medikamente}
\end{compactitem}

}{}{}{}

\SlideCompensationNoParskip{1}


\TextSlide{Extrasystolen}
{%
\index{Extrasystolen|TLike}
\label{topic:extrasystolen}
Extrasystolen sind \enquote{Extraschläge} des Herzens, sie können
vereinzelt oder gehäuft auftreten, eventuell auch
\enquote{salvenartig}. Je nach Häufigkeit haben Extrasystolen
unterschiedliche Auswirkungen auf den Kreislauf und können
symptom- bzw. bedeutungslos sein, oder auch Beschwerden
verursachen. Wenn die Extrasystolen die normale
Herzpumpfunktion massiv stören, können Sie auch
lebensbedrohlich sein.
}
{%
\begin{itemize}
  \item Extraschläge
  \item Vereinzelt od. gehäuft, evtl. \enquote{salvenartig}
  \item Unterschiedliche Auswirkungen auf Kreislauf: symptomlos bis
  lebensbedrohlich möglich
\end{itemize}

}{}{}{}


\opt{STATUS-EXPERIMENTAL}{
\begin{table}[b]
\begin{WidePar}
\Captionof{table}{Bilderserie: Herzrhythmus(störungen)}
\begin{tabularx}{\linewidth}{XXXX}
% \includegraphics[width=\linewidth]{./images/hirtler-aass/Herzschmerz.pdf}
\Captionof{figure}{Normaler Sinusrhythmus
\SourcePicture{}} 
&
% \includegraphics[width=\linewidth]{./images/wikipedia-cc-by/Heart_ant_wall_infarction-00800.jpg}
\Captionof{figure}{Tachykardie \SourcePicture{}} 
&
% \includegraphics[width=\linewidth]{./images/wikipedia-cc-by/Heart_coronary_artery_lesion-lq.jpg}
\Captionof{figure}{Bradykardie \SourcePicture{}}
&
% \includegraphics[width=\linewidth]{./images/wikipedia-cc-by/Hk_coronary_bionerd.jpg}
\Captionof{figure}{Extrasystolen \SourcePicture{}}
\\\addlinespace
% \includegraphics[width=\linewidth]{./images/hirtler-aass/Herzschmerz.pdf}
\Captionof{figure}{Kammerflimmern \SourcePicture{}} 
&
% \includegraphics[width=\linewidth]{./images/wikipedia-cc-by/Heart_ant_wall_infarction-00800.jpg}
\Captionof{figure}{Ventrikuläre Tachykardie \SourcePicture{}} 
&
% \includegraphics[width=\linewidth]{./images/wikipedia-cc-by/Heart_coronary_artery_lesion-lq.jpg}
\Captionof{figure}{Asystolie \SourcePicture{}}
&
% \includegraphics[width=\linewidth]{./images/wikipedia-cc-by/Hk_coronary_bionerd.jpg}
\Captionof{figure}{Pulslose Elektrische Aktivität \SourcePicture{}}
\\\bottomrule
\end{tabularx}
\end{WidePar}
\end{table}
} % STATUS-EXPERIMENTAL


\subsubsection{Tachykarde Attacke}
\label{topic:tachykarde-attacke}

\Definition{Plötzliche Tachykardie ohne erkennbare Ursache, wobei der Patient
(meistens) Symptome wahrnimmt und eine HF\,>\, 160\PerMin*
vorliegt.}


\TextSlide{\TxBeschreibung}
{%
Tachykarde Attacken sind ein häufiges Notfallbild. Da oft
chronische Erkrankungen die Ursache sind, haben die
Patienten diese Attacken immer wieder und kennen die
Symptome bereits (und oft auch die Behandlung). Die Ursachen
können sehr unterschiedlich sein und sind auch vom Fachmann
oft nicht leicht zu ermitteln.

Eine große Rolle spielt die \EE{Anamnese}: Oft sind
entsprechende Grunderkrankungen schon seit langem bekannt,
oft gibt es bereits entsprechende Arztbriefe oder Befunde.
Ein rasches Hinweisen auf diese Befunde ist für den
nachbehandelnden Arzt wichtig um eine rasche, zielführende
Behandlung einzuleiten.
}{
\begin{itemize}
  \item Oft Grunderkrankung bekannt
  \item Anamnese!
  \item Befunde, Arztbriefe!
\end{itemize}
}{}{}{}

\TextSlide{\TxSymptome}
{Es liegt klarerweise eine \EE{Tachykardie} vor. Der Puls
ist i.\,d.\,R. hoch, meist über 160\PerMin*. Der Patient
klagt oft über ein \EE{\enquote{Herzklopfen}} und ist unruhig. Manchmal wird der Anfall von
Angstgefühlen begleitet.

Oft kommen \EE{Begleitsymtome} dazu, wie zum Beispiel Atemnot oder ein leichter
Brustschmerz. Dies sind erste Zeichen, dass das Herz an
die Grenze seiner Belastbarkeit stößt.\bigskip
}{
\begin{itemize}
  \item Tachykardie
  \item Spürbares Herzklopfen
  \item Evtl. Atemnot, leichter Brustschmerz, Angstgefühl
\end{itemize}
}{}{}{}

\TextSlide{{\Lightning}Gefahr }
{%
Da das \EE{Herz plötzlich viel mehr Arbeit leisten muss},
kann es sein, dass es an seine \EE{Grenzen} stößt und
\EE{versagt} bzw. die Blutversorgung des Herzens für diese
Arbeit nicht mehr reicht und ein Angina-pectoris-Anfall oder
sogar ein Herzinfarkt\A{ KOCH: passt mit der Definition MCI
nicht überein, AP wäre besser oder? GABS: AP dazugenommen,
Infarkt aber gut möglich. } entsteht.

Weiters können aus einer tachykarden Rhythmusstörung in
andere, akut lebensbedrohende Rhythmusstörungen umschlagen.
\Commentary[Tachykarde Attacke / Gefahren] {Es gibt viele verschiedene Auslöser für tachykarde
Attacken. Manche Rhythmusstörungen sind tückisch, dass sie
von einem Moment umschlagen und einen Kreislaufstillstand
verursachen können. Zum Beispiel beim Kammerflattern besteht
oft eine \Conc{2}{1}-Überleitung vom Vorhof in die Kammern
mit einer resultierenden Frequenz um die 140\PerMin*. Kommt es
zu der gefürchteten \Conc{1}{1}-Überleitung wird die
Flatterfrequenz der Vorhöfe von \Range{250}{350}\PerMin* (!) direkt auf die
Kammern übergeleitet, dies kommt meist einem Herzstillstand
gleich. \cite{BLIM4}}

\bigskip
}{
\begin{itemize}
  \item Herz überfordert:
  \begin{compactitem}
    \item Angina pectoris
    \item Herzinfarkt
  \end{compactitem}
\end{itemize}
}{}{}{}






% \clearpage
\subsection{Mechanische Atemwegsverlegung}
\label{topic:mechanische-atemwegsverlegung}
\label{topic:fremdkoerperaspiration}

\HeadingSub{\Lang{Syn.} Fremdkörperaspiration}


\Layout{20}{\TxBeschreibung}
{%
\index{Aspiration}%
Wenn ein Fremdkörper den Atemweg teilweise oder vollständig 
\E{blockiert} oder \E{behindert}, spricht man von einer
mechanische Verlegung. Abhängig von der Größe des
Fremdkörpers können sowohl die oberen als auch die unteren
Atemwege betroffen sein. Wenn der Fremdkörper in die
\E{unteren Atemwege} gelangt, spricht man von einer
(Fremdkörper-) \T{Aspiration}.

Nicht nur feste Körper, sondern auch Flüssigkeiten können
Auslöser einer Verlegung oder Aspiration sein. Besonders
Blut, Magensaft, Erbrochenes oder Speisebrei stellen
diesbezüglich eine Gefahr dar.
Besonders gefährdet hinsichtlich einer Atemwegsverlegung sind
Kinder, ältere Personen, bewusstseinseingschränkte Personen,
Patienten mit Schluckstörungen oder Allergiker (z.\,B.
Anschwellen der Schleimhäute nach einem Bienenstich).
}
{%
\begin{itemize}
  \item Verlegung durch Fremdkörper
  \item Obere und untere Atemwege möglich
  \item Aspiration: untere Atemwege
  \item Besonders gefährdet:
  \begin{itemize}
    \item Kinder oder alte Personen
    \item Bewusstseinsgetrübte/-lose Personen
    \item Patienten mit Schluckstörungen
    \item Allergiker
  \end{itemize}
\end{itemize}
}
{./images/hirtler-aass/Fremdkoerper.pdf}
{Atemwegsverlegung (Schema). Der linke Bissen steckt in der Luftröhre fest, der
andere befindet sich in der Speiseröhre.}
{Hirtler.}


% \clearpage
\TextSlide{\TxSymptome}
{%
Die Symptome sind abhängig  von der Schwere der Verlegung.
Man unterscheidet zwischen einer \EE{schweren} und einer
\EE{milden Atemwegsverlegung}
\citep{Erc2005Sec2,Erc2005Sec6,Erc2005DeSec2,Erc2005DeSec6}.

\begin{itemize}
  \item \EE{Schwere Verlegung}:
  Bei einer schweren Verlegung der Atemwege ist der Patient
  \EE{unfähig zu reden oder effektiv zu husten}. Eventuell kann
  er nur durch Nic\-k\-en oder Gesten kommunizieren. In besonders
  schweren Fällen kann es zu einem \EE{Atemstopp} und nach
  kurzer Zeit zu einer \EE{Bewusstlosigkeit}, gefolgt von einem
  \EE{Kreislaufstillstand} kommen.
  \item \EE{Milde Verlegung}: Bei der milden Verlegung kann
  der Patient noch sprechen und atmen. Er klagt meist über
  Atemnot und ein Fremdkörpergefühl, oft versucht er zu husten
  oder zu würgen. Oft ist auch ein pfeifendes Atemgeräusch
  während der Ein- und/oder Austamung zu hören
  (in-/exspiratorischer \E{Stridor}).
\end{itemize}


% \bigskip
% \SlideCompensationNoParskip{1}
}{
\begin{itemize}
  \item Schwere Verlegung:
  \begin{itemize}
    \item Unfähigkeit zu Sprechen oder effektiv zu husten
    \item Evtl. Bewusstlosigkeit → Kreislaufstillstand
  \end{itemize}
  \item Milde Verlegung:
  \begin{itemize}
    \item \EE{Kann Sprechen \& Husten}.
    \item Atemnot
    \item Fremdkörpergefühl
    \item Pfeifendes Atemgeräusch (Stridor)
  \end{itemize}
\end{itemize}
} 
{} 
{} 
{}

\Fazit{Die (Un-) Fähigkeit zu sprechen und zu husten ist das
wesentliche Kriterium um die Schwere einer Atemwegsverlegung
einzuschätzen.}

\Text{Schläge zwischen die Schulterblätter}
{%
Durch kräftige Schläge auf den Rücken zwischen die
Schulterblätter kann  ein festsitzender Fremdkörper evtl.
gelockert und ausgehustet werden.}
{}{}{}{}


\Layout{20}{Heimlich-Manöver}
{%
\label{topic:heimlich-manoever}%
\Lang{Syn.} Heimlich-Handgriff. Beim Heimlich-Manöver
versucht man durch Kompression des Bauchraums den Fremdkörper
durch den entstandenen Überdruck aus den Atemwegen zu
entfernen.

Der Helfer umfasst von hinten den Oberbauch des Patienten und
bildet mit der einen Hand eine Faust und umschliesst diese
mit der anderen Hand. Nun legt er die Faust unterhalb der
Rippen und des Brustbeins an und zieht sie dann ruckartig
gerade nach hinten zu sich.

\Ca{Verletzungsgefahr!} Durch Anwendung des
Heimlich-Handgriffes kann es zu schweren inneren Verletzungen
kommen, die Anwendung ist daher auf bedrohliche Situationen
beschränkt.
}{%
\begin{itemize}
  \item Kompression des Bauchraumes → Überdruck
  \item Ansatzpunkt unterhalb des Rippenbogens, ruckartige Bewegung
  \item \Ca{Verletzungsgefahr!}
\end{itemize}
}
{./images/wikipedia-pd/Abdominal_thrusts3.jpg}
{Heimlich-Manöver.}
{U.\,S. Naval Hospital, Okinawa, Japanese intern Naoko Uehara teaches
proper technique for treating a choking victim with the assistance of Hospital
Corpsman Second Class Ross Francis. Amanda M. Woodhead, USMC / WmCo, PD.}


% {./images/wikipedia-pd/US_medic_teaches_the_Heimlich_manuever_to_laughing_Afghans_00500.jpg}
% {Heimlich-Manöver.}
% {\protect\enquote{Air Force Staff Sgt. Abraham Jara provides instruction in the proper
% application of the Heimlich maneuver to teachers at the Kharana Girls Middle
% School, Dara district, Panjshir province, Afghanistan.} Photo by 
% Capt. John T. Stamm, USAF / WmCo, PD.}






 \SlideCompensationNoParskip{2}


\begin{WidePar}
\MassnahmenWide{Spezielle Maßnahmen}{Mechanische Atemwegsverlegung}{}
{%
\begin{itemize}
  \item Atemwege soweit möglich freimachen
%   \item \TxBeiVit \TxMassVitBes
  \begin{itemize}
    \item 
    Lagerung: Situationsgerecht, z.\,B. Oberkörper erhöht
%     
    Bei Schlag
    zwischen die Schulterplätter soll der Patient abwärts schauen
    \item  \ce{O2}-Gabe hochdosiert (\VonBis{10}{15}\LPerMin*)
  \end{itemize}
  \item Maßnahmen abhängig von der Beurteilung der Schwere der
  Atemwegsverlegung\:
\end{itemize}

\noindent\begin{tabularx}{\linewidth}{XXX}
\toprule\addlinespace
\multicolumn{3}{c}{\TabSub{Beurteilung der Schwere der Atemwegsverlegung}}
\\\addlinespace\midrule\addlinespace
\TabSub{Schwer}

Kann nicht sprechen

Hustet ineffektiv
&
&
\TabSub{Mild}

Kann sprechen

Hustet effektiv
\\\addlinespace\cmidrule(lr){1-2}\cmidrule(lr){3-3}\addlinespace
\TabSub{Ohne Bewusstsein}
&
\TabSub{Bei Bewusstsein}
&
\TabSub{Bei Bewusstsein}
\\\addlinespace\cmidrule(lr){1-1}\cmidrule(lr){2-2}\cmidrule(lr){3-3}\addlinespace
5 Initialbeatmungen

Reanimation
&
5 Schläge zwischen Schulterblätter

5 Heimlich-Manöver (\textgreater{}\,1\,\MetricUnit{a}) oder Thoraxkompressionen
(\textless{}\,1\,\MetricUnit{a})

Wiederholung 
&
Husten lassen

Überwachung
\\\addlinespace\bottomrule\addlinespace
\end{tabularx}

Bei Schlägen
    zwischen die Schulterplätter soll der Patient
    nach Möglichkeit abwärts schauen (besonders relevant
    bei Kindern).

Durch Anwendung des Heimlich-Handgriffes kann es zu
schweren inneren Verletzungen kommen. Ein Patient ist
nach Anwendung unbedingt zu hospitalisieren!
\citep{Erc2005DeSec2}
}


\end{WidePar}

% \Text{\dsliterary}{%
\nocite{Wuttig2009}
% }{}{}{}{}

\clearpage
\subsection{Asthma bronchiale}
\label{topic:asthma-bronchiale}

\Definition{Chronische Erkrankung mit mehr oder weniger
  häufigen Episoden von Anfällen mit massiver Atemnot oder
  Atembehinderung durch Verengung der Bronchien.}

% \DefinitionExp{Asthma_bronchiale}

\TextSlide{\TxBeschreibung{}}
{
Im akuten Asthmaanfall kommt es zu einer Verengung der
Bronchien und vermehrter Schleimbildung, die \EE{Ausatmung}
wird massiv erschwert und der Patient hat das Gefühl nicht
genug Luft zu bekommen.
}{%
Wiederholte Anfälle von Atemnot durch

\begin{itemize}
  \item Verengte Bronchien
  \item Vermehrte Schleimbildung
  \item Behinderte Ausatmung
\end{itemize}
}
{}
{}
{}





Asthma ist bei den meisten Patienten oft schon bekannt und
kann bzw. muss erfragt werden. Viele Patienten haben auch
einen oder mehrere Sprays oder Inhalationsgeräte zur Dauer-
oder Akutbehandlung. 
Besonders betroffen vom Asthma bronchiale sind Patienten
mit bekannten Allergien. Folglich kann alles, was bei dem
Patienten eine Allergie verursacht auch einen Asthmaanfall
auslösen (Pollen, Gräser, Katzen, Zigarettenrauch, \ldots). 

Ist bei Eintreffen des Rettungsdienstes noch keine
Besserung der Atemnot eingetreten, ist diese Atemnot als
massiv und lebensbedrohend einzuschätzen. Es ist nicht
damit zu rechnen, dass der Zustand von alleine wieder
vergeht und bedarf einer medikamentösen Behandlung durch einen
Notarzt. 

\Z{Unterschätze nie die Schwere eines Asthmaanfalles!
Gefährlich sind besonders die \protect\enquote*{ruhigen},
wenig klagenden Patienten, diese haben oft ein anderes
Luftnot-Empfinden. Diese Patienten sind Hochrisiko-Patienten
und haben auch mit der \ce{O2}-Gabe Probleme! \ce{O2}-Gabe nur
bis max. 95\PC* Sauerstofsättigung!
\citep{Olschewski2009}\A{\footnote{Vortrag Univ-Prof.Dr. Horst Olschewski, 5. AGN-Kongress 2009; Notfall Rettungsmed (2009) 12:322}}}


\TextSlide{\TxSymptome}
{%
Das Leitsymptom ist die Atemnot. Oft ist ein pfeifendes,
brummendes oder giemendes Atemgeräusch, besonders während
der Ausatmung (\E{Exspiration}, \E{exspiratorischer
Stridor}) wahrnehmbar. Die Exspiration ist ausserdem
verlängert. Trockener Husten ist häufig.
Der Patient nimmt meist automatisch und unbewusst eine
Position ein, die seine Atemhilfsmuskulatur unterstützt: Er
sitzt aufrecht und stützt sich oft nach vorne oder hinten
ab.
Die Atemfrequenz ist erhöht, oft auch die Herzfrequenz
(\E{Tachypnoe} und \E{Tachykardie}). Die Sauerstoffsättigung
kann erniedrigt sein.

\subparagraph{\Ca{Schwerer Anfall}}

Im schweren Anfall ist oft \EE{kein}, oder nur mehr ein
\EE{leises Atemgeräusch} hörbar (da kaum mehr Luft bewegt
wird, \Z{\protect\enquote{silent chest}}). Es kommt zu einer
Kreislaufdepression, Herzfrequenz und Blutdruck stürzen ab.
Dies kann zu Bewusstseinsstörungen führen. Die
Sauerstoffsättigung ist stark vermindert und der Patient ist
zyanotisch.


}{
\begin{itemize}
  \item Akute Atemnot mit trockenem Husten
  \item Exspiratorisches Atemgeräusch (Stridor)
  \item Einsatz der Atemhilfsmuskulatur
  \item Verlängerte Ausatmungsphase
  \item HF, AF $\uparrow$
  \item Schwerer Anfall:
  \begin{itemize}
    \item Kein oder leises Atemgeräusch
    \item HF, RR $\downarrow$
    \item Sp\ce{O2} $\downarrow\downarrow$
    \item Bewusstseinsstörungen
  \end{itemize}
\end{itemize}
}
{}
{}
{}

% \Text{\TxKomplikationen}
% {
% \begin{itemize}
%   \item Zyanose
%   \item Status asthmaticus
%   \item Panik
% \end{itemize}
% }{
% 
% }
% {}
% {}
% {}


\TextSlide{Status Asthmaticus}{%
\index{Status!asthmaticus}
\label{topic:status-asthmaticus}
Der
\protect\enquote{Status} ist eine gefürchtete Komplikation des
Asthmaanfalles bei der die Atemnot nicht wieder vergeht. Es
besteht die Gefahr, dass der Patient nach einiger Zeit keine
Kraft mehr zum Atmen hat (Erschöpfung). Der Status
asthmaticus kann tödlich enden.
}{
\begin{itemize}
  \item Lebensgefährliche Komplikation:
  \item Anfall hört nicht mehr auf
\end{itemize}
}{}{}{}

\SlideCompensationNoParskip{2}
\Layout{22}{}
{%
Patientin mit einem akuten Asthmaanfall.

Sie sitzt auf einer Treppe und stützt sich nach hinten mit
den Armen ab. Die Erstmaßnahmen bei vital bedrohten
Patienten wurden ergriffen:
\begin{compactitem}
  \item Situationsrechte Lagerung
  \item Sauerstoffgabe
  \item Reanimationsbereitschaft\footnote{Defibrillator eingeschaltet,
  sonstiges Zubehör nicht am Foto sichtbar}
  \item Notarztnachforderung\footnote{Paramedic
  vor Ort, am Foto nicht sichtbar.}
  \item Vitalwerte erheben und Monitoring:
  Blutdruck wird gerade gemessen, die Patientin
  ist am Monitor angeschlossen
\end{compactitem}
}{

}
{./images/asthma-paramedic.jpg}
{Patientin mit einem akuten Asthmaanfall.\label{fig:asthma-paramedic}}
{Gabriel}


\subsection{Schlaganfall, Apoplektischer Insult, Hirninfarkt}
\index{Schlaganfall}
\index{Insult}
\index{Apoplexie}
\index{Apoplektischer Insult}
\label{topic:tia}
\label{topic:insult}
\label{topic:schlaganfall}

\HeadingSub{\Lang{Syn.} Apoplektischer Insult, Apoplexie}

% \Definition{Siehe Arten.}

\TextSlide{Arten}
{%
Grundsätzlich gibt es 2 Formen, die aber vor Ort nicht
unterschieden werden können:
\begin{itemize}
  \item \EE{\T{Trockener Insult}}:
  \I[Ischämie!trockener Insult]{Ischämie}. Verschluss eines
  Hirngefäßes mit anschließendem Absterben der versorgten Hirnregion
  \item \EE{\T{Blutiger Insult}}: \I{Hirnblutung}, z.\,B.
  durch Platzen eines Aneurysmas eines Hirngefäßes, kann spontan
  oder aufgrund eines Traumas auftreten.
\end{itemize}

}
{
\begin{itemize}
  \item \enquote{Trockener} Insult
  \item \enquote{Blutiger} Insult
\end{itemize}

}{}{}{}

\TextSlide{\T{TIA}}
{%
\index{TIA}
\index{Transitorisch-Ischämische Attacke}
\label{topic:tia}
TIA ist die Abkürzung für Transitorisch-Ischämische Attacke.
Von einer TIA spricht man, wenn die Symptome  innerhalb von
24 Stunden vergehen. Solange noch Symptome bestehen, kann
man nicht von einer TIA sprechen, der Patient ist wie bei
jedem anderen Schlaganfall zu behandeln.
\bigskip
}
{
% TODO MED: Slide fehlt!
\begin{itemize}
  \item Transitorisch-Ischämische Attacke.
  \item Symptome <\,24\Hour*.
\end{itemize}

}{}{}{}


\TextSlide{\TxSymptome}
{%
\label{topic:hirndruckzeichen}
Die diversen \EE{Funktionsareale}
(Sprachzentren, Sehzentrum, motorische Zentren, \ldots)
sind in verschiedenen Regionen des Hirns verteilt. Je
nachdem welches Gebiet durch die Ischämie oder Blutung
betroffen ist, können die Symptome sehr unterschiedlich
sein.  Typische Störungen sind:

\begin{itemize}
  \item \EE{Motorik}: Störungen der Kraft und Beweglichkeit
  sind sehr häufig.
  Charakteristisch für einen
  Schlaganfall sind sog. \T{Halbseitenzeichen}: Es kommt
  dabei zu einer \E{einseitigen Kraftminderung}, dies kann
  bis zu einer \E{schlaffen, kompletten Lähmung} führen. Ein
  \E{herabhängender Mundwinkel} ist ein Zeichen einer
  Lähmung der
  Gesichtsmuskulatur. Durch die Lähmungen kann es zu
  \EE{Schluckstörungen} kommen, es besteht dann eine erhöhte
  Aspirationsgefahr.
  
  Der \T{Herdblick} ist eine Blickdiagnose: Es kann zu
  Blicklähmungen kommen, der Patient blickt in Richtung des
  \enquote*{Herdes des Geschehens}.
  
  \item \EE{Sprache}: Störungen der Sprache sind ebenfalls
  sehr häufig. Oft nimmt man eine \EE{verwaschene Sprache}
  wahr, welche (fälschlich) an Trunkenheit erinnert.
  \Z{Weiters kann es zu einer \enquote{wirren Sprache}
  kommen: Worte fehlen dabei oder Silben werden vertauscht.
  Im Extremfall kann die Sprache gänzlich \enquote{verloren}
  gehen. Auf der anderen Seite kann die Sprache zwar
  erhalten, aber das Sprach\E{verständnis} geschädigt sein.}
  
  \item \EE{Wahrnehmung}:
  Es kann zu Wahrnehmungsstörungen wie z.\,B.
  \EE{Sehstörungen}, plötzliche Erblindung, Doppelbilder,
  Gesichtsfeldausfälle kommen.
  
  \item \EE{Bewusstsein}: Störungen des Bewusstseins sind
  gefährlich und müssen bei Insult-Patienten immer erwartet
  werden. Der Bewusstseinszustand kann sich auch plötzlich
  verschlechtern.
  
  \item \EE{Allgemein}: Weiters können unspezifische
  Symptome wie z.\,B. Schwindel und Kopfschmerzen auftreten. Ein
  plötzlicher, \EE{donnerschlagartiger Kopfschmerz} ist ein
  klassisches Symptom für eine spontane Hirnblutung.
  
  \item \EE{\T{Hirndruckzeichen}}: Bei einer Erhöhung des
  Hirndrucks, z.\,B. aufgrund einer Hirnblutung, kommt es zu
  einer typischen Kombination von Symptomen: 
  
  \parpic[4cm][r]{\includegraphics[width=4cm]{./images/trauma/pupillendifferenz.jpg}}
  Meist das erste Symptom sind \EE{Kopfschmerzen}, gefolgt
  von Übelkeit. In weiterer Folge treten
  \EE{Bewusstseinsstörungen} auf. Der Patient ist
  \EE{hyperton, aber bradykard} (\EE{RR\,\Sign{↑},
  HF\,\Sign{↓}}), der Puls ist langsam, aber äußerst stark
  spürbar (\I{Druckpuls}). Die \EE{Atmung} kann \EE{unregelmäßig} sein. Weiters
  ist die \EE{Ungleichheit der Pupillen} \Z{(Anisokorie)\index{Anisokorie|Z}}, bzw.
  eine verlangsamte Lichtreaktion
  typisch.\ZFootnote{\E{Erhöhter Hirndruck}:In schweren
  Fällen kommt es zu sog. Strecksynergismen, d.\,h.  der
  Patient reagiert bei Schmerzreiz nur mit einem
  ungerichtetem Strecken der Extremitäten.}
  \nocite{StriebelAnaesthesie2} %[Bd.2, S.1344]
  
  \Z{Hirndruckzeichen werden nicht nur durch
  Hirnblutungen verursacht, auch z.\,B. beim \Z{Hirnödem}
  oder Schädel-Hirn-Trauma (\Abk{SHT}) können sie
  vorliegen.}
\end{itemize}

Präklinisch kann man nicht sicher zwischen einer Ischämie
(trockener Insult) und einer Blutung unterscheiden.
Auf jeden Fall ist die \EE{Bewusstseinslage} sehr engmaschig
zu überwachen; eine rasche Verschlechterung ist jederzeit
möglich! 
}
{
Je nach betroffener Hirnregion sehr unterschiedlich.
\begin{itemize}
  \item Halbseitenzeichen (Schwäche, Lähmung)
  \item Herdblick
  \item Verwaschene Sprache
  \item Sehstörungen
  \item Bewusstseinsstörungen
  \item Kopfschmerzen
  \item Hirndruckzeichen
  \begin{itemize}
    \item Bewusstseinsstörungen
    \item RR\,\Sign{↑} \& HF\,\Sign{↓}
    \item Ungleichweite Pupillen, verlangsamte Lichtreaktion
    \item Strecksynergismen
  \end{itemize}
\end{itemize}
}{}{}{}


\TextSlide{Differentialdiagnose}
{
Die wichtigste Differentialdiagnose ist eine Störung des
\EE{Blutzucker}-Stoffwechsels. Sowohl die Hyper-, als auch
die Hypoglykämie können neurologische Symptome auslösen, daher
muss der Blutzucker immer gemessen werden! Ebenso ist immer
an \EE{Vergiftungen} zu denken.

Die Abgrenzung zu anderen neurologischen Erkrankungen ist
ohne weitergehende Untersuchungen oft schwierig bzw. schwer
möglich. 
}{
\begin{itemize}
  \item Blutzucker ${\downarrow}$, ${\uparrow}$
  \item Sonstige neurologische Erkrankung.
  \item Vergiftungen
\end{itemize}

}{}{}{}


% \clearpage
\subsection{Diabetes Mellitus}
\label{topic:diabetes}
\label{topic:diabetes-mellitus}

\Definition{Der \T{Diabetes mellitus} ist eine chronische
Stoffwechselerkrankung, bei der die Blutzuckerregulation gestört ist.} 



\Text{Exkurs: Der Blutzuckerwert}
{%
%
%\vspace{1em}
%
\begin{center}
\FontTableBase
\begin{tabular}{cccl}
\toprule
\multicolumn{3}{c}{\TabHeading{Wert [{\MgDl}]}}
&
\TabHeading{Beurteilung (nüchtern)}
\\\midrule
0 & -- & 20 & Lebensgefahr
\\\addlinespace
21 & -- & 60 & stark erniedrigt, Bewusstseinsstörungen
\\\addlinespace
61 & -- & 80 & erniedrigt, leichte Symptome
\\\addlinespace
81 & -- & 100 & Nüchtern-Blutzucker
\\\addlinespace
101 & -- & 200 & erhöht (nüchtern)
\\\addlinespace
201 & -- & 399 & stark erhöht
\\\addlinespace
400 & -- & ?? & sehr stark erhöht
\\\addlinespace\bottomrule
\end{tabular}
\end{center}

Durch die Erkrankung oder die unsachgemäße Behandlung kann es zu 
folgenden Notfällen kommen:

\begin{itemize}
  \item \EE{{Hyp\underline{\E{o}}glykämie}}: \enquote{Unterzucker}
  \item \EE{{Hyp\underline{\E{er}}glykämie}}: \enquote{Überzucker}
\end{itemize}
}
{}
{./icons/under-construction.png}
{Interpretation des
BZ-Wertes: Richtwerte.}
{}

\Text{Insulin}
{%
\T{Insulin} ist ein Hormon und wird in den
\enquote{Beta}-Zellen der Bauchspeicheldrüse (Pankreas)
produziert. Es ist für den Zuckerstoffwechsel wichtig. Es
sorgt dafür, dass Zucker (\T{Glukose}) aus dem Blut in die
Zellen aufgenommen wird und senkt somit den
Blutzuckerspiegel.
Wenn man künstlich Insulin spritzt, besteht die Gefahr eines
Unterzuckers (zu viel Insulin gespritzt oder zu wenig
gegessen ${\rightarrow}$  Spritz-/ Essabstand)
er\-fra\-gen.

\Fazit{Insulin senkt den Blutzuckerspiegel.} 

Hinsichtlich des Insulins sind zwei
Störungen häufig:
\begin{enumerate}
  \item Es wird zu wenig Insulin gebildet.
  \item Der Körper hat einen gesteigerten Bedarf an Insulin,
  weil er schon zu sehr daran gewöhnt ist. Der Körper
  reagiert auf die \enquote{normale} Dosis nicht mehr so wie er
  sollte, er braucht größere Mengen Insulin um den gleichen
  Effekt wie ein Gesunder zu erreichen (Gewöhnungseffekt),
  irgendwann kann aber der Körper nicht mehr Insulin
  produzieren.
\end{enumerate}
Dementsprechend unterscheidet man 2 Typen des Diabetes
mellitus.
}{

}{}{}{}







\TextSlide{Diabetes mellitus: Typen}
{
\begin{itemize}
  \item \EE{Typ I}:  \Z{(IDDM,
  insuline dependent DM)} 
  
  Die Insulinproduktion ist gestört
  und der Körper produziert zu wenig Insulin. Der Patient
  muss als Ersatz künstliches Insulin spritzen.  Die Patienten
  sind eher jung (Jugendliche, junge Erwachsene) und
  sportlich.
  
  \Z{Ursache für die Störung ist eine Störung des
  Immunsystems, bei der es zu einer Zerstörung der
  Insulin-produzierenden ${\beta}$-Zellen des Pankreas
  kommt. Durch den Produktionsausfall kommt es zu einem
  absoluten Mangel.}
  
  \item \EE{Typ II}:
  \Z{(NIDDM, non insuline dependent DM)} 
  
  Durch eine
  dauerhaft zu hohe Zuckerzufuhr kommt es im Körper zu
  Gewöhnungseffekten und einem Wirkungsverlust des vom
  Körper produzierten Insulins \Z{(Insulinresistenz)}. Als
  Folge muss der Körper immer mehr Insulin produzieren um
  den gleichen Effekt zu erreichen. Wenn das
  Produktionslimit erreicht ist, kommt es zu einer Erhöhung
  des Blutzuckerspiegels und zum Typ-II-Diabetes.
  
  Der Patient mit Typ-II-Diabetes benötigt oft kein Insulin.
  Die Erkrankung kann zuerst mittels Lebensstiländerungen
  und \EE{Medikamenten} behandeln. Erst im Endstadium muss
  mitunter auch Insulin gespritzt werden.
  
  Der Typ-2-Diabetes betrifft eher Patienten im mittleren
  Alter und aufwärts, aufgrund der zunehmend fett- und
  kohlehydratreichen Ernährungsgewohnheiten werden die
  Patienten aber immer jünger.
\end{itemize}
}{
\begin{itemize}
  \item Typ I
  \begin{itemize}
    \item Insulinproduktion gestört
    \item Insulin muss gespritzt werden
    \item Eher Jugendliche und junge Erwachsene
  \end{itemize}
  \item Typ II
  \begin{itemize}
    \item Lebensstil, Ernährung
    \item Insulin-Gewöhnungseffekt
    \item Eher ab mittlerem Alter
    \item Lebensstiländerung, Medikamente, Insulin
  \end{itemize}
\end{itemize}
}{}{}{}


\TextSlide{Spätfolgen}
{%
Ein hoher Blutzuckerspiegel über Jahre schädigt den gesamten
Körper. Es kommt zu:
\begin{itemize}
  \item \EE{Gefäßschäden}: 
  \begin{itemize}
    \item Gefäßverschlüsse (kleine und große Gefäße)
    \item \EE{Beine}: \enquote{Diabetischer Fuß}, infolge der
    Mangeldurchblutung in der Peripherie heilen auch kleine
    Wunden nicht mehr, in weiterer Folge muss immer weiter
    amputiert werden.
    \item \EE{Herz}: Die Herzkranzgefäße (Koronar\-arterien)
    sind ebenso betroffen. Diabetiker sind Risikopatienten
    für Herz-Kreislauf-Erkrankungen wie zB Herz\-infarkt.
    \item \EE{Hirn}: Wie beim Herzen kommt es zu
    Gefäßschäden, daher ist das Risiko für
    Schlaganfälle erhöht.
    \item \EE{Niere}: Die sehr kleinen Gefäße der Niere
    werden geschädigt. Im weiteren Verlauf stellt die Niere
    ihre Arbeit ein, es kommt zu einer
    \enquote{\T{chronischen Niereninsuffizienz}}, die im
    Endstadium mittels Blutwäsche (Dialyse) behandelt werden
    muss. Dialyse-Patienten machen einen großen Teil der
    Krankentransporte aus.
  \end{itemize}
  \item  \EE{Nervenschäden} \Z{(periphere Polyneuropathie)}:
  Gefühlslosigkeit vor allem in den Extremitäten, aber auch am
  gesamten Körper. Krankheiten, die sonst schmerzhaft sind,
  können beim Diabetiker aufgrund der Nervenschäden ohne
  Schmerzen ablaufen, z.\,B. ein schmerzloser, \EE{stummer Herzinfarkt}.
  \item  \EE{Augenschäden}: Bis hin zur Erblindung
  
\end{itemize}

\Cave{Achtung: \enquote{Stummer}, schmerzarmer/-loser Herzinfarkt beim Diabetiker!}
}{
\begin{itemize}
  \item Gefäßschäden (Herz, Hirn, Nieren, Extremitäten, \ldots)
  \item Nervenschäden
  \item \Ca{Cave \enquote{Stummer Infarkt}!}
  \item Augenschäden
\end{itemize}
}{}{}{}








\subsubsection{Hypoglykämie}
\label{topic:hypoglykaemie}
\index{Hypoglykämie}


\Definition{Maßgebliche Erniedrigung des Blutzuckerspiegels.}

\TextSlide{Ursachen}
{%
Die häufigsten Ursachen für eine plötzliche Hypoglykämie
sind Diätfehler und/oder eine falsche Anwendung von Insulin
oder bestimmten anderen Diabetes-Medikamenten. 

Aufgrund
besonderer Anstrengungen kann es zu einem gesteigerten
Zuckerbedarf des Körper kommen, wodurch der
Blutzuckerspiegel sinkt. Dies kommt häufig im Rahmen von
Begleiterkrankungen, wie z.\,B. grippalen Infekten, vor.
Auch die Alkohol-Intoxikation kann zu einer Hypoglykämie
führen\opt{NIVDOC}{ (Das beim Ethanolabbau entstehende
\Abk{NADH} hemmt die
Glukoneogenese.)} \citep{LoefflerBiochemie7}. 
\bigskip
}{%
\begin{itemize}
  \item Zu viel Insulin (spritzen ohne zu essen)
  \item Fasten (spritzen ohne zu Essen)
  \item Besondere Anstrengungen, Begleiterkrankungen,
  insb. Infekte
  \item Alkohol-Intoxikation
\end{itemize}
}{}{}{}






\TextSlide{\TxSymptome}
{%
Der Patient mit Hypoglykämie ist meistens unruhig,
desorientiert und agitiert, oft auch unkooperativ und
aggressiv. Oft fällt ein Zittern und eine schweißige Haut
auf. Auch Kopfschmerzen sind häufig. Eine Tachykardie und
Tachypnoe ist häufig feststellbar. Manchmal klagt der
Patient über Übelkeit und Heißhunger.
Je tiefer der Blutzuckerspiegel sinkt, desto deutlicher
werden die \EE{Bewusstseinsstörungen}. Von der anfänglichen
Agitiertheit/Desorientiertheit kann es rasch zu einer
Somnolenz, zum Sopor und zum Koma kommen!

Eine Hypoglykämie kann fast alle neurologischen  Störungen
aus\-lösen/imitieren! Häufig sind Kopfschmerzen und
Bewusstseinsstörungen, evtl. auch Krampfanfälle. Bei jeder
neurologischen Symptomatik muss man daher an eine
Hypoglykämie denken!

\Cave{Bei der Hypoglykämie können Symptome auftreten, die
den Patient wie alkoholisiert wirken lassen! \E{Hüte dich
vor vorschnellen
Diagnosen!}}

\Fazit{${\rightarrow}$ BZ-Messung bei jeder neurologischen Störung. }

}{
\begin{itemize}
  \item Unruhe, Zittern, Schweißausbruch
  \item Aggressives Verhalten, unkooperativ
  \item Neurologische Symptomatik
  \begin{itemize}
    \item Bewusstseinsstörungen, Verwirrtheit, Kopfschmerzen, Krampfanfälle,
    Bewusstlosigkeit
  \end{itemize}
  \item Heißhunger, Übelkeit.
  \item Tachykardie, Tachypnoe
  \item Haut feucht/blass
\end{itemize}
}{}{}{}












% \MassnahmenNG{Spezielle Maßnahmen}{Hypoglykämie}{}{
% \begin{itemize}
%   \item Vitale Bedrohung einschätzen. \TxBeiVit \TxMassVit
%   
%   \item BZ messen! Zur Diagnosestellung bzw. \E{vor} und \E{nach}
%   jeder Zuckergabe
%   \item Zucker zuführen: Wenn der Patient bewusstseinsklar
%   ist, Traubenzucker oder Zuckerwasser geben. Danach \enquote{feste
%   Speisen} nachessen lassen, da der Traubenzucker nur eine
%   kurze Wirkdauer hat.
%   
%   Wenn der Patient nicht mehr bewusstseinsklar ist, kann die
%   Zuckergabe nur über eine venöse Gabe erfolgen (getrübte
%   Patienten dürfen aufgrund der Aspirationsgefahr nichts
%   mehr essen).
%   
% \end{itemize}
% }


% \clearpage % TODO temp clearpage
\subsubsection{Hyperglykämie}
\label{topic:hyperglykaemie}
\index{Hyperglykämie}



\Definition{Maßgebliche Erhöhung des Blutzuckerspiegels.}

\Text{Hyperglykämie: Zustand oder Notfall?}
{%
Eine Erhöhung des Blutzuckerspiegels ist nicht automatisch
ein Notfall. Die Bedrohlichkeit ist z. B. stark davon
abhängig, wie rasch sich diese Erhöhung entwickelt hat. So
hat ein Patient, bei dem seit Jahren ein unerkannter
Diabetes mellitus besteht auch mit sehr hohen
Blutzuckerwerten (akut) keine Probleme. Der Körper hat sich
über die Zeit daran \E{gewöhnt}.

Ein anderer Patient, bei dem der Zustand plötzlich
eingetreten ist (z.\,B. Erstmanifestation eines Diabetes
mellitus Typ I bei einem 19-jährigen, oder wenn das Insulin
nicht gespritzt wurde, \ldots), kann mit dem gleichen
Blutzuckerwert bereits im Koma liegen.

\Fazit{Beurteile nicht nur den Messwert, sondern immer auch den Zustand des
Patienten.} }
{←}
{}{}{}



\TextSlide{Auswirkungen}
{
Im Grunde bereitet der Zucker im Akutfall dem Körper zwei
Probleme. Einerseits wirkt der Zucker wie ein Elektrolyt auf
den \EE{Flüssigkeitshaushalt} ein und stört diesen. Es kommt
zu einer ungünstigen Wasserumverteilung. Andererseits
mangelt es den Körperzellen an Zucker wenn das Insulin fehlt
um es in die Zellen hineinzubekommen, es entsteht ein
\EE{Nährstoffmangel}.
}{
\begin{itemize}
  \item Störung des Wasser- und Elektrolythaushaltes
  \begin{itemize}
    \item Zucker wirkt wie ein Elektrolyt
  \end{itemize}
  \item Störung der Nährstoffversorgung
\end{itemize}
}{}{}{}

\SlideCompensation{1}

\TextSlide{Ursachen}
{%
Die Hyperglykämie betrifft praktisch nur Diabetes mellitus-Patienten. Bei
Patienten, die einen noch unerkannten DM~Typ~I haben, ist
ein hyperglykämischer Notfall oft die Erstmanifestation (die
Diagnose Diabetes ist noch nicht bekannt!). Sonst sind
Diätfehler oder eine mangelnde Insulinzufuhr häufige
Ursachen. Auch Stress und Begleiterkrankungen wie z.\,B.
akute Infektionen kommen in Betracht, da hier der
Insulinbedarf steigt.
}{
\begin{itemize}
  \item Erstmanifestation eines DM~Typ~I
  \item Diätfehler
  \item Insulinzufuhr $\downarrow$
  \item Stress
  \item Infektionen (Insulinbedarf ${\uparrow}$)
\end{itemize}
}{}{}{}









\TextSlide{\TxSymptome}
{%
Der Patient hat i.\,d.\,R. ein erhöhtes Durstgefühl und
gleichzeitig eine \EE{vermehrte Harnproduktion} (Zucker wird
ausgeschieden). Er klagt über Abgeschlagenheit, in weiterer
Folge kann es zu \EE{Bewusstseinsstörungen} bis hin zum
\EE{Koma} (s.u.) kommen. Die Geschwindigkeit der
Veschlechterung ist je nach Patient sehr unterschiedlich.

\subparagraph{Koma} 
\label{topic:koma-diabetisches}
\label{topic:diabetisches-koma}
\Z{Wie oben bereits beschrieben wurde, kann es sowohl zu
einem Problem des Flüssigkeitshaushaltes und/oder der
Nährstoffversorung der Zellen kommen. Je nachdem welches
Problem im Vordergrund steht, kommt es zu \E{unterschiedlichen Symptomen}.}


Ist vorwiegend der Flüssigkeitshaushalt gestört kommt es zum
\Z{sog. hyperosmolaren} Koma, dessen Leitsymptom
erwartungsgemäß die Bewusstseinsstörung ist.
Steht der Nährstoffmangel im Vordergrund passiert etwas
Interessantes: Der Körper produziert sog. \Z{Ketonkörper}. Da
diese \E{sauer} sind, kommt es zu einer stoffwechselbedingten
\EE{Übersäuerung} (metabolische
\EE{Azidose}\ZFootnote{Genauer:
\I[Ketoazidose!diabetische]{diabetischen Ketoazidose}}). 
Er \EE{hyperventiliert}, um \ce{CO2} abzuatmen und somit
den pH wieder zu normalisieren (Puffersystem). Es ist daher
bei diesen Patienten eine \EE{schnelle, tiefe Atmung}
feststellbar (\Z{\I{Kußmaul-Atmung}}).

Mitunter kann auch ein blumig/fruchtiger oder Azeton-artiger
Atemgeruch wahrgenommen werden. Dieser darf nicht mit einer
\enquote{Alkoholfahne} verwechselt werden!
}{
\begin{itemize}
  \item Durstgefühl
  \item Vermehrte Harnproduktion
  \item Abgeschlagenheit, Bewusstseinstrübung bis
  \item Bewusstlosigkeit (Koma)
  \begin{itemize}
    \item Problem Flüssigkeitshaushalt
    \begin{itemize}
      \item \Z{Hyperosmolares Koma}
    \end{itemize}
    \item Problem Nährstoffmangel
    \begin{itemize}
      \item \Z{Ketoazidotisches Koma}
      \item Azidose → Hyperventilation zur Kompensation
    \end{itemize}
  \end{itemize}
  \item Azeton-Geruch (Nicht mit Alkoholfahne verwechseln!)
\end{itemize}
}{}{}{}








\subsection{Abdominelle Erkrankungen}



\TextSlide{Anamnese und Untersuchung}{%
Es ist daher wichtig
typische Symptome zu erheben und Alarmzeichen sofort zu
erkennen.
Wesentlich ist bei allen unklaren Bauchbeschwerden die
Erhebung einer exakten \I{Stuhlanamnese}:

\begin{itemize}
  \item Auffälligkeiten? (Geruch, Farbe, Konsistenz, Frequenz)
  \item Blutbeimengungen
  \item Wann erfolgte der letzte Stuhlgang?
  \item Winde?
  \item Genaue Beschreibung von Harn und eventuell Erbrochenem.
  \item Bei Frauen im gebärfähigem Alter Regelanamnese.
\end{itemize}

Der Patient muss bis zur endgültigen Klärung der Situation
\EE{nüchtern} belassen werden!
Wenn es der Patient toleriert, muss eine Betrachtung und
Abtastung des \so{unbekleideten} Abdomens beim liegenden
Patienten erfolgen.
}{%
\begin{itemize}
  \item Stuhl
  \begin{itemize}
    \item Wann letzter?
    \item Auffälligkeiten
    \item Farbe, Konsistenz
  \end{itemize}
  \item Erbrochen?
  \item Blut?
  \item Regel?
\end{itemize}
}{}{}{}


\TextSlide{Der Bauch -- unendliche Weiten}
{%
Gleich vorweg: In vielen
Fällen wird der Grund des Bauchschmerzes nicht mit
hinreichender Sicherheit feststellbar sein und es bietet
sich keine sichere Verdachtsdiagnose an. Statt einer
Diagnose wird dann das \EE{Beschwerdebild beschrieben} (Ort
des Schmerzes, Art des Schmerzes, Dauer, Verlauf,
Austrahlung, Ergebnis der abdominellen Tastuntersuchung:
Abdomen weich oder hart,
Druckschmerz, Loslassschmerz).
}{
\begin{itemize}
  \item In vielen Fällen Krankheitsbild unklar
  \item Im Zweifel: Beschreibung statt Diagnose
\end{itemize}
}{}{}{}



\subsubsection{Das Akute Abdomen}
\label{topic:akutes-abdomen}
\index{Abdomen!akutes|see{Akutes Abdomen}}


\Definition{\T{Akutes Abdomen}: Akute
Erkrankung des Bauchraumes mit plötzlichen, starken
\EE{Bauchschmerzen} und
\EE{bretthartem Bauch}, welche im weiteren Verlauf
lebensbedrohlich ist. Es sind viele Ursachen möglich!}

\Commentary[Akutes Abdomen, Definition]{Die Definition des
Akuten Abdomens ist oft nicht ganz eindeutig: Einerseits kann man 
darunter jede plötzlich auftretende
Baucherkrankung
verstehen \citep{Longmore2007,Renz-Polster2006}. 

Allerdings ver\-stehen viele Leute darunter eher eine
bedrohliche Bauch\-erkrankung, bei der man akut (schnell)
handeln muss (Ab\-wehrspannung, harter Bauch). In der
Notfallmedizin halten wir uns an letztere Definition.}


\Text{\TxBeschreibung}
{Das \E{Akute Abdomen} ist eine Kombination von bestimmten
Symptomen, die
durch viele verschiedene Erkrankungen im Bauchraum ausgelöst
werden kann. Diese zu Grunde liegenden Erkrankungen können
plötzlich entstehen, oder sich über längere Zeit entwickeln
und erst im fortgeschrittenen Stadium die Symptome des
akuten Abdomens auslösen. Ein akutes Abdomen erfordert eine
rasche (spitalsmäßige) Behandlung.
}
{

}{}{}{}

\TextSlide{\TxSymptome}
{%
Der Patient liegt meistens im Bett. Oft ist eine
Schonhaltung (angezogene Beine, gekrümmte Körperhaltung) zu
beobachten. Er klagt über \EE{Schmerzen} im Bauch, diese können
als diffus (im gesamten Bauch) oder auf ein Bauchsegment
beschränkt angegeben werden.

Die \EE{harte Bauchdecke} (\enquote{\E{bretthart}}) ist
neben den Schmerzen das Leitsymptom des Akuten Abdomens.
Es besteht oft Übelkeit oder auch Erbrechen. Der Kreislauf
kann kompensiert sein, oder Zeichen von Schwäche zeigen, bis
hin zu einer Schocksymptomatik (Blutverlust, Sepsis).
Weitere Symptome sind abhängig von der zu Grunde liegenden
Erkrankung.

Diese Auflistung der Symptome ist in der Praxis nicht
unbedingt vollständig bei einem Patienten nachzuweisen. Die
Intensität und die Anzahl der angeführten Symptome ist
einerseits vom Allgemeinzustand, vom Lebensalter, von
anderen vorbestehenden Erkrankungen des Patienten und der
Dauer des aktuellen Zustandes abhängig.

Der Patient ist nicht unbedingt unmittelbar vital
bedroht. Wesentlich nach Eintreffen beim Patienten ist die
Beantwortung der Fragen \Speech{\enquote{Wie schlecht, bzw.
wie gut geht es dem Patienten? Ist der Patient \uline{akut}
vital bedroht?}}.

} 
{
\begin{itemize}
  \item \EE{Schmerz}: diffus (ganzer Bauch) oder
  lokalisierbar
  \item \EE{Harte  Bauchdecke} (bretthart)
  \item Schonhaltung  (angezogene Beine, Patient
  krümmt sich)
  \item Übelkeit, Erbrechen
  \item Evtl. Stuhl-, Windverhalten
  \item Evtl. Fieber
  \item Evtl. Schock
\end{itemize}
}{}{}{}




\TextSlide{Ursachen}
{%
Grundsätzlich kommt fast jede Erkrankung oder Verletzung in
Betracht die Strukturen und Organe im Bauchraum betrifft,
entsprechend vielfältig sind auch die Ursachen: 

\begin{multicols}{2}
\FormatTable
\begin{itemize}
  \item Bauchfellentzündung (Peritonitis)
  \item Darmverschluss (Ileus)
  \item Blinddarmentzündung (Appendizitis)
  \item Entzündung der Bauchspeicheldrüse  (Pankreatitis)
  \item Entzündung des Magens / Magengeschwür
  \item Durchbruch / Zerreißung eines Hohlorganes
  \item Blutungen in den Bauchraum
  \item Bauchaortenaneurysma
  \item Mesenterialinfarkt
  \item andere abdominelle Erkrankungen
  \item Entzündung der Leber (Hepatitis )
  \item Gynäkologische Erkrankungen
  \item \ldots
\end{itemize}
\end{multicols}
}{
\begin{itemize}
  \item Bauchfellentzündung 
  \item Blinddarmentzündung 
  \item Entzündung der Bauchspeicheldrüse  (Pankreatitis)
  \item Durchbruch eines Hohlorganes
  \item Blutungen in den Bauchraum
  \item Bauchaortenaneurysma
  \item Gynäkologische Erkrankungen
  \item u.\,v.\,a.\,m. \ldots
\end{itemize}
}{}{}{}






\subsubsection{Darmverschluss}
\index{Darmverschluss}
\label{topic:ileus}
\HeadingSub{\Lang{Term.} \T{Ileus}
(vollständig), \Z{\T{Subileus}  (unvollständig)}}

\Definition{Unterbrechung der normalen  Darmpassage}

\Text{\TxBeschreibung}
{%
Der Verschluss kann mechanisch bedingt sein (durch
Hindernisse wie Verwachsungen oder Tumore), oder durch eine
Lähmung der Darmmuskulatur (Medikamente, Gifte, Entzündung
\ldots) verursacht sein. Bestimmte Schmerzmittel
\Z{(Opiate)}, welche oft bei Krebspatienten eingesetzt
werden, sind bekannt dafür Darmlähmungen zu erzeugen. 
}{%
\begin{itemize}
  \item Mechanische Blockade oder Lähmung
\end{itemize}
}{}{}{}


\TextSlide{\TxSymptome}
{%
Typisch für den Darmverschluss ist die mangelhafte
Stuhlproduktion. Es kommt zu diffusen Bauchschmerzen und einem
Blähbauch, sowie zu Übelkeit und Erbrechen. In seltenen
Fällen kann es auch zum Erbrechen von zurück
gestautem Kot kommen \Z{(\I[Miserere|Z]{Miserere})}.

\bigskip

\SlideCompensationNoParskip{1}
}{
\begin{itemize}
  \item Diffuse Bauchschmerzen
  \item Stuhlverhalten
  \item Übelkeit, Erbrechen (selten: Koterbrechen)
\end{itemize}
}{}{}{}


\subsubsection{Blinddarmentzündung}
\label{topic:appendizitis}
\HeadingSub{\Lang{Term.} \T{Appendizitis}}


\Definition{Entzündung des \EE{Wurmfortsatzes}
(\EE{Appendix} \Z{vermiformis}) des Blinddarms}

\TextSlide{\TxSymptome{} und Verlauf}
{%
\ACh{GABG}{2009-05-07}{NEU}Im Vordergrund der akuten Appendizitis
steht anfänglich meist ein als eher diffus empfundener,
heftiger Schmerzzustand im Bereich der
Nabelgegend\ACh{GABS}{2009-05-07}{Magengegend --> Nabelgegend}, der
innerhalb weniger Stunden in den rechten Unterbauch
wandert. Übelkeit, Brechreiz und Verstopfung
\ACh{GABS}{2009-05-07}{Obstipation --> Verstopfung} können
das Krankheitsbild begleiten. Meist ist die Körpertemperatur
bis 38\DegC{} leicht erhöht
\ACh{GABS}{2009-05-07}{gelöscht da unerheblich für RS: , wobei eine rektale und
axilläre Messung eine Differenz um ca. 1\DegC aufweist}. 
\ACh{GABS}{2009-05-07}{alt: Unbehandelt kann der weitere Verlauf durch eine
Abszeßbildung und Perforation in
die Bauchhöhle gekennzeichnet sein. neu: Unbehandelt kann es zu einer Eiterung
und einem Durchbruch in die Bauchhöhle kommen.} Unbehandelt kann es zu einer
Eiterung und einem Durchbruch in die Bauchhöhle, einer
anschließenden Bauchfellentzündung und einem akuten Abdomen
kommen (\FullPrettyRef{topic:akutes-abdomen}).

Der Verlauf einer akuten Appendizitis ist jeweils durch
individuelle Faktoren wie Lebensalter und Allgemeinzustand
des Patienten geprägt.
}{
\begin{itemize}
  \item Schmerzen im re. Unterbauch
  \item \E{Druck}schmerzen im re. Unterbauch
  \item \E{Loslass}schmerz li. Bauch
  \item Abwehrspannung
  \item Übelkeit, Erbrechen
  \item Evtl. Fieber
\end{itemize}
}{}{}{}



\subsubsection{Gallenkolik}
\label{topic:gallenkolik}
\index{Kolik!Gallen-}
\index{Gallenkolik}

\Definition{Kolik-artige Schmerzen infolge einer Blockade
der Gallenwege durch Gallensteine}


\TextSlide{\TxBeschreibung} 
{Gallensteine verursachen eine Verstopfung des
Gallengangs und folglich einen Rückstau von
Gallenflüssigkeit in die Leber. Dabei wird der Gallengang
gedehnt und es kommt zu starken Schmerzen.} 
{
\begin{itemize}
  \item Verstopfung der Gallenwege (Stein)
  \item → Rückstau und Dehnung
\end{itemize}
}{}{}{}

\SlideCompensationNoParskip{2}

\TextSlide{\TxSymptome}
{%
Durch den Rückstau und der Dehnung des Gallengangs kommt es
zu starken, auf- und abschwellenden, krampfartigen
(\EE{kolikartigen}) \EE{Schmerzen} im \EE{rechten
Oberbauch}, klassischerweise mit Ausstrahlung in die rechte
Schulterregion. Begleitendes Fieber ist möglich.

Der rechte Oberbauch ist äußerst druckschmerzhaft, der
Patient ist sehr \EE{unruhig} und agitiert.

\bigskip

\bigskip

\bigskip
% \SlideCompensationNoParskip{1}
}{
\begin{itemize}
  \item Auf- und abschwellende, krampfartige
  (kolikartige) Schmerzen im re. Oberbauch
  \item Nach fettreichen Mahlzeiten
  \item Evtl. Gelbfärbung des Augenweiß (Gelbsucht)
  \item Evtl. Übelkeit, Erbrechen
\end{itemize}
}{}{}{}



\subsubsection{Gastrointestinale
Blutungen und Ösophagusvarizenblutung}
\label{topic:oesophagusvarizenblutung}
\label{topic:gastrointestinale-blutung}
\TODO{EMHO}{2010-08-04}{andere Blutungsquellen sind
anzuführen}{}

\Definition{Blutungen im Magen-Darm-Trakt bzw. blutende \enquote{Krampfadern} in der
Speiseröhre.}\nopagebreak[4]

\TextSlide{\TxBeschreibung}
{Im Verlauf des Magen-Darm-Traktes kann es an verschiedenen
Stellen zu Blutungen kommen, die auch sehr schwer ausfallen
können. 

Blutungen im \EE{oberen Teil} des Verdauungstraktes können
durch 
\T[Oesophagusvarizen@Ösophagusvarizen]{Ösophagusvarizen} 
verursacht werden: Bei Patienten
mit einer fortgeschrittenen Lebererkrankung (Leberzirrhose) kann es zur Bildung von Krampfader-artigen Gefäßerweiterungen in der Speiseröhre
kommen. Diese Gefäße können massive Blutungen verursachen,
die tödlich enden können. Auch Blutungen im Rahmen eines
Magengeschwürs sind möglich, aber selten dermaßen stark
ausgeprägt.

\Z{Blutungen im \EE{mittleren Teil} des Magen-Darm-Traktes
können durch diverse Erkrankungen und Tumore verursacht
werden. Blutungen im \EE{Endteil} kommen oft durch Tumore
oder infolge eines Hämorrhoidalleidens zu Stande.}


}
{
\begin{itemize}
    \item Oberer Teil: Ösophagusvarizzen
    
    (Oft bei
    Lebererkrankungen, hoher Blutverlust möglich)
    \item \Z{Diverse Erkrankungen, Tumore}
\end{itemize}
}{}{}{}

\TextSlide{\TxSymptome}
{
Bei Blutungen im \EE{oberen Magen-Darm-Trakt}
(Speiseröhre, Magen) kommt es je nach Stärke zu
schwallartigem Erbrechen  von rotem, frischen Blut
(\Z{\I{Hämatemesis}}) oder häufiger zum Erbrechen von
\II[Erbrechen!Kaffeesatz-artiges]{Kaffeesatz-artigem}
Mageninhalt: Durch die Magensäure kommt es zur Veränderung
des Blutes, es wird bräunlich, ähnlich dem Kaffeesatz im
Kaffeefilter. Bei Blutungen im \EE{mittleren Teil} des
Verdauungstraktes färbt sich der Stuhl tief schwarz, man
spricht vom sog. \T{Teerstuhl} (\T{Melena}. Bei Butungen im
\EE{Endteil} wird rötliches Blut ausgeschieden. 

}{
\begin{itemize}
  \item Oberer Teil: Erbrechen von
  \begin{itemize}
    \item Blut im Schwall
    \item Kaffeesatz-artiger Mageninhalt
  \end{itemize} 
  \item Mittlerer Teil: Teerstuhl
  \item Endteil: Blutiger Stuhl
\end{itemize}
}{}{}{}

\textleaf


\Fazit{Der Patient kann aufgrund des Blutverlustes akut
vital bedroht sein, es besteht die Gefahr eines
hypovolämischen Schocks.}






\subsubsection{Magen-Darm-Grippe, Gastroenteritis,
Lebensmittelvergiftung}
\label{topic:gastroenteritis}
\label{topic:lebensmittelvergiftung}
\TODO{EMHO}{2010-08-04}{Lebensmittelvergiftung: Text
fehlt}{}

\HeadingSub{\Lang{Syn.} Brechdurchfall,
\Z{gastrointestinaler Infekt, Gastroduodenitis}}

Die Magen-Darm-Grippe ist eine chronisch oder akut
verlaufende Entzündung der Magen- und der Dünndarmschleimhaut. 

\TextSlide{\TxBeschreibung{} und Ursachen}
{%
Die Gastroenteritis wird durch Erreger wie Bakterien oder
Viren hervorgerufen. Es gibt völlig harmlose aber auch
lebensbedrohliche Verlaufsformen, dies ist abhängig vom
Erreger und letztendlich auch vom Alter des Patienten. Durch
den Durchfall und das Erbrechen kann es zu einem maßgeblichen
Verlust von Wasser und Elektrolyten kommen. Wenn der Patient
nicht ausreichend trinken kann besteht die Gefahr der
Austrocknung (\I[Exsikkose!bei gostrointestinalem
Infekt]{Exsikkose}). 
}{
\begin{itemize}
  \item Vielfältig
  \item Häufig Infektionen
\end{itemize}
}{}{}{}


\Text{\TxSymptome}
{
Klassische Zeichen einer Magen-Darm-Grippe sind Erbrechen
(\Z{\I{Emesis}}) und Durchfall (\Z{\I{Diarrhoe}}). Dazu hat
der Patient meist starke, kaum lokalisierbare
Bauchschmerzen. Infolge des 
Erbrechens und des Durchfalls kann es zu einem erheblichen 
Wasser- und Elektrolytverlust und Symptomen der Exsikkose
 kommen. 

}{
\begin{itemize}
    \item Erbrechen 
    \item Durchfall 
    \item Bauchschmerzen
    \item Exsikkosezeichen
\end{itemize}
}{}{}{}

\TextSlide{Besonders gefährdet}
{%
Besonders betroffen davon sind Kleinkinder, da sie
normalerweise einen höheren Wasseranteil am Körpergewicht
haben.  Die Symptome sind klassisch: Oberbauchkrämpfe,
die ziemlich heftig sein können, Übelkeit, Erbrechen, Durchfälle.
Exsikkose-Gefahr! 
}{
\begin{itemize}
    \item (Klein-)Kinder
    \item Alte Menschen
\end{itemize}
}{}{}{}






\subsubsection{Flüssigkeitsmangel: Exsikkose}
\ACh{WITM}{????-??-??}{NEU: Schlagworte}
\ACh{GABG}{2009-05-11}{NEU: Fließtexte}
\ACh{GABS}{2009-05-14}{Fließtexte stark modifiziert und angepasst. }

\HeadingSub{\Lang{Syn.} \T{Dehydratation}}

\Definition{\enquote{Austrocknung}: Es besteht  ein
hochgradiges Defizit von Flüssigkeit im Körper.}

Grundsätzlich wird (wurde) entweder \E{zu wenig Flüssigkeit
zugeführt} oder \E{zu viel Flüssigkeit an die Außenwelt
abgegeben} (eine Kombination aus beidem ist natürlich auch
möglich). Dafür gibt es jeweils sehr verschiedene Ursachen:

\TextSlide{Ursachen}
{%
Ältere Menschen haben oft ein mangelndes Durstgefühl und daher
einen verminderten Trinkantrieb. Auch durch die Einnahme von
Drogen (Ecstacy) kann auf den Durst \enquote{vergessen} werden.  

Andererseits kann es durch \E{Erbrechen}, \E{Durchfall}
sowie \E{starkes Schwitzen} zu einem stark erhöhten
Flüssigkeitsverlust kommen. Bei \E{Fieber} wird auch
vermehrt Feuchtigkeit abgeatmet, und natürlich können
auch \E{Medikamente} \enquote{entwässern}.
}{
\begin{itemize} %GABS: umgebaut
  \item Flüssigkeitszufuhr ${\downarrow}$
  \begin{compactitem}
    \item Verminderter Trinkantrieb
  \end{compactitem}
  \item Flüssigkeitsverlust ${\uparrow}$
  \begin{compactitem}
    \item Erbrechen, Durchfall
    \item Starkes Schwitzen
    \item Fieber (Feuchtigkeit wird vermehrt abgeatmet)
    \item Medikamente
  \end{compactitem}
\end{itemize}
}{}{}{}

Die genannten Ursachen sind natürlich nur ein kleiner Teil
von sehr vielen Möglichkeiten.

\TextSlide{\TxSymptome}
{%
Der Patient wirkt allgemein \E{geschwächt bis apathisch}, oft
\EE{verwirrt}, klagt über mehr oder minder heftiges
Durstgefühl, ev. über Wadenkrämpfe und Erbrechen. Die Haut
ist auffallend trocken, \EE{Falten sind abhebbar und bleiben
bestehen}. \EE{Mundhöhle und Zunge sind trocken}, evtl.
borkig oder mit eingetrocknetem Schleim belegt. Die Augen
sind charakteristischerweise tiefliegend, die Nase wirkt
auffallend spitz. Der RR ist mitunter sehr nieder, die HF
eher hoch. \E{Der Patient ist bis zur Klärung der Situation
nüchtern zu belassen!}
}{
\begin{itemize}
    \item RR ${\downarrow}$
    \item stehende Hautfalten
    \item trockene, belegte Zunge
    \item Elektrolytentgleisung
    \item \EE{Verwirrtheit, Schwindel, Apathie}
\end{itemize}
}{}{}{}




\subsubsection{Magen- und Zwölffingerdarmgeschwür}
\label{topic:gastritis}

\Z{\HeadingSub{\Lang{Synn.} Gastritis, Duodenitis}}

\Definition{Magengeschwür: Entzündung der
Magenschleimhaut mit Substanzdefekt.}

\TextSlide{Ursachen}
{
Die häufigste Ursache für ein Magengeschwür ist eine
chronische Infektion mit einem Bakterium \Z{(Helicobacter
pylori)} \citep{Malfertheiner2004}. Medikamente, wie
z.\,B. nicht-steroidale Antirheumatika (NSAR, z.\,B.
Voltaren{\texttrademark} (Diclofenac),
Aspirin{\texttrademark} (Acetylsalicylsäure, ASS),etc.)
schädigen die Magenschleimhaut und können ebenso bei
häufiger Einnahme zu einem Geschwür führen. Auch andere
Stoffe können auf chemisch-toxischen Weg zu einer
entsprechenden Schädigung führen. \Z{Ein
Gallensäurenrückfluß (Reflux) bzw. Autoimmunerkrankungen
können ebenso zu einem Geschwür führen.} 
}{
\begin{itemize}
    \item Bakteriell \Z{(Helicobacter pylori)}
    \item Medikamente
    \item Chemisch-toxisch
    \item \Z{Gallensäurenrückfluß}
    \item \Z{Autoimmunerkrankungen}
\end{itemize}
}{}{}{}

\TextSlide{\TxSymptome}
{Der Patient klagt über einen
\EE{stechend-schneidenden Schmerz}, welcher sehr stark sein
kann. Daneben kommt es zu Appetitlosigkeit, Völlegefühl,
Übelkeit und Erbrechen. Wenn gleichzeitig eine \EE{Blutung}
besteht, dann sind \II[Teerstuhl!Gastritis]{Teerstühle}
(schwarz, übelriechend) oder \EE{Kaffeesatz-artiges
Erbrechen} zu beobachten. 
}{
\begin{itemize}
    \item Schmerzen (Oberbauch), Übelkeit, Erbrechen
    \item Kaffeesatz-artiges Erbrechen, Teerstuhl
\end{itemize}
}{}{}{}


\subsubsection{Pankreatitis}

\HeadingSub{\Lang{Syn.} \T{Bauchspeicheldrüsenentzündung}}
\label{topic:pankreatitis}

Wie nahezu jedes Organ des Körpers kann sich auch die
Bauchspeicheldrüse entzünden. Die
Bauchspeichedrüsenentzündung (Pankreatitis) kann dabei
sowohl \E{akut}, als auch \E{chronisch} verlaufen. 

% \Z{\Lang{ICD-10} Akute Pankreatitis: K85, chronisch
% Pankreatitis: K86.1}

\TextSlide{Ursachen}
{\index{Gallensteine!Pankreatitis}
\index{Alkoholmißbrauch!Pankreatitis}
Die häufigsten Ursachen für eine Pankreatitis sind
\EE{Gallensteinleiden} und \EE{Alkoholmißbrauch}
\citep{SiewertChirurgie8-07.14-Pankreas,ScholzNotfallmedizin2}.
Daneben gibt es noch viele, allerdings nicht so häufige
Ursachen, oft bleibt die Ursache auch unbekannt
\citep{HeroldInnereMedizin2005}.

\Speech{Warum können Gallensteine eine Pankreatitis
verursachen?} Bei vielen Menschen mündet der
Pankreasausführungsgang, welcher Verdauungsenzyme in de
Zwölffingerdarm abgibt, nicht direkt in den Dünndarm,
sondern in den unteren Abschnitt des Gallengangs.
Verschließt nun ein Stein den Gallengang unterhalb der
Einmündung, so staut sich Galle und Pankreassaft in die
Bauchspeicheldrüse zurück. Es kann sogar dazu kommen, dass
sich das Pankreas durch die Verdauungsenzyme selber
an-/verdaut.

}
{\begin{itemize}
  \item Gallensteinleiden
  \item Alkoholmißbrauch
  \item \ldots
\end{itemize}}
{}
{}
{}

\TextSlide{\TxSymptome}
{Die Symptome sind oft unspezifisch, eine
Erstdiagnosestellung ist präklinisch normalerweise nicht
möglich. Charakteristisch ist bestenfalls ein
\II{gürtelförmiger
Oberbauchschmerz}\index{Oberbauchschmerz,gürtelförmiger}.
Die häufigsten Begleitsymptome sind Übelkeit und Erbrechen,
sowie Fieber. \citep{HeroldInnereMedizin2005}


\bigskip

\bigskip
}
{\begin{itemize}
  \item Oberbauchschmerzen, oft gürtelförmig
  \item Übelkeit, Erbrechen
  \item Fieber
\end{itemize}}
{}
{}
{}

\TextSlide{\TxKomplikationen}
{Je nach Schwere und Dauer der Entzündung kann es zum
Selbstverdauen der Bauchspeicheldrüse bzw. zu einer
\E{Ruptur} kommen. In Folge tritt ein
\I[Akutes Abdomen!bei Pankreatitis]{akutes Abdomen} auf
(\FullPrettyRef{topic:akutes-abdomen}) und der Patient kann
kann einen \I[Schock!septischer!bei Pankreatitis]{septischen
Schock} erleiden.

\bigskip

}
{\begin{itemize}
  \item Ruptur
  \item Akutes Abdomen
  \item Septischer Schock
\end{itemize}}
{}
{}
{}




\subsubsection{Mesenterialinfarkt}
\label{topic:mesenterialinfarkt}

Beim Mesenterialinfarkt kommt es zum Verschluss einer den
Darm versorgenden Arterie
\Z{(\I[Mesenterialarterie|Z]{Mesenterialarterie}\ZFootnote{Das
\I[Mesenterium!Mesenterialinfarkt]{Mesenterium}
(\I{Darmgekröse}) ist die Aufhängung von Teilen des
Dünndarms (\FullPrettyRef{topic:mesenterium}). Darin
verlaufen venöse und arterielle Gefäße zum und vom Darm.})}
mit einem Absterben des entsprechenden Darmabschnittes. Er
ist somit ein \I[Ischämie!Mesenterialinfarkt]{ischämischer}
Notfall.

\TextSlide{\TxSymptome}
{Die Symptome sind eher unspezifisch. Der Patient klagt oft
(aber nicht immer) über \EE{abdominelle Schmerzen}, welche
oft dumpf und diffus sind. Es kann für einige Stunden ein
\E{stilles Intervall} folgen, in welchem die Schmerzen eher
gering sind. 

Nach ungefähr \VonBis{12}{24}\Hour* wird ein
\EE{Darmverschluss} bemerkbar und die Darmwand verliert
zunehmend an Dichtigkeit, es kann zu einem \EE{Durchbruch}
(Ruptur) kommen. Es folgt eine \EE{Bauchfellentzündung} und
ein \II[Akutes Abdomen!beim Mesenterialinfarkt]{akutes
Abdomen} (\FullPrettyRef{topic:akutes-abdomen}). Es kann zu
\EE{blutigen, übel riechenden Stuhlabgängen} kommen.
\cite{SiewertChirurgie8-06-Gefaesschirurgie}
% 
% \bigskip

}
{\begin{itemize}
  \item Abdominalschmerzen
  \item Übelkeit, Erbrechen
  \item $\sim$ 6\Hour* stilles Intervall
  \item Nach $\sim$ \VonBis{12}{24}\Hour* Durchbruch und
  Bauchfellentzündung, Ileus
  \item Übelriechende, blutige Stuhlabgänge
\end{itemize}}
{}
{}
{}




\subsubsection{Bauchfellentzündung}
\label{topic:bauchfellentzuendung}
\label{topic:peritonitis}

\HeadingSub{\Lang{Term.} \T{Peritonitis}}

\Definition{Entzündung und Keimbesiedelung des 
Bauchfells}

\Text{\TxBeschreibung}{%
Durch die Infektion und Entzündung des Bauchfelles kommt es
zu Symptomen eines Akuten Abdomens. Oft ist die Perforation
eines Hohlorganes  (Magendurchbruch, Darmperforation,
Blinddarmdurchbruch, \ldots) Ursache der Entzündung. 
}{

}{}{}{}

\TextSlide{Ursache}{
\begin{itemize}
    \item Perforation eines Hohlorganes (Darm, Magen,
    Gallenblase, \ldots)
    \item Fremdkörper
\end{itemize}
}{
\begin{itemize}
    \item Perforation eines Hohlorganes
    \item Infektion über Blut- / Lymphwege oder eingewandert
    über Fremdkörper
\end{itemize}
}{}{}{}

\TextSlide{\TxSymptome}{
\begin{itemize}
    \item Symptome eines Akuten Abdomens
    \item Schmerzen (diffus)
    \item Evtl. reflektorischer Darmverschluss
    \item Evtl. Sepsis- und/oder Schockzeichen
\end{itemize}

\bigskip

}{
\begin{itemize}
    \item Symptome eines Akuten Abdomens
    \item Schmerzen (diffus)
    \item Evtl. reflektorischer Darmverschluss
    \item Evtl. Sepsis- und/oder Schockzeichen
\end{itemize}
}{}{}{}


\subsection{Systemische Hitzeschäden}
\label{topic:system-Hitzeschaeden}

\subsubsection{Hitzekollaps}
\label{topic:hitzekollaps}

\TextSlide{\TxBeschreibung{} und \TxSymptome{}}{%
Bei einem Wärmestau, bei dem die zugeführte Wärme
größer ist, als die abgegebene, kann es zu einem so
genannten Hitzekollaps kommen. 

Die \EE{Blutgefäße} in der Haut \EE{erweitern} sich und das
Blut versackt, der \EE{Blutdruck fällt ab} und die
Durchblutung des Gehirns nimmt kurzfristig ab. Die Folge ist
eine \EE{kurze Ohnmacht}, den Patient wird \enquote{schwarz vor
Augen}. Lagert man ihn flach, oder eventuell mit erhöhten
Beinen, und kühlt ihn, bzw. bringt ihn an einen kühlen Ort,
so bessert sich sein Zustand recht rasch.
}{
\begin{itemize}
  \item Wärmestau
  \item Erweiterung der Blutgefäße in der Haut (maximal)
  \item RR\,${\downarrow}$ , Durchblutung des Gehirns
  kurzzeitig ${\downarrow}$
  \item kurze Ohnmacht, \enquote{schwarz vor Augen}
\end{itemize}
}{}{}{}




% \MassnahmenNG{Spezielle Maßnahmen}{Hitzekollaps}{}{
% 
% \begin{itemize}
%   \item Flachlagerung: ${\rightarrow}$ Gehirndurchblutung
%   wieder ${\uparrow}$, bewusstseinsklar
%   \item nach Sekunden gebessert, \enquote{regelt sich von
%   selbst}
%   \item Patient sitzen lassen, Flüssigkeit
% \end{itemize}
% }


\subsubsection{Hitzeerschöpfung}
\label{topic:hitzeerschoepfung}

\TextSlide{\TxBeschreibung}
{
Bei der Hitzeerschöpfung ist die \EE{Wärmeaufnahme
dauerhaft größer als die Wärmeabgabe}. Der Körper versucht
dies durch starkes Schwitzen zu kompensieren, dadurch kommt 
es zu einem \EE{starken Wasser- und Elektrolytverlust}. Der
\EE{Blutdruck fällt}, es kommt zu einer \EE{Tachykardie}, die
Körpertemperatur bleibt jedoch noch normal. 
Hier ist es wichtig den Patienten zu
kühlen.

\vspace{1cm}

}{
\begin{itemize}
  \item Wärmeaufnahme {\textgreater}  Wärmeabgabe
  \item Starkes Schwitzen, Wasser- und  Elektrolytverlust
  \item Starkes Durstgefühl
  \item RR ${\downarrow}$, Tachykardie
  \item Körpertemperatur normal
\end{itemize}
}{}{}{}




% \MassnahmenNG{Spezielle Maßnahmen}{Hitzeerschöpfung}{}{
% 
% \begin{itemize}
%   \item Kühlung, Schatten
%   \item Flüssigkeitszufuhr, wenn bewusstseinsklar
%   \item Neurocheck inkl. BZ-Messung
% \end{itemize}
% }



\subsubsection{Hitzschlag}
\label{topic:hitzschlag}

\TextSlide{\TxBeschreibung{} und \TxSymptome{}}{
Beim Hitzschlag ist nun auch die \EE{Regulation der Körpertemperatur gestört},
es sind auch \E{sehr hohe Werte} (bis 43\,\DegC{})
möglich. Dies bedeutet \EE{akute Lebensgefahr}. 

Der Patient hat eine zunächst trockene, rote und weiße,
später grau-bläuliche Haut, der Kopf ist oft hochrot.
Weiters kommt es zu \EE{neurologischen Symptomen}, wie massiven
Kopfschmerzen, Schwindel, Erbrechen und
Bewusstseinsstörungen; bis hin zur Bewusstlosigkeit. Der Patient 
kann Schockzeichen zeigen, evtl. auch eine schnelle, flache
Atmung.
}{
\begin{itemize}
  \item Kopfschmerz, Schwindel, Erbrechen
  \item Haut zunächst trocken, rot und heiß, später grau-blau, Kopf
  hochrot
  \item Stark  erhöhte  Körpertemperatur (bis 
  43{\,\DegC} möglich, \EE{Lebensgefahr!})
  \item Bewusstseinsstörung, Bewusstlosigkeit
  \item Schockzeichen
  \item schnelle, flache Atmung
\end{itemize}
}{}{}{}



\subsubsection{Maßnahmen bei systemischen Hitzeschäden}

\MassnahmenNG{Spezielle Maßnahmen}{Hitzekollaps,
Hitzeerschöpfung, Hitzschlag}{}{
\FormatTable

\begin{tabularx}{\linewidth}{XXX}
\TabHeading{Hitzekollaps}
&
\TabHeading{Hitzerschöpfung}
&
\TabHeading{Hitzschlag}
\\\addlinespace\midrule\addlinespace
% \begin{minipage}{\linewidth}
\begin{itemize}
  \item Flachlagerung: ${\rightarrow}$ Gehirndurchblutung
  wieder ${\uparrow}$, bewusstseinsklar
  \item nach Sekunden gebessert, \enquote{regelt sich von
  selbst}
  \item Patient sitzen lassen, Flüssigkeit
\end{itemize}
% \end{minipage}
&
% \begin{minipage}{\linewidth}
\begin{itemize}
  \item Kühlung, Schatten
  \item Flüssigkeitszufuhr, wenn bewusstseinsklar
%   \item Neurocheck inkl. BZ-Messung
\end{itemize}
% \end{minipage}
&
% \begin{minipage}{\linewidth}
\begin{itemize}
  \item \TxMassVit
  \item Beengende Kleidungsstücke lockern
  \item Flachlagerung an kühlem Ort (Schatten), Beine und
  Kopf hochlagern.
  \item Kühlung von \E{außen} (Fächer, mit Eiswürfeln
  abreiben etc.)
  \item Kühlung von \E{innen}: für \EE{ansprechbare}
  Patienten viel kühle
  (alkoholfreie) Flüssigkeit, evtl. Elektrolyt-Getränke
%   \item Neurocheck inkl. BZ-Messung
\end{itemize}
% \end{minipage}
\\\addlinespace\bottomrule
\end{tabularx}
}

\clearpage 
\subsubsection{Sonnenstich}
\label{topic:sonnenstich}

\Definition{Hitzeschaden durch Hirnerwärmung infolge von
direkter Sonneneinstrahlung auf den \textbf{ungeschützten
Kopf}. }

\TextSlide{\TxBeschreibung{} und \TxSymptome{}}{
Der Sonnenstich ist eine Sonderform des Hitzeschadens: 
er entsteht durch \EE{direkte Sonneneinstrahlung} auf den
ungeschützten Kopf. Dadurch steigt die Temperatur im Kopf,
es kommt zu einer Reizung der Hirnhäute. 

Die Folge sind Symptome wie bei
einer Hirnhautentzündung (Meningitis) beziehungsweise andere
\EE{neurologische Symptome}, wie z.\,B. Kopfschmerz, sind die
Folge. Die \EE{Gesichts- und Kopfhaut ist hochrot und heiß}, der 
restliche Körper bleibt dagegen eher kühl. Der Patient
wirkt abgeschlagen und klagt über heftige 
Kopfschmerzen, Schwindel und Übelkeit, auch Unruhe, 
Verwirrtheitszustände und Nackensteifigkeit sind zu
beobachten. In schweren Fällen kann es sogar bis zu 
Krampfanfällen und Bewusstlosigkeit kommen.

\vspace{4cm}

}{
\begin{itemize}
  \item Direkte Sonneneinstrahlung auf \textbf{ungeschützten
  Kopf} ${\rightarrow}$ Reizung der Hirnhäute  
  \item Symptome ähnlich einer
  Hirnhautentzündung (Nackensteifigkeit) 
  \item Evtl. kombiniert mit anderen Hitzeschäden
  \item Gesichts- \& Kopfhaut hochrot, heiß, restlicher Körper kühl
  \item Abgeschlagenheit, heftiger Kopfschmerz, Schwindel
  \item Unruhe, Verwirrtheitszustände, Brechreiz
  \item Schwere Fälle: Krampfanfälle, Bewusstlosigkeit
\end{itemize}
}{}{}{}





\MassnahmenNG{Spezielle Maßnahmen}{Sonnenstich}{}{

\begin{itemize}
%   \item Vitale Bedrohung einschätzen
%   \TxBeiVit \TxMassVit
  \item Flachlagerung an kühlem Ort (Schatten), Kopf hochlagern
  \item Beengende Kleidungsstücke lockern
  \item Kühlung von außen (Fächer, mit Eiswürfeln abreiben, kalte
  Umschläge auf die Stirn etc.)
  \item Kühlung von innen: für \EE{ansprechbare} Patienten kühle
  (alkoholfreie) Flüssigkeit, evtl. Elektrolyt-Getränke
%   \item Schocklagerung u. -bekämpfung. % REVIEW 2010-05-19
%   \item Neurocheck inkl. BZ-Messung
\end{itemize}
}


\clearpage % TEMP temp clearpage
\subsection{Unterkühlung}
\label{topic:unterkuehlung}

\Definition{Bei einer Körperkerntemperatur unter
36\DegC* spricht man von Unterkühlung.}

\begin{wrapfigure}{r}{4cm}
\includegraphics[width=4cm]{./images/teaser/huette_winter_mariazell.jpg}
\end{wrapfigure}
Ähnlich
wie beim Schock kommt es auch  bei der Unterkühlung zur
Zentralisation  des Blutes, um die  Körperkerntemperatur
aufrecht erhalten  zu können. Je niedriger diese 
Kerntemperatur jedoch sinkt, desto  schwerwiegender sind die
Folgen!

Es besteht die Gefahr, dass kaltes Blut aus der Peripherie
zum Körper hin fließt und sich mit dem noch relativ warmen
zentralisierten Blut vermischt. Es kommt dann zu einer
weiteren Unterkühlung des Körperstammes, im schlimmsstem Fall
zu
einem \T{Bergungstod} (Blutfluss in Folge von Bewegungen während des
Bergungsversuches).

\paragraph{Phasen}

\begin{itemize}
  \item \T{Kampfphase}: 
  \VonBis{36,5}{34}\DegC*, Kältezittern,
  Herzfrequenz steigt
  ({\textgreater}\,100\,\nicefrac{x}{min}), schnelle Atmung,
  Schmerzen.
  \item \T{Erschöpfungsphase}:
  \VonBis{34}{30}\DegC*, Bewusstseinstrübung,
  Kältestarre, Herzfrequenz sinkt wieder
  ({\textless}\,60\,\nicefrac{x}{min}), zu langsame Atmung,
  keine Schmerzen mehr. 
  \item \T{Kältenarkose}:
  \VonBis{30}{27}\DegC*, Bewusstlosigkeit,
  unregelmäßiger Herzschlag, Atempausen.
  \item {\textless}\,27\DegC* ${\rightarrow}$ klinisch tot. 
\end{itemize}




\MassnahmenNG{Spezielle Maßnahmen}{Schwere Unterkühlung
(\textless\,34\DegC*)}{}{

\begin{itemize}
%   \item \TxMassVit
  \item Den Patient so wenig als möglich (nur soviel wie unbedingt nötig)
  bewegen (drohender Bergungstod).
  \item Schutz vor weiterer Abkühlung, \Ca{keine \underbar{aktive}
  Erwärmung!}
  \begin{itemize}
    \item Woll- und Aludecken! Mütze!
  \end{itemize}
\end{itemize}
}



Bei Unterkühlung gilt:

\Fazit{\enquote{Nobody is dead until warm and dead!}
Schwer unterkühlte Patienten können auch noch nach Stunden (!) ohne
Kreislauftätigkeit erfolgreich reanimiert werden!}





\clearpage % TEMP temp clearpage
\subsection {Vergiftungen}
\label{topic:vergiftung}



\Definition{\T{Vergiftung} bedeutet eine (meist  dosisabhängige)
organische  Funktionsbeeinträchtigung bzw. Schädigung des
Organismus durch eine aufgenommene Substanz.}

\dictum[Hohenheim von Theoprast]{Die Dosis macht das Gift.}

Vergiftungen sind eine interessante Sache: Ein an sich
nützlicher Stoff kann sich in eine todbringende Waffe
verwandeln, wenn er im Übermaß zugeführt wird. So ist zum
Beispiel ein Überleben ohne Salz nicht möglich. Führt man
jedoch 20\,dag zu, endet das meist tödlich. Sonst nützliche
Medikamente können schon bei einer Dosisabweichung im
Mikrogramm-Bereich eine fatale Wirkung haben.

Andererseits können Vergiftung in Abhängigkeit vom
ursächlichen Stoff so ziemlich alle Erscheinungsformen
darbieten. Oft findet man nur \EE{unspezifische Symptome}
vor, wie z.\,B. Erbrechen, Durchfall, Schmerzen im
Magen-Darm-Trakt, Bewusstseinsstörungen, Krampfanfälle oder
Entgleisung der Vitalparameter.
Es ist weder möglich, noch sinnvoll, auf sämtliche mögliche
Vergiftungen im Detail einzugehen. Im Folgenden werden
einige häufig anzutreffende Vergiftungen sowie
Herangehensweisen für andere, weniger häufige, vorgestellt
und besprochen. 





\Text{Der Weg in den Körper -- Aufnahme des Giftes}{%
Die Aufnahme der Substanz kann über verschiedene Wege
erfolgen:

\begin{itemize}
  \item \EE{Oral}: über den Mund
  \item \EE{Inhalativ}: Einatmen
  \item \EE{Stechen, spritzen}: in die Vene (intravenös,
  z.\,B.
  Drogenkonsum), in den Muskel, in die %(Unter-)
  Haut (z.\,B. Insektenstich)
  \item \EE{Über die Haut}: transkutan; fettlösliche
  Substanz kann  Hautbarriere durchdringen (z.\,B.:  organische
  Lösungsmittel,  Insektenschutzmittel etc.)
\end{itemize}
}{}{}{}{}





% \subparagraph*{Vorgehensweise}
% 
% 5-Fingerregel:
% 
% \begin{enumerate}
%   \item Vitalparameter
%   \item Giftentfernung
%   \item Antidot-Therapie (Notarzt)
%   \item Gift-Asservierung
%   \item Transport
% \end{enumerate}

\TextSlide{Vergiftungsinformationszentrale}{%
Bei der Vergiftungsinformationszentrale kann Expertenrat
eingeholt werden, sofern der verantwortliche Stoff bekannt
ist. Unter Umständen kann wird hier auch bei der
Identifizierung des Stoffes geholfen.
\Fazit{\EE{Vergiftungsinformationszentrale: {\ttfamily ~01 /
406 43 43}}} }{
Expertenrat

01 / 406 43 43
}{}{}{}






Gegen die meisten Vergiftungen kann man präklinisch,
ursächlich wenig unternehmen. Das Hauptaugenmerk liegt
daher auf: 




\MassnahmenNG{Allgemeine Maßnahmen und
Herangehensweise}{Vergiftungen}{\label{m:am-vergiftung}}{
\begin{enumerate}
  \item Selbstschutz
  \item Beendigung des Giftkontaktes
  \item Evtl. Dekontamination bei
  Gefahrstoffen\footnote{z.\,B. Insektizide, Lösungsmittel
  in der chemischen Industrie, \ldots}
  \item Sicherung der Vitalfunktionen
%   \item Entscheidung, ob es sich um einen akut vital
%   bedrohten Patienten handelt
  \item Quellenforschung (Mistkübel
  durchsuchen)
  \item Nüchtern lassen
  \item NICHT absichtlich Erbrechen
  herbeiführen
  \item Asservierung (Packungen, Erbrochenes in
  Nierentasse mitnehmen!)
\end{enumerate}
}



\subsubsection{Vergiftungen mit Alkohol}

\Layout{20}{Einleitung}{%
Die Vergiftung mit Alkohol ist häufig anzutreffen, weil
\enquote{sozial akzeptiert}.
Wenn der Betroffene \enquote{über's Ziel hinaus schießt},
kann es schnell zu lebensbedrohlichen Zuständen kommen.
}
{}
{./images/TAE/dif-w2-zelt-1.jpg}
{}
{}


\Z{Die akute Alkoholvergiftung ist von der chronischen
Alkoholabhängigkeit (Alkoholkrankheit) zu unterscheiden.}

\TextSlide{Wirkung}
{%
Alkohol  wirkt enthemmend, z.\,T. stimulierend und
stimmungsmodulatorisch. Weiters kommt es zu einer Sprach- 
und Gangunsicherheit. 

\bigskip

\bigskip

}{%
\begin{itemize}
    \item Enthemmend, stimmungsmodulatorisch
    \item Sprach- und Gangunsicherheit
    \item Bewusstseinstrübung
\end{itemize}
}{}{}{}

\TextSlide{\TxGefahren}
{%
Mit zunehmender Menge kommt es zu einem totalen
Kontrollverlust, Bewusstseinstrübung,
Schmerzunempfindlichkeit, Temperaturregulationsstörung und
Inkontinenz\footnote{\emph{Inkontinenz:} Unfähigkeit Harn
und/oder Stuhl zu halten.}. Im Extremfall kann es zu einer 
vegetativen Entgleisung und einem tief komatösen Zustand 
kommen. Spätestens dann besteht akute Lebensgefahr. Weitere 
Gefahren sind Atemdepression, Aspiration, Ausfall der 
Schutzreflexe und -- besonders in der kalten Jahreszeit --
Unterkühlung.
}{%
\begin{itemize}
    \item Koma
    \item Atemdepression
    \item  Ausfall der Schutzreflexe
    \item Aspiration
    \item Schmerzunempfindlichkeit
    \item Unterkühlung
\end{itemize}
}{}{}{}


Es ist unbedingt zu erheben, ob noch andere Substanzen
(Medikamente, Drogen, \ldots) eingenommen wurden
(\E{Mischintoxikation}). Aufgrund der oft reduzierten
Schmerzempfindlichkeit ist besonders auf Verletzungen zu
achten.
Die \EE{vitale Bedrohung} ist vor allem anhand der
Vitalparameter, der körperlichen Untersuchung und des
Bewusstseinszustandes zu beurteilen. 





\Text{Chronische Schäden}
{
\begin{itemize}
    \item Leberschäden, Zirrhose
    \item Aszites (Wasserbauch)
    \item Krampfadern in der unteren Speiseröhre  (Blutungsgefahr)
    \item Verminderte Blutgerinnung
    \item Vitaminmangel
    \item Geistiger Abbau, Gedächtnisschwäche
    \Z{(Korsakow-Syndrom, Wernicke-Enzephalopathie)}
    \item Anämie
    \item Pankreatitis
    \item Fetales Alkoholsyndrom (Schädigung des Fötus in der Schwangerschaft)
    \cite{Jones1973,Guerri2009}
\end{itemize}
}{	

}{}{}{}













\subsubsection{Vergiftungen durch Medikamente}

\Layout{61}{}
{%
Überdosierungen von Medikamenten können komplett
unterschiedliche Symptome erzeugen, je nach eingenommenen
Medikament. Manchmal können die Symptome sogar erst Stunden
nach der Einnahme auftreten und tödlich enden. Vergiftungen
mit Medikamenten gehören daher immer abgeklärt und die
Patienten hospitalisiert. 

Neben unabsichtlichen Fehleinnahmen oder Überdosierungen
werden Arzneimittel oft auch in Selbstmordabsicht
eingenommen. Die Umstände der Einnahme müssen daher
abgeklärt werden.
}{

}{./images/teaser/imgp0482_00800.jpg}{}{Gabriel}

\Z{
\paragraph{Paracetamol (Mexalen\texttrademark,
APA\texttrademark, \ldots)} Paracetamol kann in hoher Dosis
zu Leberversagen führen. Symptome treten erst nach Stunden
auf. Es gibt ein Gegenmittel, dieses muss jedoch
rechtzeitig gegeben werden. Ohne Behandlung kommt es zum
Leberversagen und in Folge zum Tod. 

\paragraph{Digitalis} Digitalis ist ein auf das Herz
wirkendes Medikament und wird seit langer Zeit zur
Behandlung der Herzinsuffizienz und von
Herzrhythmusstörungen eingesetzt. Es ist jedoch sehr
schwierig die richtige Dosis für einen Patienten zu finden,
daher kommt es -- trotz korrekter Einnahme -- häufig zu
Überdosierungen. Die Folge sind Herzrhythmusstörungen und
andere EKG-Veränderungen und \enquote{Farbensehen}.
}

\subsection{Vergiftungen mit Suchtmitteln und
\enquote*{Drogen}}

\subsubsection{Opiate: Heroin \& Co.}

Heroin und Morphium gehören zu der Gruppe der Opiate und
sind im Drogenmilieu weit verbreitet, die Anwendung erfolgt
i.\,d.\,R. intravenös. 

\TextSlide{\TxSymptome{} der Opiatvergiftung}
{%
Die gefährlichste Wirkung der Opiate ist die
Atemdepressionen, d.\,h. die Atemfrequenz wird dosisabhängig
immer weiter reduziert bis es schließlich zu einem
Atemstillstand kommt. Die Patienten zeigen
dementsprechend Zeichen einer Hypoxie
(Zyanose). Charakteristisch ist die starke Verengung der
Pupillen \E{(stecknadelkopfgroße
Pupillen)}\index{Pupillen!stecknadelkopfgroße}. Oft
findet man alte Einstichstellen in der Armbeuge. Da in der
\enquote{Szene} oft unsauberes oder getauschtes Spritzbesteck
verwendet wird, besteht eine erhöhte Erkrankungsrate an HIV
und Hepatitis C. Eine weitere Gefahr geht vom Spritzbesteck
aus, wenn es nicht sicher verwahrt wird (\Ca{Vorsicht vor
der Nadel!})

\bigskip

\bigskip

\bigskip
}{
\begin{itemize}
    \item Atemdepression
    \item Hypoxie
    \item Stecknadelkopfgroße Pupillen
    \item Oft Begleiterkrankungen (HIV, Hep. C) 
    \item \Ca{Vorsicht Spritzbesteck!}
    \item Gegenmittel verfügbar, wirkt aber nicht so lang
\end{itemize}
}{}{}{}

\Text{Gegenmittel}
{%
Jedes Notarztmittel führt ein Opiat-Gegenmittel mit sich,
dieses ist sehr effektiv.\footnote{\Z{Naloxon. Im Handel
als: \emph{Narcanti}\texttrademark,
\emph{Naloxon-ratiopharm}\texttrademark, u.\,a.}} Es ist
jedoch wichtig zu wissen, dass dieses Gegenmittel nicht so
lange wirkt wie das Opiat, d.\,h. es kann wieder zu einer
verminderten Atmung kommen, wenn die Wirkung des
Gegenmittels nachlässt. Daher soll der Transport nur mit
Arztbegleitung erfolgen.

}{

}{}{}{}


\subsubsection{'Clubber' -- Ecstacy, Schwammerl,
Special K \ldots}

\TextSlide{\TxBeschreibung}
{Als \T{Partydrogen} bezeichnet man Amphetaminpräparate oder
ähnliche Stimulantien. Zu den bekanntesten Drogen dieser Art
gehören  \I{Ecstasy} (\I{XTC}, \I{E}, \I{X}\Z{; MDMA
(3,4-Methylen\-dioxy\-meth\-amphet\-amin)}). Verwendet wird
die Substanz häufig als \EE{Partydroge}: Der Patient erlebt
ein gesteigertes Glücksgefühl, ein Gefühl von
unerschöpflicher Energie und ein gesteigertes Erleben der
Umwelt.

\bigskip

\bigskip

\bigskip

}{
\begin{itemize}
  \item Ecstasy (XTC, E, X).
  \item Gesteigertes \enquote{Glücksgefühl}.
  \item Gefühl von unerschöpflicher Energie.
  \item Gesteigertes Erleben der Umwelt.
\end{itemize}
}{}{}{}

\TextSlide{\TxSymptome}
{%
Die Nebenwirkungen sind jedoch
gravierend und können lebensgefährlich sein. Man kann sie
unterteilen in körperliche und psychische Nebenwirkungen. 
Psychisch kommt es dosisabhängig zu Halluzinationen
und Panik-Attacken. 
Körperlich können Patienten tachykarde
Herzrhythmusstörungen zeigen, Krampfanfälle, hypertensive
Krisen, massiver Körpertemperaturanstieg über
40\textdegree{} und Muskelzerfall.
Bei einer ausgeprägten Vergiftung kann eine vitale
Bedrohung bestehen. 

\vspace{2cm}

}{
\begin{itemize}
  \item Nebenwirkungen Psychisch
  \begin{itemize}
    \item Halluzinationen
    \item Panik-Attacken
  \end{itemize}
  \item Nebenwirkungen Körperlich
  \begin{itemize}
    \item Tachykarde Herzrhythmusstörungen
    \item Hypertensive Krisen
    \item Krampfanfälle
    \item Massiver Temperaturanstieg über
    {\textgreater}\,40\DegC*
  \end{itemize}
\end{itemize}
}{}{}{}



%„Liquid Ecstasy”: Gamma-Butyrolacton-Entzugsdelir mit
%Rhabdomyolyse und dialysepflichtiger Niereninsuffizienz
\cite{Supady2009}








\clearpage % TEMP temp clearpage
\subsection{Gase}

\subsubsection{Vergiftungen mit Stickgasen}

\MassnahmenNG{Allgemeine Maßnahmen}{Stickgasvergiftungen}{\label{m:am-stickgase}}{


\begin{enumerate}
  \item Auf Selbstschutz achten!
  \item Patienten retten, aber nur wenn es der
  Eigenschutz zulässt.
%   \item \ce{O2}-Berieselung mit Maske ist in keinem Fall
%   ein ausreichender Atemschutz!
%   
  \item Spezialkräfte anfordern (Schwerer Atemschutz).
  I.\,d.\,R. Feuerwehr
  \item Weitere Opfer ausschließen (lassen):
  Nachbarwohnungen, etc.
%   \item Sicherung der Vitalfunktionen, Beurteilung ob der
%   Patient vital bedroht ist.
%   \item \TxBeiVit \TxMassVitBes
%   \begin{itemize}
%     \item \ce{O2} hochdosiert, 15\LPerMin*
%     \item Lagerung: situationsgerecht.
%   \end{itemize}
%   \item Traumacheck (Folgeverletzungen)
%   \item Neurocheck inkl. Blutzuckerbestimmung
%   (Differentialdiagnosen, Beurteilung der Beeinträchtigung im Verlauf)
\end{enumerate}
}


\subsubsection{Kohlenmonoxid (\ce{CO}) -- oder: Von defekten
Heizlüftern, Gasthermen und Autoabgasen}

\Layout{60}{Einleitung}{%
Kohlenmonoxid (\ce{CO}) ist ein farb- und geruchloses Gas.
Esentsteht bei unvollständiger Verbrennung (\ce{O2}-Zufuhr
${\downarrow}$), dies passiert vor allem bei schlecht
gewarteten Gasheizungen und Durchlauferhitzern. \ce{CO}
bindet 300 mal besser an den roten Blutfarbstoff Hämoglobin
als \ce{O2}, daher verdrängen schon kleine Mengen den
Sauerstoff aus dem Blut! Die Folge ist logischerweise eine
Unterversorgung mit Sauerstoff.
}
{}
{./images/gabriel-sebastian-aass/Atemschutz_IMG_2043_00800.jpg}
{}
{GABS}




\TextSlide{Anamnese und Szene}
{
\begin{itemize}
    \item Mehrere, unerklärliche Opfer, leblose Haustiere
    \item Gasofen, Durchlauferhitzer
    \item \emph{\enquote{Immer beim duschen!}}
\end{itemize}
}{
\begin{itemize}
    \item Mehrere, unerklärliche Opfer, tote Tiere
    \item Gasofen, Durchlauferhitzer
    \item \emph{\enquote{Immer beim duschen!}}
\end{itemize}
}{}{}{}

\TextSlide{\TxSymptome}
{% 
Kohlenmonoxid-beladenes Blut ist -- wie
Sauerstoff-beladenes -- hellrot. Der Patient ist daher eher
nicht bleich oder zyanotisch, sondern die Haut ist oft rosig.
Pulsoxymeter älterer Bauart verwechseln
Kohlenmonoxid-beladenes und Sauerstoff-beladenes Blut, daher
wird ein falsch hoher Sauerstoffsättigungswert angezeigt!

Es stehen die Symptome des \EE{Sauerstoffmangels} im
Vordergrund (Hypoxie): Der Patient leidet an Atemnot,
allerdings \E{ohne} Zyanose, begleitet von
Allgemeinsymptomen wie Schwindel, Kopfschmerz, Übelkeit und
Erbrechen, evtl. besteht auch eine Tachykardie.
Es können sich rasch Bewusstseinsstörungen bis hin zur
Bewusstlosigkeit entwickeln, auch Krämpfe sind gut möglich.
}{
\begin{itemize}
  \item Dyspnoe  \E{ohne}  Zyanose
  \item Tachykardie
  \item Schwindel, Kopfschmerz, Übelkeit, Erbrechen
  \item Bewusstseinsstörungen, Koma
  \item Krämpfe
  \item Pulsoxymeter zeigt falsche Werte an!
\end{itemize}
}{}{}{}

\Layout{27}
{}
{}
{}
{./images/2006-01-31-gaswienheute2.png}
{\EE{\ce{CO}-Vergiftung}:
Regel\-mäß\-ig kommt es zu Unfällen durch defekte Gasthermen
oder Heizlüfter. Bei diesem Einsatz war ein RTW des ASB
Floridsdorf-Donaustadt beteiligt. Die dritte Patientin
wurde erst \EE{nachdem die Feuerwehr \E{sämtliche}\,(!)
Wohnungen des Wohnhauses aufgebrochen hatte} in der
\E{darüberliegenden Wohnung} gefunden.}
{Wien heute, \url{http://wien.orf.at}}


\subsubsection{Kohlendioxid (\ce{CO2}) -- Tod im Weinkeller}


Kohlendioxid ist ein farb- und geruchsloses Gas. Es ist
schwerer als Luft und sammelt sich z.\,B. in Mulden,
Weinkellern, Silos, etc. Es fällt als Gär- oder
Verbrennungsgas gehäuft in der Landwirtschaft und der
Industrie an.


\TextSlide{\TxSymptome}
{%
\begin{itemize}
  \item Kopfschmerz, Schwindel, Übelkeit
  \item Krämpfe
  \item Zyanose
  \item Hyperventilation, da der hohe
  \ce{CO2}-Spiegel im Blut das Atemzentrum stimuliert
\end{itemize}
}{
\begin{itemize}
  \item Kopfschmerz, Schwindel, Übelkeit
  \item Krämpfe
  \item Zyanose
  \item Hyperventilation
\end{itemize}
}{}{}{}


\Cave{
Bei Gasunfällen hat der \EE{Selbstschutz} Vorrang.
Oft ist die Bergung nur mit \EE{schwerem
Atemschutz} möglich ${\rightarrow}$ FEUERWEHR!}

\Fazit{Ein toter Helfer ist keine Hilfe.}

\subsubsection{Reizgase}


\Text{Typen}
{%
\begin{description}
  \item[Soforttyp] (z.\,B. Tränengas,...) Augenrötung, starker Hustenreiz
  \item[Verzögerungstyp] \Z{(z.\,B. Nitrosegase)} kaum
  Hustenreiz, Latenzphase (kaum Symptome),
  \EE{\VonBis{3}{24}\Hour* später: toxisches Lungenödem};
  bei Nichtbehandlung evtl. irreversible Lungenschäden
\end{description}
}{

}{}{}{}

\Layout{46}{Vorkommen}
{}{}{%
\begin{tabulary}{\linewidth}{LL}\addlinespace\toprule\addlinespace
\noindent\TabHeading{Vorkommen} & \noindent\TabHeading{Art}
\\\addlinespace\midrule\addlinespace
Chemische Industrie 
&
Nitrosegase,  Schwefelwasserstoffe, Ammoniak
\\\addlinespace
Kunststoffindustrie 
&
Chlorwasserstoffe,  Formaldehyd
\\\addlinespace
Düngemittelherstellung 
&
Ammoniak,  Nitrosegase, Schwefelwasserstoffe
\\\addlinespace
Haushalt 
&
Chlorgas (Kombination von best.  Haushaltsreinigern mit
Essig), Lacke,  Imprägniersprays, ...
\\\addlinespace\bottomrule
\end{tabulary}
}{Vorkommen von Reizgasen.}{}





\MassnahmenNG{Allgemeine
Maßnahmen}{Reizgasvergiftungen}{\label{m:am-reizgase}}{


\begin{enumerate}
  \item Auf Selbstschutz achten!
  \begin{itemize}
    \item Patienten retten, aber nur wenn es der
    Eigenschutz zulässt.
%     \item \ce{O2}-Berieselung mit Maske ist in keinem Fall
%     ein ausreichender Atemschutz!
  \end{itemize}
  \item Gefahr erkennen: Welcher Stoff? Expertenrat einholen!
  \item Wenn erforderlich: Absperrmaßnahmen durchführen
  \item Wenn erforderlich:  Spezialkräfte anfordern
%   \item Vorgehen analog zu Abschnitt \enquote{Gefahrengutunfall}
%   (\prettyref{topic:gefahrengutunfall})
  \item Weitere Opfer ausschließen (lassen):
  Nachbarwohnungen, etc
%   \item Sicherung der Vitalfunktionen, Beurteilung ob der
%   Patient vital bedroht ist.
%   \item \TxBeiVit \TxMassVitBes
%   \begin{itemize}
%     \item \ce{O2} hochdosiert, 15\LPerMin*
%     \item Lagerung: situationsgerecht
%   \end{itemize}
%   \item Traumacheck (Folgeverletzungen)
%   \item Neurocheck inkl. Blutzuckerbestimmung
%   (Differentialdiagnosen, Beurteilung der Beeinträchtigung im Verlauf)
\end{enumerate}
}




\subsubsection{Kampfgase}

\Layout{61}{}
{%
Kurz und bündig: Hier gehören Spezialisten her.
Selbstschutz und Expertenrat einholen.
}{
}{./images/gabriel-sebastian-aass/schild-atemschutzmaske-abnehmen.jpg}{}{}


\clearpage
\subsection{Wovon man nicht trinken sollte \ldots}

\subsection{Einnahme von Säuren und Laugen}
\label{topic:veraetzung-intoxikation}



\Definition{\T{Verätzung}: Lokale Schädigung der
Oberfläche aufgrund einer irreversiblen Zerstörung (Denaturierung) von
Eiweißstoffen.}

% 
% Betroffene Organe
% 
% \begin{itemize}
%     \item Magen-Darm-Trakt
%     \item Haut
%     \item Augen
% \end{itemize}

\TextSlide{\TxSymptome}
{
\begin{itemize}
    \item Ätzspuren (Mund/Rachen), Abrinnspuren
    \item Würgen
    \item Erbrechen
    \item Schmerzen
\end{itemize}
}{%
\begin{itemize}
  \item Ätz-, Abrinnspuren
  \item Würgen
  \item Erbrechen
  \item Schmerzen
\end{itemize}
}{}{}{}

\TextSlide{\TxGefahren}
{%
\begin{itemize}
  \item Ruptur/Perforation Speiseröhre oder Magen
  \item Blutung
\end{itemize}
}{
\begin{itemize}
  \item Ruptur/Perforation Speiseröhre oder Magen
  \item Blutung
\end{itemize}
}{}{}{}


\SlideCompensationNoParskip{2}

\MassnahmenNG{Spezielle Maßnahmen}{Einnahme von Säuren oder Laugen}{}{
\begin{itemize}
%   \item Allgemeine Maßnahmen bei Vergiftungen
%   (\prettyref{m:am-vergiftung})
  \item Schluckweise Wasser zur Verdünnung trinken lassen
  \item \EE{Niemals absichtlich Erbrechen  herbeiführen!}
  \item \Z{Keine Neutralisationsversuche! (Milch
  obsolet!)}
\end{itemize}
}


% Verätzung der Haut \prettyref{topic:veraetzung}



\subsubsection{Schaumbildner (Wasch-/Putzmittel)}


\TextSlide{{\TxBeschreibung}, {\TxSymptome} und {\TxGefahren}}{%
Schaumbildner vermindern die Oberflächenspannung, dadurch kommt es
zur Bildung von Schaum. Werden Schaumbildner
eingenommen, besteht durch die Schaumbildung eine \EE{hohe
Aspirationsgefahr} bzw. eine hohe Gefahr der \EE{Verlegung}
der Atemwege.
Zusätzlich kann der Patient über Übelkeit klagen bzw. erbrechen.

\bigskip

\bigskip

\bigskip
}{
\begin{itemize}
  \item Schaumbildung
  \begin{itemize}
    \item Aspirationsgefahr
    \item Verlegung der Atemwege
  \end{itemize}
  \item Übelkeit, Erbrechen
\end{itemize}
}{}{}{}




\MassnahmenNG{Spezielle Maßnahmen}{Einnahme von
schaumbildenden Substanzen}{}{
\begin{itemize}
%     \item Allgemeine Maßnahmen bei Vergiftungen
%     (\prettyref{m:am-vergiftung})
    \item Patienten nüchtern lassen
    \item \EE{Niemals absichtlich Erbrechen 
    herbeiführen!}
\end{itemize}
}





% 
% 
% \subsection{\Z{Reserviert}\A{Vergiftungen mit
% Insektiziden}}
% % TODO Vergiftungen mit Insektiziden
% 
% \subsection{\Z{Reserviert}\A{Vergiftungen durch giftige
% Tiere}}
% % TODO Vergiftungen durch giftige Tiere
% 
% 
% 
}

\VARIANT{VAR-STANDARD}{\section*{Weiterführende Literatur}}


\cite{BuchfelderErsteHilfe4,Erc2005Sec9,Erc2005Sec2,Erc2005Sec8,Erc2005Sec6,WirthEH,HolaubekASB16h,KoehnleinEH,KeggenhoffEH,KarutzKursbuchEH}


% \ChapterBibliography
% 
% 
% {
% {\textsc{Böhmer, }}{Roman }{\textsc{/ Schneider,
% }}{Thomas }{\textsc{/ Wolcke }}{Benno: Reanimation
% kompakt}{\textsuperscript{2}}{. Nach den internationalen
% ERC-Richtlinien 2005 basierend auf den CoSTR 2005 (ILCOR/AHA) (Mainz:
% Naseweis Verlag, 2006)}}
% 
% 
% 
% 
% {
% {\textsc{Hackstock, }}{Horst }{\textsc{/
% Kühberger, }}{Rudolf: Arbeiter-Samariter-Bund Österreichs
% Bundes-schulungsreferat: Leitfaden für Vortragende. Lebensrettende
% Sofortmaßnahmen am Ort des Verkehrsunfalls (Wien:
% Arbeiter-Samariter-Bund Österreichs intern, Mai 2006)}}
% 
% 
% 
% 
% {
% {\textsc{Hackstock, }}{Horst }{\textsc{/
% Kühberger, }}{Rudolf: Arbeiter-Samariter-Bund Österreichs
% Bundes-schulungsreferat: Leitfaden für Vortragende.
% Breitenschulungskurs (Wien: Arbeiter-Samariter-Bund Österreichs
% intern, Sept. 2008)}}
% 
% 
% 
% 
% {
% {\textsc{Hochmeister}}{, Manfred /
% }{\textsc{Grassberger}}{, Martin /
% }{\textsc{Stimpfl}}{, Thomas: forensische Medizin. für
% Studium und Praxis (Wien: Maudrich Verlag, 2007)}}
% 
% 
% 
% 
% {
% {\textsc{Holaubeck}}{,
% Karl/}{\textsc{Nistler}}{,
% Kurt/}{\textsc{Scherer}}{, Nicolas: Erste Hilfe. 16
% Stunden für das Leben (Wien: Arbeiter-Samariter-Bund Österreichs,
% 2005)}}
% 
% 
% 
% 
% {
% {\textsc{Karutz,}}{ Harald / }{\textsc{von
% Buttlar}}{, Manfred: Kursbuch Erste
% Hilfe}{\textsuperscript{2}}{ (München: Deutscher
% Taschenbuch Verlag, 2008)}}
% 
% 
% 
% 
% {
% {\textsc{Keggenhoff, }}{Franz: Erste Hilfe. Das offizielle
% Handbuch. Sofortmaßnahmen bei Babys, Kindern und Erwachsenen
% (München: Südwest Verlag, 2007)}}
% 
% 
% 
% 
% {
% {\textsc{Köhnlein, }}{Edzard }{\textsc{/ Weller
% }}{Siegfried [Hrsg.]: Erste
% Hilfe}{\textsuperscript{10}}{ (Stuttgart: Georg Thieme
% Verlag, 2004)}}
% 
% 
% 
% 
% {
% {\textsc{Köhnlein, }}{Edzard }{\textsc{/ Weller
% }}{Siegfried [Hrsg.]: Erste Hilfe. Aktuelle Empfehlungen zur
% kardiopulmonalen Reanimation (Stuttgart: Georg Thieme Verlag, 2006)}}
% 
% 
% 
% 
% {
% {\textsc{Pschyrembel}}{, Willibald: Klinisches
% Wörterbuch (Berlin: Walter de Gruyter Verlag, 2004)}}
% 
% 
% 
% 
% {
% {\textsc{Wirth}}{, Armin: Erste Hilfe
% unterwegs}{\textsuperscript{3}}{. Effektiv und praxisnah
% (Bielefeld: Reise-Know-How Verlag Peter Rump, 2007)}}

\ChapterEnd
